% !TeX root = ../Book3.tex
\BookPart{基础复分析与正规族}{b3:part:01}
% 本编导读（章级标题，不计入编号）
\BookPartGuide

本编研究复导数所带来的解析性，以及积分公式怎样控制函数的局部行为与函数列的极限。复平面虽只有两个实维度，全纯性却使函数的内部值与边界值紧密相连。零点的孤立性、奇点附近的展开，以及函数族的紧致性，都将由这种联系逐步建立。

\ProseParagraphBreak

所需基础包括一元微积分、实平面的微分、级数收敛及紧致性。\chapref{b3:ch:01}从差商开始，建立全纯与解析的等价性，并给出初等函数和局部结构。\chapref{b3:ch:02}把局部积分论证推广到一般曲线和适当的定义域，建立相应的积分估计。\chapref{b3:ch:03}用 Laurent 展开描述奇点，并通过留数与辐角变化计数零点。\chapref{b3:ch:04}转向函数族，研究局部一致收敛与紧致性。各章之间的主要依赖关系如下：
\[
\begin{tikzcd}[column sep=2.8em,row sep=2em]
\text{第1章}\arrow[r]
 & \text{第2章}\arrow[r]\arrow[rd]
 & \text{第3章}\arrow[d]\\
 & & \text{第4章}
\end{tikzcd}
\]
第四章所需的两类结果分别是积分估计和零点计数：前者控制函数族的变化，后者解释极限下零点的稳定性。第一章为证明解析性而先建立圆盘内的积分公式，第二章再将其推广到更一般的定义域。

\ProseParagraphBreak

第一次阅读宜依章序推进，先掌握定义域、收敛方式与分支约定，再使用积分和级数公式。若目标是计算留数与实积分，应先完成前三章；若目标是理解共形映射的存在问题，还需要第四章对函数族极限的控制。选读内容可在主线熟悉后回看。

\ProseParagraphBreak

John B. Conway 的《Functions of One Complex Variable I》\cite{Conway1978Functions}与 Theodore W. Gamelin 的《Complex Analysis》\cite{Gamelin2001Complex}对幂级数、解析函数和复积分采用了不同的讲述次序。学习本编时，可以比较两书的定义和例子；对照证明时，还应注意各书此前已经建立了哪些积分结论。

% !TeX root = ../../Book3.tex
\chapter{全纯函数与幂级数}
\label{b3:ch:01}
\BookChapterNumber{b3:ch:01}
% 本章摘要（不计入编号）
\begin{chapterabstract}
复导数要求差商从所有方向趋近同一个值。这个局部条件最终会迫使函数具有收敛的幂级数展开，带来远强于实可微的刚性。本章从复差商与 Cauchy--Riemann 方程开始，证明幂级数的收敛与微分性质，再借助圆盘内的积分公式建立全纯与解析的等价性。随后构造指数、对数、三角函数和 Möbius 变换，最后由零点分解得到恒等定理、局部逆函数定理与开映射定理。
\end{chapterabstract}

% 本章目标（不计入编号）
\begin{learninggoals}
\begin{itemize}
\item\label{b3:item:01-goal-derivative}能区分复可微、实可微与全纯，正确使用 Cauchy--Riemann 方程，并解释非零复导数的几何意义。
\item\label{b3:item:01-goal-series}能计算收敛半径与闭子圆盘上的截断误差，判断逐项微分的范围，并说明局部积分公式如何推出幂级数展开。
\item\label{b3:item:01-goal-elementary}能在明确的定义域和分支约定下使用初等复函数，核对 Möbius 变换在有限点与无穷远点的行为。
\item\label{b3:item:01-goal-local}能利用零点阶数和连通性判断函数的唯一性，构造局部标准坐标并计算附近原像的数目，区分局部可逆、全局单射与开映射。
\end{itemize}
\end{learninggoals}

实函数可以无限次可微，却不由一点的全部导数确定。复函数会发生什么变化？先比较 \(f(z)=z^2\) 与 \(g(z)=\overline z\)：前者的差商沿所有方向有同一极限，后者沿实方向和虚方向给出相反的值。要求所有方向相容，看似只增加一个微分条件，却会迫使函数具有收敛幂级数，并使它在一小段曲线上的值决定整个连通区域中的值。这条推理需要两座桥：先用积分控制微分所不能直接控制的余项，再用幂级数的首个非零项描述局部行为。实线性映射的分解、几何级数的尾项以及局部开根的计算，将分别把这些步骤具体化。

\ProseParagraphBreak

比较实微分与复微分，有助于理解这种刚性：一般实线性映射可以分别改变两个方向，复线性映射却只能同时旋转和伸缩。局部求逆时，第一卷的实逆函数定理也会在这个二维图景中发挥作用。

\ProseParagraphBreak

\subsecref{b3:subsec:010204}先证明三角形上的 Goursat 定理与圆盘内的 Cauchy 积分公式，再由此推出全纯函数的局部幂级数展开。\chapref{b3:ch:02}将讨论一般曲线与定义域的拓扑条件。熟悉这些结果的读者仍需留意本卷关于对数主值和保向共形的约定。

\ProseParagraphBreak

John B. Conway 的《Functions of One Complex Variable I》\cite[Chapter III, \S1, pp. 30--31]{Conway1978Functions}从绝对收敛和几何级数出发讨论收敛半径，可与\secref{b3:sec:0102}的证明比较。Theodore W. Gamelin 的《Complex Analysis》\cite{Gamelin2001Complex}则依次讲述初等函数、解析函数、复积分和幂级数，与本章先建立幂级数性质、再构造初等函数的安排不同。

% 本章来源实读范围见参考文献原文件/来源记录.md；证明义务以现行正文和工作记录/审核进度.md为准。
% Conway 公开样章实读印刷30--31页；Gamelin仅核验官方书目与目录。
% 跨卷调用稳定标签与当前编号登记于跨卷引用表.md。

% !TeX root = ../../Book3.tex
\section{复可微与全纯}
\label{b3:sec:0101}

% 本节从差商与实全微分独立证明基本法则，不调用跨卷编号。
% 反例与旋转伸缩图为按本节辨析目标自行组织的表达。
% 未使用全纯函数导数连续、局部幂级数、开映射或局部逆函数结论。

把复平面识别为实平面之后，复函数可以看成两个实变量的向量值函数。复导数的存在比两个偏导数的存在要求更强：无论增量沿哪个方向趋于零，差商都必须趋于同一个复数。本节用实微分刻画这一条件，并说明非零复导数为何对应于局部的旋转与伸缩。

\subsection{复导数}
\label{b3:subsec:010101}

以下写 \(z=x+\mathrm{i}y\)，其中 \(\mathrm{i}^2=-1\)，并用 \(\overline z=x-\mathrm{i}y\) 和 \(\lvert z\rvert=(x^2+y^2)^{1/2}\) 分别表示共轭与模。所有邻域和极限均取复平面的欧氏距离。设 \(U\subseteq\mathbb C\) 为开集；这一条件保证每个 \(a\in U\) 附近的任意充分小增量都能参与差商。
\symidx{\(\overline z\)：复数 \(z\) 的共轭}
\symidx{\(\lvert z\rvert\)：复数 \(z\) 的模}

\begin{definition}\label{b3:def:01-complex-derivative}
设 \(f:U\rightarrow\mathbb C\)，且 \(a\in U\)。若有限极限
\[
 f'(a)\coloneqq\lim_{\substack{h\to0\\h\ne0}}
 \frac{f(a+h)-f(a)}{h}
\]
存在，则称 \(f\) 在 \(a\) 处\term[fu-kewei]{复可微}[complex differentiable]，称 \(f'(a)\) 为该点的\term[fu-daoshu]{复导数}[complex derivative]。
\end{definition}
\symidx{\(f'(a)\)：函数在 \(a\) 处的复导数}

极限存在等价于存在复数 \(c\)，使
\begin{equation}\label{b3:eq:01-complex-linear-approximation}
 f(a+h)=f(a)+ch+o(\lvert h\rvert).
\end{equation}
这里 \(o(\lvert h\rvert)\) 表示一个复余项，其模除以 \(\lvert h\rvert\) 趋于零；下文的 \(O(|h|^k)\) 或 \(O(h^k)\) 表示模不超过某个常数倍 \(|h|^k\) 的余项。上述等价性来自 \(\lvert rh^{-1}\rvert=\lvert r\rvert/\lvert h\rvert\)。系数只能是 \(c=f'(a)\)。这一表达立即说明复可微蕴含连续，也说明微分的首项是乘以一个固定复数，而不能是任意实线性变换。

\ProseParagraphBreak

例如 \(f(z)=z^2\) 满足 \(f(a+h)-f(a)=2ah+h^2\)，故处处复可微，且 \(f'(a)=2a\)。共轭函数则完全不同：当 \(h\) 为非零实数时，\(\overline h/h=1\)；当 \(h\) 为非零纯虚数时，\(\overline h/h=-1\)。因此 \(z\mapsto\overline z\) 在任何一点都不复可微，尽管它作为实平面上的映射处处具有连续的各阶偏导数。

\begin{proposition}\label{b3:prop:01-complex-derivative-rules}
设 \(f,g\) 在 \(a\) 复可微，\(\alpha,\beta\in\mathbb C\)。则线性组合与乘积在 \(a\) 复可微，且
\[
 (\alpha f+\beta g)'(a)=\alpha\cdot f'(a)+\beta\cdot g'(a),
 \quad
 (fg)'(a)=f'(a)\cdot g(a)+f(a)\cdot g'(a).
\]
若 \(g(a)\ne0\)，则商在 \(a\) 的某邻域有定义，并且
\[
 \left(\frac fg\right)'(a)
 =\frac{f'(a)\cdot g(a)-f(a)\cdot g'(a)}{g(a)^2}.
\]
\end{proposition}
\begin{proof}
线性组合的差商直接按极限的线性性处理。乘积的增量写成
\[
 f(a+h)\cdot\bigl(g(a+h)-g(a)\bigr)
 +g(a)\cdot\bigl(f(a+h)-f(a)\bigr),
\]
除以 \(h\) 后，利用 \(f\) 在 \(a\) 连续即可取极限。若 \(g(a)\ne0\)，连续性保证充分小时 \(\lvert g(a+h)\rvert\geq\lvert g(a)\rvert/2\)。于是
\[
 \frac{g(a+h)^{-1}-g(a)^{-1}}h
 =-\frac{g(a+h)-g(a)}{h\cdot g(a+h)\cdot g(a)}
 \longrightarrow-\frac{g'(a)}{g(a)^2}.
\]
再应用乘积公式得到商法则。
\end{proof}

复合运算的求导由\term[lianshi-faze]{链式法则}[chain rule]控制。直接把两个差商相乘时，可能遇到内层函数在邻近点取相同值的情形；用线性近似可以同时处理这些点。

\begin{proposition}[链式法则]\label{b3:prop:01-complex-chain-rule}
若 \(f\) 在 \(a\) 复可微，\(F\) 定义在 \(f(a)\) 的开邻域内，并在 \(f(a)\) 复可微，则复合函数在 \(a\) 的某邻域有定义，且
\[
 (F\circ f)'(a)=F'\bigl(f(a)\bigr)\cdot f'(a).
\]
\end{proposition}
\begin{proof}
令 \(b=f(a)\)，并写
\[
 \begin{aligned}
 f(a+h)-b&=f'(a)h+r(h),&r(h)&=o(\lvert h\rvert),\\
 F(b+k)-F(b)&=F'(b)k+s(k),&s(k)&=o(\lvert k\rvert).
 \end{aligned}
\]
把最后一式在 \(k=0\) 时的余项取为零。第一个展开给出 \(\lvert k\rvert=O(\lvert h\rvert)\)，因此代入 \(k=f(a+h)-b\) 后有 \(s(k)=o(\lvert h\rvert)\)，即使某些非零 \(h\) 对应 \(k=0\) 也成立。复合函数的线性首项遂为 \(F'(b)\cdot f'(a)h\)。其邻域定义由 \(f\) 在 \(a\) 的连续性保证。
\end{proof}

\subsection{Cauchy--Riemann 方程}
\label{b3:subsec:010102}

写 \(f=u+\mathrm{i}v\)，把 \(u,v\) 看作两个实变量的实值函数。作为实映射，\(f\) 在 \(a\) 可微是指存在实线性映射 \(A:\mathbb R^2\rightarrow\mathbb R^2\)，使
\[
 f(a+h)-f(a)=Ah+o(\lvert h\rvert).
\]
该映射称为\term[frechet-weifen]{Fréchet 微分}[Fréchet differential]，记作 \(Df(a)\)。若存在，它的矩阵由偏导数确定为
\[
 Df(a)=
 \begin{pmatrix}
 u_x(a)&u_y(a)\\
 v_x(a)&v_y(a)
 \end{pmatrix}.
\]
下文把这一条件简称为“实可微”。复数乘法还要求此矩阵具有特殊形状。
\symidx{\(Df(a)\)：复函数作为实映射的 Fréchet 微分}

\begin{theorem}[Cauchy--Riemann]\label{b3:thm:01-cauchy-riemann-criterion}
函数 \(f=u+\mathrm{i}v\) 在 \(a\) 复可微，当且仅当它在该点实可微，并满足\term[cauchy-riemann-fangcheng]{Cauchy--Riemann 方程}[Cauchy--Riemann equations]
\begin{equation}\label{b3:eq:01-cauchy-riemann}
 u_x(a)=v_y(a),\quad u_y(a)=-v_x(a).
\end{equation}
成立时，
\[
 f'(a)=u_x(a)+\mathrm{i}v_x(a)=v_y(a)-\mathrm{i}u_y(a).
\]
\end{theorem}
\begin{proof}
若 \(f'(a)=\alpha+\mathrm{i}\beta\)，由\eqref{b3:eq:01-complex-linear-approximation}得到实微分
\[
 Df(a)=
 \begin{pmatrix}\alpha&-\beta\\ \beta&\alpha\end{pmatrix}.
\]
比较矩阵元素即得方程及导数公式。反过来，若实微分存在且满足\eqref{b3:eq:01-cauchy-riemann}，令 \(c=u_x(a)+\mathrm{i}v_x(a)\)，则这个实微分恰好是乘法 \(h\mapsto ch\)。实可微的余项因此给出\eqref{b3:eq:01-complex-linear-approximation}，从而复导数存在。
\end{proof}

定理中的实可微不能删去。定义
\[
 f(x+\mathrm{i}y)=
 \begin{cases}
 \dfrac{x^3+\mathrm{i}y^3}{x^2+y^2},&(x,y)\ne(0,0),\\[4pt]
 0,&(x,y)=(0,0).
 \end{cases}
\]
由 \(x^6+y^6\leq(x^2+y^2)^3\) 得 \(\lvert f(z)\rvert\leq\lvert z\rvert\)，所以它在原点连续。沿两坐标轴计算，得到 \(u_x(0)=v_y(0)=1\)、\(u_y(0)=v_x(0)=0\)，满足 Cauchy--Riemann 方程。但当 \(t\ne0\) 时，
\[
 \frac{f(t)}t=1,\quad
 \frac{f(t+\mathrm{i}t)}{t+\mathrm{i}t}=\frac12.
\]
差商没有极限。这个例子同时具有连续性与满足方程的四个偏导数，仍缺少对任意小增量的统一线性近似。

\begin{corollary}\label{b3:cor:01-continuous-partial-criterion}
若 \(u,v\) 的四个一阶偏导数在 \(a\) 的某邻域存在，并在 \(a\) 连续，则 \(f\) 在 \(a\) 复可微当且仅当 Cauchy--Riemann 方程在 \(a\) 成立。
\end{corollary}
\begin{proof}
只须证明实可微性。写 \(a=x+\mathrm{i}y\)。对 \(w=u\) 或 \(w=v\)，把增量从 \((x,y)\) 到 \((x+s,y+t)\) 分成先水平后竖直的两段。一元中值定理给出
\[
 w(x+s,y+t)-w(x,y)
 =s\cdot w_x(\xi,y)+t\cdot w_y(x+s,\eta),
\]
其中 \(\xi\) 介于 \(x\) 与 \(x+s\) 之间，\(\eta\) 介于 \(y\) 与 \(y+t\) 之间；零长度的一段直接记为零。减去 \(s\cdot w_x(a)+t\cdot w_y(a)\)，由两偏导在 \(a\) 连续，余项为 \(o\bigl((s^2+t^2)^{1/2}\bigr)\)。对两个分量合并后应用\cref{b3:thm:01-cauchy-riemann-criterion}。
\end{proof}

这里的连续偏导条件是检验复可微的一条充分途径。直接从差商定义出发时，不需要先假定导数随基点连续。

\subsection{全纯函数}
\label{b3:subsec:010103}

\begin{definition}\label{b3:def:01-holomorphic}
设 \(U\subseteq\mathbb C\) 为开集。若 \(f:U\rightarrow\mathbb C\) 在 \(U\) 的每一点复可微，则称它在 \(U\) 上\term[quan-chun]{全纯}[holomorphic]，称其为\term[quan-chun-hanshu]{全纯函数}[holomorphic function]，全体这样的函数记作 \(\mathcal H(U)\)。若 \(f\) 在某个包含 \(a\) 的开邻域上全纯，则称它在 \(a\) 附近全纯。在整个复平面上全纯的函数称为\term[zheng-hanshu]{整函数}[entire function]。
\end{definition}
\symidx{\(\mathcal H(U)\)：开集 \(U\) 上的全纯函数类}

全纯性对一个开集内的所有点提出要求，强于在一个指定点复可微。例如 \(f(z)=\lvert z\rvert^2\) 在原点的差商为 \(\overline h\)，所以 \(f'(0)=0\)。若 \(a=x+\mathrm{i}y\ne0\)，沿实轴与虚轴增量所得的差商极限分别为 \(2x\) 与 \(-2\mathrm{i}y\)，两者相等只可能发生在 \(a=0\)。故该函数在原点复可微，却不在原点的任何邻域上全纯。

\ProseParagraphBreak

由\cref{b3:prop:01-complex-derivative-rules}，多项式都是整函数，有理函数在分母不为零的开集上全纯；全纯函数的线性组合与乘积仍全纯。\cref{b3:prop:01-complex-chain-rule}还保证定义范围内的复合全纯。商法则只允许在分母不为零处使用。至于某个零点是否能通过重新赋值补上，需要另作极限分析，不能单凭代数表达式下结论。

\ProseParagraphBreak

在本卷中，非空连通开集称为\term[quyu]{区域}[domain]。连通条件尚未出现在全纯函数的定义里，但它让局部结论能够传播。例如，区域上的实值全纯函数只能是常数。

\begin{proposition}\label{b3:prop:01-real-valued-holomorphic}
若 \(U\) 为区域，\(f\in\mathcal H(U)\)，且 \(f(U)\subseteq\mathbb R\)，则 \(f\) 为常数。
\end{proposition}
\begin{proof}
写 \(f=u+\mathrm{i}v\)，则 \(v=0\)。由 Cauchy--Riemann 方程，\(u_x=u_y=0\)。对 \(U\) 中任意足够小的开矩形，沿其水平线段及竖直线段分别使用一元中值定理，得到 \(u\) 在该矩形上恒定。因此 \(u\) 局部恒定。固定 \(a\in U\)，集合 \(E=\{\,z\in U\mid u(z)=u(a)\,\}\) 非空，局部恒定性使 \(E\) 与其在 \(U\) 内的补集都为开集。由 \(U\) 连通，\(E=U\)。
\end{proof}

全纯函数的定义只含一次复导数。\secref{b3:sec:0102}将证明它还能在每一点附近展开成幂级数，因而具有任意阶导数。这个结论是需要证明的正则性，不能在当前的判别中倒过来充当假设。

\subsection{共形性}
\label{b3:subsec:010104}

若实线性映射 \(A:\mathbb C\rightarrow\mathbb C\) 还满足 \(A(\mathrm{i}h)=\mathrm{i}\cdot A(h)\)，则它与复数数乘相容，称为\term[fu-xianxing-yingshe]{复线性映射}[complex-linear map]。令 \(c=A(1)\)，利用实线性可得
\[
 A(x+\mathrm{i}y)=x\cdot A(1)+y\cdot A(\mathrm{i})=(x+\mathrm{i}y)c.
\]
因此复平面上的复线性映射恰好是乘以一个复数。若 \(c\ne0\)，写 \(c=\rho(\cos\theta+\mathrm{i}\sin\theta)\)，其中 \(\rho=\lvert c\rvert>0\)，则
\[
 A=\rho
 \begin{pmatrix}
 \cos\theta&-\sin\theta\\
 \sin\theta&\cos\theta
 \end{pmatrix}.
\]
这表示先旋转角 \(\theta\)，再把长度都乘以 \(\rho\)。\cref{b3:fig:01-complex-similarity}给出一组向量的变化；两个方向的夹角保持不变。

\begin{figure}[htbp]
\centering
\begin{tikzpicture}[>=stealth,scale=0.95]
 \begin{scope}
  \draw[->,AnalysisGray] (-0.45,0)--(2.05,0);
  \draw[->,AnalysisGray] (0,-0.4)--(0,1.85);
  \fill (0,0) circle (1.4pt);
  \node[below left] at (0,0) {\(0\)};
  \draw[->,thick,AnalysisBlue] (0,0)--(1.45,0) node[below] {\(h_1\)};
  \draw[->,thick,AnalysisBlue] (0,0)--(0.725,1.256) node[above right] {\(h_2\)};
  \draw[AnalysisBlue] (0.45,0) arc[start angle=0,end angle=60,radius=0.45];
  \node at (0.71,0.38) {\(\alpha\)};
 \end{scope}
 \draw[->,thick] (2.25,0.8)--(4.10,0.8)
  node[midway,above] {\(h\mapsto ch\)};
 \begin{scope}[xshift=5.2cm]
  \draw[->,AnalysisGray] (-0.55,0)--(2.0,0);
  \draw[->,AnalysisGray] (0,-0.4)--(0,2.1);
  \fill (0,0) circle (1.4pt);
  \node[below left] at (0,0) {\(0\)};
  \draw[->,thick,AnalysisBlue] (0,0)--(1.575,0.909) node[right] {\(ch_1\)};
  \draw[->,thick,AnalysisBlue] (0,0)--(0,1.819) node[above right] {\(ch_2\)};
  \draw[AnalysisBlue] (0.433,0.25) arc[start angle=30,end angle=90,radius=0.5];
  \node at (0.41,0.72) {\(\alpha\)};
  \draw[AnalysisGray] (0.85,0) arc[start angle=0,end angle=30,radius=0.85];
  \node at (1.12,0.25) {\(\theta\)};
 \end{scope}
\end{tikzpicture}
\caption[复线性微分的旋转与伸缩]{复线性微分的旋转与伸缩。图中取 \(\theta=\pi/6\)、\(\rho>1\)，相交方向的有向夹角 \(\alpha\) 不变。}
\label{b3:fig:01-complex-similarity}
\end{figure}

\begin{definition}\label{b3:def:01-conformal}
设 \(f\) 在 \(a\) 实可微。若 \(Df(a)\) 是非零的旋转伸缩，即具有上述 \(\rho>0\) 的形式，则称 \(f\) 在 \(a\) \term[gongxing]{共形}[conformal]。本卷的这一用语约定保持定向；在开集的每一点共形时，称该映射在此开集上共形。
\end{definition}

对实线性可逆映射，这个定义也等价于保持有向夹角。事实上，标准基 \(1,\mathrm{i}\) 的像必须垂直且定向为正，可写成 \(r e,s\mathrm{i}e\)，其中 \(r,s>0\)、\(\lvert e\rvert=1\)。再考察 \(1\) 与 \(1+\mathrm{i}\) 的夹角 \(\pi/4\)，得到 \(s/r=1\)，故两个伸缩因子相同。反向的角度保持直接由旋转伸缩形式得到。

\begin{proposition}\label{b3:prop:01-conformal-criterion}
函数 \(f\) 在 \(a\) 共形，当且仅当它在 \(a\) 复可微且 \(f'(a)\ne0\)。此时
\[
 \det Df(a)=\lvert f'(a)\rvert^2>0.
\]
若两条实参数曲线在 \(a\) 相交，且在交点具有非零切向量，则它们的像在该点仍具有非零切向量，二者的有向夹角保持不变。
\end{proposition}
\begin{proof}
由\cref{b3:thm:01-cauchy-riemann-criterion}，复可微恰好表示实微分是复数乘法；乘数不为零时，它恰好是非零旋转伸缩。矩阵形式给出行列式等式。

设两条曲线满足 \(\gamma_j(0)=a\)、\(\gamma_j'(0)=v_j\ne0\)。把 \(\gamma_j(t)=a+tv_j+o(\lvert t\rvert)\) 代入\eqref{b3:eq:01-complex-linear-approximation}，得到
\[
 (f\circ\gamma_j)'(0)=f'(a)v_j\ne0.
\]
两个像向量相除所得的比值为 \(v_2/v_1\)，与原向量的比值相同，因而有向夹角相同。
\end{proof}

共轭映射虽保持夹角的大小，却把有向夹角反号，不满足这里的保向约定。导数为零时也不能沿用上述结论：对 \(f(z)=z^2\)，从原点出发的射线辐角会变为原来的两倍，而在原点可微的实参数曲线的像切向量都被微分压成零向量。最后，共形条件本身是局部的微分条件，并未断言整个定义域上的单射性；例如 \(z^2\) 在 \(\mathbb C\setminus\{0\}\) 处处共形，却把 \(z\) 与 \(-z\) 映到同一点。单射性以及局部逆函数的存在将在\subsecref{b3:subsec:010403}中分别讨论。

% !TeX root = ../../Book3.tex
\section{幂级数}
\label{b3:sec:0102}

多项式可以直接求导，而无穷级数的求导必须受到收敛条件约束。幂级数的特殊之处在于：只要把讨论限制在收敛圆的内部，级数本身及它的各阶导数就都有统一控制。更深的一步是反过来证明，每个全纯函数局部都来自这样的级数。这个结论不能仅由 Cauchy--Riemann 方程的形式推出。

\subsection{收敛半径}
\label{b3:subsec:010201}

固定 \(a\in\mathbb C\) 及复数列 \((c_n)_{n\geq0}\)。形如
\[
\sum_{n=0}^\infty c_n(z-a)^n
\]
的级数称为以 \(a\) 为中心的\term[miji-shu]{幂级数}[power series]。常数项为 \(c_0\)，在 \(z=a\) 处级数总是收敛。记 \(D(a,r)=\{z\in\mathbb C:|z-a|<r\}\)，并约定 \(D(a,\infty)=\mathbb C\)。
\symidx{\(D(a,r)\)：中心为 \(a\)、半径为 \(r\) 的开圆盘}

\ProseParagraphBreak

先观察收敛向中心传播的原因。若级数在某个 \(z_1\ne a\) 处收敛，置 \(\rho=|z_1-a|\)，则它的通项有界，存在 \(M_\rho<\infty\) 使
\[
 |c_n|\rho^n\le M_\rho\quad(n\ge0).
\]
对任何 \(0<r<\rho\)，在 \(|z-a|\le r\) 上都有
\[
 |c_n(z-a)^n|\le M_\rho(r/\rho)^n.
\]
右端是可求和的几何数列。因此，一个非中心点的收敛就保证更小闭圆盘上的绝对一致收敛；这个论证没有假定原级数在 \(z_1\) 绝对收敛。控制范围天然呈圆形，是因为各项的绝对值只含距离 \(|z-a|\)。

\ProseParagraphBreak

下面的\term[cauchy-hadamard-gongshi]{Cauchy--Hadamard 公式}[Cauchy--Hadamard formula]从系数确定收敛的范围。

\Needspace{6\baselineskip}
\begin{theorem}[Cauchy--Hadamard]\label{b3:thm:01-power-radius}
设 \(a\in\mathbb C\)、\(c_n\in\mathbb C\ (n\ge0)\)，考虑幂级数 \(\sum_{n=0}^{\infty}c_n(z-a)^n\)。令 \(q=\limsup_{n\to\infty}|c_n|^{1/n}\)，并规定 \(1/0=\infty\)、\(1/\infty=0\)。则 \(R=1/q\) 满足：级数在 \(|z-a|<R\) 时绝对收敛，在 \(|z-a|>R\) 时发散。对任意 \(0<r<R\)，级数在闭圆盘 \(\overline{D(a,r)}\) 上绝对一致收敛。
\end{theorem}
\begin{proof}
若 \(0<r<R\)，可选 \(b>q\) 使 \(br<1\)。由上极限的定义，对充分大的 \(n\) 有 \(|c_n|\leq b^n\)，故在 \(|z-a|\leq r\) 上，级数尾部由几何级数 \(\sum_n(br)^n\) 控制。它既证明绝对收敛，也给出与 \(z\) 无关的尾部误差估计。

若 \(s=|z-a|>R\)，则可选 \(b<q\) 使 \(bs>1\)，并有无穷多个 \(n\) 满足 \(|c_n|^{1/n}>b\)。对这些指标，\(|c_n(z-a)^n|>(bs)^n\)，所以通项不趋于零。上述选择在 \(q=0\) 或 \(q=\infty\) 时仍按相应方向适用：前者任意有限闭圆盘上收敛，后者每个非中心点处发散。
\end{proof}

对复系数幂级数 \(\sum_{n\ge0}c_n(z-a)^n\)，数 \(R=(\limsup_{n\to\infty}|c_n|^{1/n})^{-1}\)（倒数按扩展非负实数约定）称为\term[shoulian-banjing]{收敛半径}[radius of convergence]。定理没有判定边界。半径同为 \(1\) 的三个级数
\[
\sum_{n=0}^\infty z^n,\quad
\sum_{n=1}^\infty\frac{z^n}{n},\quad
\sum_{n=1}^\infty\frac{z^n}{n^2}
\]
分别展示不同情形：第一个在单位圆上处处发散；第二个在 \(z=1\) 处发散、在 \(z=-1\) 处由交错级数判据收敛；第三个在整个闭单位圆盘上绝对一致收敛，因为尾部由 \(\sum n^{-2}\) 控制。因此，不能把开圆盘内的结论直接延到边界。

\ProseParagraphBreak

两个极端值也会实际出现。若 \(c_n=n!\)，则对 \(n\ge2\)，最后至少 \(n/2\) 个因子不小于 \(n/2\)，故
\[
 (n!)^{1/n}\ge(n/2)^{1/2}\longrightarrow\infty.
\]
相应级数的半径为零，只有中心点收敛。若 \(c_n=n^{-n}\)（\(n\ge1\)），则 \(|c_n|^{1/n}=1/n\to0\)，半径为无穷大。这些约定并非只为统一记号；没有正收敛半径的形式级数，不能直接当作邻域中的函数使用。

\ProseParagraphBreak

闭子圆盘上的控制还能给出具体误差。若 \(0<r<\rho<R\)，并取 \(M_\rho=\sup_{n\ge0}|c_n|\rho^n<\infty\)，则对每个整数 \(N\ge0\)，有
\begin{equation}\label{b3:eq:01-power-tail}
 \sup_{|z-a|\le r}\left|f(z)-\sum_{n=0}^{N}c_n(z-a)^n\right|
 \le M_\rho\frac{(r/\rho)^{N+1}}{1-r/\rho}.
\end{equation}
这是把从 \(N+1\) 开始的控制几何级数相加所得的，其中 \(f\) 表示原级数的和。先固定比值 \(r/\rho<1\)，再令 \(N\to\infty\)，误差便统一趋零；若让 \(r\uparrow\rho\)，分母变小，这项估计就不再提供统一的边界控制。

\ProseParagraphBreak

收敛半径还依赖展开中心。例如对 \(a\neq1\)，几何级数给出
\[
\frac1{1-z}=\frac1{1-a}\sum_{n=0}^\infty
\left(\frac{z-a}{1-a}\right)^n,
\quad |z-a|<|1-a|.
\]
由系数直接计算，半径恰为 \(|1-a|\)。同一个函数在不同中心的展开，能够覆盖的圆盘大小并不相同。

\subsection{逐项求导}
\label{b3:subsec:010202}

一致收敛只控制函数值，一般并不自动保留导数。例如实区间上的 \(g_N(t)=\sin(Nt)/N\) 一致趋于零，但 \(g_N'(0)=1\)，并不趋于极限函数在零点的导数。幂级数允许逐项求导，需要利用它在较小圆盘上的额外结构：微分所增加的 \(n,n(n-1),\ldots\) 能被几何衰减吸收。证明时先控制形式导数级数，再确认其和确实是原函数的导数。

\Needspace{6\baselineskip}
\begin{theorem}\label{b3:thm:01-power-differentiation}
设 \(a\in\mathbb C\)、\(c_n\in\mathbb C\ (n\ge0)\)，幂级数 \(f(z)=\sum_{n=0}^\infty c_n(z-a)^n\) 的收敛半径为 \(R\in(0,\infty]\)。则 \(f\) 在 \(D(a,R)\) 上全纯，并可任意次\term[zhuxiang-qiudao]{逐项求导}[termwise differentiation]：
\begin{equation}\label{b3:eq:01-power-derivatives}
f^{(k)}(z)=\sum_{n=k}^\infty\frac{n!}{(n-k)!}c_n(z-a)^{n-k}
\quad(k\in\mathbb Z,\ k\geq1).
\end{equation}
各导数级数在任意闭子圆盘上一致收敛，收敛半径仍为 \(R\)。特别地，\(c_k=f^{(k)}(a)/k!\)。
\end{theorem}
\begin{proof}
固定 \(0<r<\rho<R\)。由于级数在 \(a+\rho\) 收敛，通项有界，故存在 \(M\) 使 \(|c_n|\rho^n\leq M\)。于是对 \(|z-a|\leq r\)，有
\[
\left|\frac{n!}{(n-k)!}c_n(z-a)^{n-k}\right|
\leq M\rho^{-k}n^k\left(\frac r\rho\right)^{n-k}.
\]
右侧关于 \(n\geq k\) 可求和，这可由相邻项比值趋于 \(r/\rho<1\) 验证。因而所有形式导数级数在闭子盘上都一致收敛，其和连续。

先处理一阶。记部分和为 \(S_N\)，形式一阶导数级数的和为 \(g\)。当线段 \([z,z+h]\) 留在一个闭子盘中时，对多项式沿该线段积分得到
\[
S_N(z+h)-S_N(z)=h\int_0^1 S_N'(z+th)\,\mathrm{d}t.
\]
两组级数的一致收敛允许令 \(N\to\infty\)，故对 \(h\neq0\)，
\[
\frac{f(z+h)-f(z)}h=\int_0^1g(z+th)\,\mathrm{d}t.
\]
由 \(g\) 连续，右端在 \(h\to0\) 时趋于 \(g(z)\)。所以 \(f'=g\)。对导数级数重复同一论证，即得各阶公式。

估计已经说明导数级数的半径不小于 \(R\)。设一阶导数级数的半径为 \(R_1\)。对任意 \(0<s<R_1\)，该级数在距离中心为 \(s\) 的点绝对收敛，故
\[
\sum_{n=1}^\infty|c_n|s^n
\leq s\sum_{n=1}^\infty n|c_n|s^{n-1}<\infty.
\]
这说明每个 \(s<R_1\) 都不超过原级数的收敛半径，因而 \(R_1\le R\)。逐阶应用得到相同半径；当 \(R=\infty\) 时，前面的闭子盘估计已覆盖任意有限半径。最后代入 \(z=a\)，高于 \(k\) 次的项均为零，得到系数公式。
\end{proof}

公式\eqref{b3:eq:01-power-derivatives}也证明同一中心的幂级数表示唯一。若两个级数在中心的某个邻域内表示同一函数，对等式逐次求导并在中心求值，就得到每个系数相同。

\ProseParagraphBreak

例如在 \(|z|<1\) 内，对几何级数求导可得
\[
 \frac1{(1-z)^2}=\sum_{n=1}^{\infty}nz^{n-1}.
\]
两个级数的半径都为一，但这并不主张它们在边界有相同的收敛情况。相反，
\[
 L(z)=\sum_{n=1}^{\infty}\frac{z^n}{n^2}
\]
在闭单位圆盘上绝对一致收敛，而其形式导数在 \(z=1\) 变成调和级数。逐项求导定理适用于开圆盘中的每个点，并不把边界上的连续性提升为边界可微性。

\ProseParagraphBreak

还可以先逐项积分来构造一个导数已知的函数。若 \(f(z)=\sum_{n\ge0}c_n(z-a)^n\) 的半径为 \(R>0\)，则
\[
 F(z)=\sum_{n=0}^{\infty}\frac{c_n}{n+1}(z-a)^{n+1}
\]
在 \(D(a,R)\) 中收敛；对闭子盘，原级数的绝对控制再乘以有界距离即可控制新级数。刚证的求导定理给出 \(F'=f\)、\(F(a)=0\)。一般地，设 \(U\subseteq\mathbb C\) 为开集、\(h:U\to\mathbb C\)，满足 \(H'=h\) 的全纯函数 \(H:U\to\mathbb C\) 称为 \(h\) 的\term[yuanhanshu]{原函数}[primitive]。此处从已知幂级数构造了一个原函数，并没有断言一般开集上的全纯函数都有全局原函数。

\subsection{解析函数}
\label{b3:subsec:010203}

\Needspace{6\baselineskip}
\begin{definition}\label{b3:def:01-analytic}
设 \(\Omega\subseteq\mathbb C\) 为开集，\(f:\Omega\to\mathbb C\)。若对每个 \(a\in\Omega\)，都存在 \(r>0\)、\(D(a,r)\subseteq\Omega\) 及一个幂级数，使
\[
f(z)=\sum_{n=0}^\infty c_n(z-a)^n\quad(z\in D(a,r)),
\]
则称 \(f\) 在 \(\Omega\) 上\term[jiexi]{解析}[analytic]，也称它为\term[jiexi-hanshu]{解析函数}[analytic function]。
\end{definition}

解析性是每个点附近的局部性质，不要求在整个定义域上使用一个中心的级数。例如 \(1/(1-z)\) 在 \(\mathbb C\setminus\{1\}\) 的每个点都可用前面的几何级数展开，但以原点为中心的那一个级数只覆盖单位圆盘。定义要求的是可在各点重新展开，不是把一个展开式越过自己的收敛圆周。由\cref{b3:thm:01-power-differentiation}，解析函数必全纯，且它的局部展开就是以 \(f^{(n)}(a)/n!\) 为系数的\term[taylor-ji-shu]{Taylor 级数}[Taylor series]。

\Needspace{6\baselineskip}
\begin{proposition}\label{b3:prop:01-power-recenter}
设 \(a\in\mathbb C\)、\(c_n\in\mathbb C\)，幂级数 \(f(z)=\sum_{n=0}^{\infty}c_n(z-a)^n\) 的收敛半径为 \(R>0\)。则 \(f\) 在整个收敛圆盘 \(D(a,R)\) 上解析。更具体地，若 \(|b-a|<R\)，则对 \(|z-b|<R-|b-a|\)，
\[
f(z)=\sum_{k=0}^\infty\left(
\sum_{n=k}^\infty\binom nk c_n(b-a)^{n-k}\right)(z-b)^k.
\]
\end{proposition}
\begin{proof}
先固定 \(b\)，再取 \(0<r<R-|b-a|\)。对整个 \(|z-b|\le r\)，数 \(q=|b-a|+r\) 都小于 \(R\)。二项式公式给出同一个绝对值总和控制
\[
\sum_{n=0}^\infty\sum_{k=0}^n
\binom nk |c_n|\cdot|b-a|^{n-k}\cdot r^k
=\sum_{n=0}^\infty|c_n|q^n<\infty.
\]
每个固定 \(k\) 的系数级数也绝对收敛：上式中的第 \(k\) 列是可求和非负总和的一部分，再除以 \(r^k>0\) 即可。因此可以先按 \(n\) 展开 \(((b-a)+(z-b))^n\)，再按 \(k\) 合并：有限和重排不改变数值，绝对收敛使未计入部分的误差趋于零。误差控制对整个 \(|z-b|\le r\) 相同，故得到的展开在各闭子盘上一致收敛。内层系数按逐项微分公式等于 \(f^{(k)}(b)/k!\)，不依赖所选的 \(z\) 或 \(r\)。由于 \(r<R-|b-a|\) 任意，得到陈述中的开圆盘。这个半径只是由原收敛圆盘保证的下界；实际展开半径可以更大。
\end{proof}

实变量中的无限次可微并不保证局部 Taylor 展开等于原函数。考虑实函数
\[
 \psi(t)=
 \begin{cases}
 e^{-1/t^2},&t\ne0,\\
 0,&t=0.
 \end{cases}
\]
在 \(t\ne0\) 处，每阶导数都是 \(1/t\) 的某个多项式乘以 \(e^{-1/t^2}\)。对任意非负整数 \(m\)，有 \(|t|^{-m}e^{-1/t^2}\to0\)：令 \(s=1/t^2\to\infty\)，再由实指数级数中任意一个足够高次的正项控制 \(e^s\)，即可验证。依次把各阶导数在零点补为零，并用差商中多出的一个 \(1/t\) 因子作同样估计，便得到 \(\psi\in C^\infty(\mathbb R)\) 且 \(\psi^{(k)}(0)=0\) 对所有 \(k\ge0\) 成立。它的 Taylor 级数恒为零，函数却在 \(t\ne0\) 时严格为正。

\ProseParagraphBreak

复变量中要排除这种现象，必须把一个点的微分条件与它周围的函数值联系起来。下面的圆周积分公式承担这一任务；证明从复导数存在出发，不预设 \(f'\) 连续。

\subsection{全纯即解析}
\label{b3:subsec:010204}

三角形可以反复细分，因而能把一条边界的积分缩到局部微分余项所控制的尺度。圆盘的凸性又使两条线段之间恰好围出三角形，三角形积分消失便能给出圆盘内的原函数。接下来把这一方法用于差商，消去它在中心的一点例外，就得到圆周积分公式。\chapref{b3:ch:02}将把积分范围扩展到一般曲线，并讨论定义域的拓扑条件。

\ProseParagraphBreak

对分段 \(C^1\) 曲线 \(\gamma:[\alpha,\beta]\to\mathbb C\) 和其像上的连续函数 \(h\)，定义\term[fu-quxian-jifen]{复曲线积分}[complex line integral]为
\[
\int_\gamma h(\zeta)\,\mathrm{d}\zeta
\coloneqq\int_\alpha^\beta h\bigl(\gamma(t)\bigr)\cdot\gamma'(t)\,\mathrm{d}t.
\]
右端分为实部、虚部的一元积分。曲线长度为 \(\ell(\gamma)=\int_\alpha^\beta|\gamma'(t)|\,\mathrm{d}t\)，从而
\begin{equation}\label{b3:eq:01-contour-estimate}
\left|\int_\gamma h(\zeta)\,\mathrm{d}\zeta\right|
\leq\ell(\gamma)\cdot\max_{\zeta\in\gamma}|h(\zeta)|.
\end{equation}
反向行进使积分变号，连接曲线使积分相加；这些性质直接来自一元换元与积分可加性。本小节的三角形和圆周都取逆时针方向。
\symidx{\(\int_\gamma h(\zeta)\,\mathrm{d}\zeta\)：复曲线积分}

\ProseParagraphBreak

\term[goursat-dingli]{Goursat 定理}[Goursat's theorem]先给出三角形边界上的积分消失。

\Needspace{6\baselineskip}
\begin{lemma}[Goursat，三角形情形]\label{b3:lem:01-goursat-local}
设 \(T\subset\mathbb C\) 是一个非退化闭三角形，\(U\subset\mathbb C\) 为包含 \(T\) 的开集，\(f:U\to\mathbb C\) 全纯。边界 \(\partial T\) 取逆时针方向，则
\[
\int_{\partial T}f(\zeta)\,\mathrm{d}\zeta=0.
\]
\end{lemma}
\begin{proof}
将三边中点相连，把 \(T\) 分成四个相似小三角形。小三角形边界上的积分之和等于原积分，因为内部边以相反方向出现两次。故至少一个小三角形上的积分绝对值不小于原积分绝对值的四分之一。反复选择，得到嵌套闭三角形 \(T_n\)，其直径和周长均缩为原来的 \(2^{-n}\)，并有
\[
\left|\int_{\partial T}f\,\mathrm{d}\zeta\right|
\leq4^n\left|\int_{\partial T_n}f\,\mathrm{d}\zeta\right|.
\]
这些闭三角形非空且嵌套，并包含于紧集 \(T\)；紧性保证它们有公共点，直径趋零又保证公共点只有一个，记作 \(p\)。由于 \(U\) 包含整个闭三角形，即使 \(p\) 落在某条边上，\(f\) 仍在它的开邻域内复可微。在 \(p\) 处展开为
\[
f(\zeta)=f(p)+f'(p)(\zeta-p)+\varepsilon(\zeta)(\zeta-p),
\quad\varepsilon(\zeta)\to0\quad(\zeta\to p).
\]
在 \(p\) 处令 \(\varepsilon(p)=0\)。常数项与一次项在闭三角形上的积分为零，因为它们分别是 \(f(p)\zeta\) 和 \(f'(p)(\zeta-p)^2/2\) 的导数，沿各边使用一元微积分基本定理后端点项相消。对余项应用\eqref{b3:eq:01-contour-estimate}，得
\[
\left|\int_{\partial T}f\,\mathrm{d}\zeta\right|
\leq\ell(\partial T)\cdot\operatorname{diam}(T)
\cdot\sup_{\zeta\in T_n}|\varepsilon(\zeta)|.
\]
令 \(n\to\infty\)，右端趋于零，结论成立。
\end{proof}

细分使积分的可能损失为 \(4^n\)，而周长与直径各缩小 \(2^{-n}\)；三者恰好抵消，最后留下趋零的微分余项。这也解释了为什么只说“很小的三角形上积分很小”还不够：必须记录它相对于边长的阶数。证明只在所选极限点使用微分展开，没有对导数作积分，也未假定偏导数连续。

\Needspace{6\baselineskip}
\begin{lemma}\label{b3:lem:01-punctured-primitive}
设 \(D=D(a_0,R)\subset\mathbb C\)，其中 \(a_0\in\mathbb C\)、\(R>0\)，且 \(p\in D\)。若 \(h:D\to\mathbb C\) 连续，在 \(D\setminus\{p\}\) 上全纯，则存在全纯函数 \(H:D\to\mathbb C\) 满足 \(H'=h\)。因此 \(h\) 沿 \(D\) 内任意分段 \(C^1\) 闭曲线的积分为零。
\end{lemma}
\begin{proof}
先证明 \(h\) 沿每个闭三角形 \(T\subset D\) 的边界积分为零。把 \(T\) 均匀分成 \(4^n\) 个直径为原来 \(2^{-n}\) 的小三角形。不接触 \(p\) 的小三角形由\cref{b3:lem:01-goursat-local}积分为零。重复中点细分所得网格中，一个点若在小三角形内部，只属于一个闭三角形；若在边内，至多属于两个；若在顶点，至多属于六个。因而接触 \(p\) 的小三角形数目与细分次数无关。这些可能不为零的积分之和，绝对值不超过
\[
6\cdot2^{-n}\ell(\partial T)\cdot\max_T|h|.
\]
内部边仍两两抵消。\(h\) 连续保证最大值有限，令 \(n\to\infty\) 即得三角形积分为零。若 \(p\notin T\)，直接用前一引理即可。

固定 \(a\in D\)，试图把 \(h\) 沿路径积累起来。圆盘内每个点都能由一条线段从 \(a\) 到达，因此先定义 \(H(w)=\int_{[a,w]}h(\zeta)\,\mathrm{d}\zeta\)。圆盘凸性保证涉及的线段和三角形都留在 \(D\) 中。三角形积分为零给出
\[
H(w+v)-H(w)=\int_{[w,w+v]}h(\zeta)\,\mathrm{d}\zeta
=v\int_0^1h(w+tv)\,\mathrm{d}t.
\]
三点共线时，同一等式由线段积分的可加性得到。除以 \(v\neq0\)，由连续性令 \(v\to0\)，便有 \(H'(w)=h(w)\)。由于 \(H'=h\) 连续，\(H\) 作为实映射具有连续微分。沿每一段 \(C^1\) 曲线应用实参数链式法则，有 \((H\circ\gamma)'=h(\gamma)\gamma'\)；一元微积分基本定理使积分成为端点差，各段连接点的项相消，闭曲线的起点与终点又相同，所以积分为零。
\end{proof}

前一引理暂时允许一个点处未能确认复可微。它没有先声称 \(h\) 在该点也全纯，而是直接用连续性让接触该点的短边积分趋零；随后用三角形积分消失构造 \(H\)。这个余地使下一步能够处理差商在分母为零处的补定义。

\ProseParagraphBreak

接着建立\term[cauchy-jifen-gongshi]{Cauchy 积分公式}[Cauchy's integral formula]在圆盘内的形式。

\Needspace{6\baselineskip}
\begin{lemma}[Cauchy 积分公式，圆盘情形]\label{b3:lem:01-disk-cauchy}
设 \(a\in\mathbb C\)、\(R>0\)，\(f\) 在 \(D(a,R)\) 上全纯。若 \(0<r<R\)，则对 \(|z-a|<r\)，
\begin{equation}\label{b3:eq:01-disk-cauchy}
f(z)=\frac1{2\pi\mathrm{i}}\int_{|\zeta-a|=r}
\frac{f(\zeta)}{\zeta-z}\,\mathrm{d}\zeta.
\end{equation}
圆周取逆时针方向。
\end{lemma}
\begin{proof}
固定 \(z\)，令
\[
h(\zeta)=\begin{cases}
\dfrac{f(\zeta)-f(z)}{\zeta-z},&\zeta\neq z,\\[4pt]
f'(z),&\zeta=z.
\end{cases}
\]
当 \(\zeta\to z\) 时，商式按复导数定义趋于 \(f'(z)\)，所以 \(h\) 在这个特殊点连续；在其余点，连续性与全纯性都由商法则得到。由\cref{b3:lem:01-punctured-primitive}，\(h\) 的圆周积分为零。因此只需证明
\[
\int_{|\zeta-a|=r}\frac{\mathrm{d}\zeta}{\zeta-z}=2\pi\mathrm{i}.
\]
在该圆周上，几何级数
\[
\frac1{\zeta-z}=\sum_{n=0}^\infty
\frac{(z-a)^n}{(\zeta-a)^{n+1}}
\]
一致收敛，因为 \(|z-a|/r<1\)。具体地，前 \(N+1\) 项之后的尾项模在整个圆周上不超过
\[
 \frac1r\frac{(|z-a|/r)^{N+1}}{1-|z-a|/r},
\]
趋零速度与圆周上的点无关。由\eqref{b3:eq:01-contour-estimate}，积分误差至多再乘圆周长 \(2\pi r\)，故可逐项积分。\(n=0\) 项沿 \(\zeta=a+r(\cos t+\mathrm{i}\sin t)\)、\(0\leq t\leq2\pi\) 参数化，积分为 \(2\pi\mathrm{i}\)；\(n\geq1\) 时，\((\zeta-a)^{-n-1}\) 有原函数 \(-\frac1n(\zeta-a)^{-n}\)，其闭曲线积分为零。把结果代回 \(h\) 的零积分，得到公式。
\end{proof}

公式\eqref{b3:eq:01-disk-cauchy}表明，边界圆周上的函数值决定了圆盘内部的函数值。将核函数 \((\zeta-z)^{-1}\) 展开为几何级数，就能得到全纯函数的幂级数表示。

\Needspace{6\baselineskip}
\begin{theorem}[全纯与解析的等价性]\label{b3:thm:01-holomorphic-analytic}
设 \(U\subseteq\mathbb C\) 为开集，\(f:U\to\mathbb C\)。则 \(f\) 全纯，当且仅当它解析。若 \(a\in U\)、\(R>0\)、\(D(a,R)\subseteq U\)，且 \(f\) 全纯，则
\begin{equation}\label{b3:eq:01-holomorphic-taylor}
f(z)=\sum_{n=0}^\infty\frac{f^{(n)}(a)}{n!}(z-a)^n
\quad(|z-a|<R).
\end{equation}
此外，给定 \(0<\rho<r<R\)，令 \(M_r=\max_{|\zeta-a|=r}|f(\zeta)|\)。对每个整数 \(N\ge0\)，有
\begin{equation}\label{b3:eq:01-taylor-remainder}
 \sup_{|z-a|\le\rho}
 \left|f(z)-\sum_{n=0}^{N}\frac{f^{(n)}(a)}{n!}(z-a)^n\right|
 \le M_r\frac{(\rho/r)^{N+1}}{1-\rho/r}.
\end{equation}
\end{theorem}
\begin{proof}
解析推出全纯已由逐项求导证明。反之，设 \(f\) 全纯，固定 \(0<r<R\)。在\eqref{b3:eq:01-disk-cauchy}中展开刚才的几何级数；在任意 \(|z-a|\leq\rho<r\) 上，它关于圆周变量一致收敛，乘以连续有界的 \(f(\zeta)\) 后仍可逐项积分。于是
\[
f(z)=\sum_{n=0}^\infty c_n(z-a)^n,
\quad c_n=\frac1{2\pi\mathrm{i}}\int_{|\zeta-a|=r}
\frac{f(\zeta)}{(\zeta-a)^{n+1}}\,\mathrm{d}\zeta.
\]
若 \(M_r=\max_{|\zeta-a|=r}|f(\zeta)|\)，则 \(|c_n|\leq M_r/r^n\)，所以所得幂级数至少在半径 \(r\) 内收敛。由逐项求导及系数唯一性，\(c_n=f^{(n)}(a)/n!\)，不依赖选择的 \(r\)。任意 \(|z-a|<R\) 都可放在一个 \(r<R\) 的圆盘内，故得到\eqref{b3:eq:01-holomorphic-taylor}。由系数界，余项在 \(|z-a|\le\rho\) 上的模不超过 \(M_r\sum_{n>N}(\rho/r)^n\)，求和即得\eqref{b3:eq:01-taylor-remainder}。对一般开集，在每一点内取一个这样的圆盘即可。
\end{proof}

全纯函数因此拥有任意阶复导数，各阶导数仍全纯且连续。由 Cauchy--Riemann 关系，\(\partial_x f=f'\)、\(\partial_y f=\mathrm{i}f'\)；再对每个全纯的导数重复这个等式，可得
\[
 \partial_x^j\partial_y^k f=\mathrm{i}^{\,k}f^{(j+k)}
 \quad(j,k\ge0).
\]
所以实部、虚部也都是无限次连续可微的。幂级数表示把最初的一次复可微提升为任意阶连续可微。

\ProseParagraphBreak

在实际使用中，展开半径要区分“已保证的范围”和“最大的范围”。若只知道 \(f\) 在 \(D(a,R)\) 全纯，定理保证中心 Taylor 级数至少覆盖这个圆盘，不能推出半径恰为 \(R\)。例如同一函数 \(1/(1-z)\) 即使暂时只在 \(D(0,1/2)\) 上讨论，它的中心展开仍有半径一。另一方面，取 \(r<R\) 的圆周，是为了圆周本身处于定义域内部；不曾假定函数在 \(|z-a|=R\) 有连续边值。

\ProseParagraphBreak

本小节先证明了局部积分结论。\chapref{b3:ch:02}将在相应的拓扑条件下，把三角形上的结论推广到一般曲线，并系统讨论积分公式的推论；本章接下来则先用解析性构造初等函数、研究零点和局部映射的形状。

% !TeX root = ../../Book3.tex
\section{初等全纯函数}
\label{b3:sec:0103}

幂级数提供了构造全纯函数的直接方法。本节先把实变量的指数与三角函数延伸到复平面，再处理指数函数的局部逆。对数的困难来自不同点可以有同一个指数值：局部求逆没有障碍，沿闭曲线作全局选择却可能失败。最后讨论的 Möbius 变换则把代数运算与圆、直线的几何联系起来。

\subsection{指数函数}
\label{b3:subsec:010301}

\begin{definition}
\label{b3:def:01-elementary-exp}
对 \(z\in\mathbb C\)，令
\[
\exp z\coloneqq\sum_{n=0}^{\infty}\frac{z^n}{n!},
\quad e^z\coloneqq\exp z.
\]
这个函数称为复\term[zhishu-hanshu]{指数函数}[exponential function]。\symidx{\(\exp z=e^z\)：复指数函数}
\end{definition}

对每个固定的 \(R>0\)，相邻控制项 \(R^{n+1}/(n+1)!\) 与 \(R^n/n!\) 的比值趋于零。因此级数在每个闭圆盘上一致绝对收敛；由\cref{b3:thm:01-power-differentiation}，它在整个复平面上全纯，并可任意次逐项求导。实轴上的定义与实变量指数函数相同，这也可由实 Taylor 公式及其余项估计核对。下面的\term[euler-gongshi]{Euler 公式}[Euler's formula]进一步说明复指数与实三角函数之间的联系。

\begin{proposition}[Euler 公式]
\label{b3:prop:01-elementary-exp}
复指数函数满足
\[
(e^z)'=e^z,
\quad e^{z+w}=e^z\cdot e^w,
\quad e^z\neq0.
\]
若 \(x,y\in\mathbb R\)，则
\begin{equation}
\label{b3:eq:01-elementary-euler}
e^{x+\mathrm{i}y}=e^x\cdot(\cos y+\mathrm{i}\sin y),
\quad |e^{x+\mathrm{i}y}|=e^x.
\end{equation}
此外，\(e^z=e^w\) 当且仅当 \(z-w\in2\pi\mathrm{i}\mathbb Z\)。
\end{proposition}

\begin{proof}
逐项求导后把指标减一，得到导数公式。两个指数级数的绝对值级数均收敛，故其乘积可按总次数重新求和；总次数为 \(n\) 的项为
\[
\sum_{k=0}^{n}\frac{z^kw^{n-k}}{k!(n-k)!}
=\frac{(z+w)^n}{n!}.
\]
这里重新求和的合法性来自绝对尾项控制：保留足够大的有限矩形后，遗漏项绝对值之和可由两个收敛的非负级数的尾和控制到任意小。于是得到加法公式。令 \(w=-z\)，可知 \(e^z\cdot e^{-z}=1\)，从而指数函数没有零点。

把 \(e^{\mathrm{i}y}\) 的偶数次项与奇数次项分别收集，得到实余弦与实正弦的 Taylor 级数。实正弦、余弦的各阶导数绝对值不超过 \(1\)，Taylor 余项趋于零，因此
\[
e^{\mathrm{i}y}=\cos y+\mathrm{i}\sin y.
\]
再用加法公式及实恒等式 \(\cos^2y+\sin^2y=1\)，便得到\eqref{b3:eq:01-elementary-euler}。若 \(e^{z-w}=1\)，模长公式先给出 \(\operatorname{Re}(z-w)=0\)，实三角函数的取值再给出 \(\operatorname{Im}(z-w)\in2\pi\mathbb Z\)。反向结论直接代入即可。
\end{proof}

式\eqref{b3:eq:01-elementary-euler}把复指数分成径向伸缩与转动两部分。特别地，\(2\pi\mathrm{i}\) 是指数函数的周期，而指数函数的导数处处不为零。因此，导数不消失只能保证局部的可逆性，不能保证在整个定义域上单射。

\subsection{对数函数}
\label{b3:subsec:010302}

对非零复数 \(z\)，方程 \(e^w=z\) 的解的实部必须等于 \(\log|z|\)，虚部则是 \(z\) 的一个辐角；所有解恰好相差 \(2\pi\mathrm{i}\) 的整数倍。要把这些解组成一个函数，需要连续地选取辐角。

\begin{definition}
\label{b3:def:01-elementary-log}
在割平面 \(D\coloneqq\mathbb C\setminus(-\infty,0]\) 上，取满足
\[
z=|z|e^{\mathrm{i}\theta},\quad -\pi<\theta<\pi
\]
的唯一实数 \(\theta\)，记为 \(\operatorname{Arg}z\)，称为\term[fujiao-zhuzhi]{辐角主值}[principal argument]。函数
\[
\operatorname{Log}z\coloneqq\log|z|+\mathrm{i}\operatorname{Arg}z
\]
称为\term[duishu-zhuzhi]{对数主值}[principal logarithm]。\symidx{\(\operatorname{Arg}z\)：取值于 \((-\pi,\pi)\) 的辐角主值}\symidx{\(\operatorname{Log}z\)：割平面上的对数主值}
\end{definition}

这里实数的对数仍记为 \(\log\)。本书把负实轴连同原点从主值的定义域中删去；另有把非零复数的辐角主值取在 \((-\pi,\pi]\) 的约定。\footnote{参见王以忠《工程复变函数与积分变换》（清华大学出版社，2021）\href{https://www.tup.tsinghua.edu.cn/upload/books/yz/091462-01.pdf}{出版社样章}印刷第 8 页。该处使用 \((-\pi,\pi]\) 选取点值；本节为讨论全纯分支而删去负实轴。}给负实数指定端点辐角，不能使主值在那里连续。

\begin{proposition}
\label{b3:prop:01-elementary-principal-log}
对数主值在 \(D\) 上全纯，并且
\[
\operatorname{Log}1=0,
\quad e^{\operatorname{Log}z}=z,
\quad (\operatorname{Log}z)'=\frac1z.
\]
当且仅当 \(-\pi<\operatorname{Im}w<\pi\) 时，有 \(e^w\in D\) 且 \(\operatorname{Log}(e^w)=w\)。
\end{proposition}

\begin{proof}
写 \(z=x+\mathrm{i}y\)、\(r=(x^2+y^2)^{1/2}\)。在 \(D\) 上有 \(r+x>0\)，实半角公式给出
\[
\operatorname{Arg}(x+\mathrm{i}y)
=2\arctan\frac{y}{r+x}.
\]
右边取值于 \((-\pi,\pi)\)，其余弦为 \(x/r\)、正弦为 \(y/r\)，所以确实是定义中的唯一辐角。这个表达式还证明辐角主值是光滑的。对其求实偏导，得到
\[
(\operatorname{Arg}z)_x=-\frac{y}{r^2},
\quad (\operatorname{Arg}z)_y=\frac{x}{r^2}.
\]
另一方面，\(\log r\) 的两个偏导为 \(x/r^2\) 与 \(y/r^2\)。它们满足 Cauchy--Riemann 方程，由\cref{b3:cor:01-continuous-partial-criterion}可知主对数全纯，且其复导数为
\[
\frac{x}{r^2}-\mathrm{i}\frac{y}{r^2}=\frac1z.
\]
其余断言由定义及\cref{b3:prop:01-elementary-exp}得出：指数值唯一确定实部，并把虚部确定到模 \(2\pi\)；主值恰选取位于 \((-\pi,\pi)\) 的代表。
\end{proof}

主值不能无条件保留实对数的乘法规则。例如取 \(z=e^{3\pi\mathrm{i}/4}\)，则 \(z^2=-\mathrm{i}\)，所以
\[
\operatorname{Log}(z^2)=-\frac{\pi\mathrm{i}}2,
\quad 2\operatorname{Log}z=\frac{3\pi\mathrm{i}}2.
\]
两者相差一个完整的周期。类似地，当正数 \(r\) 固定时，从负实轴的上侧和下侧趋近 \(-r\)，主对数分别趋于 \(\log r+\pi\mathrm{i}\) 与 \(\log r-\pi\mathrm{i}\)。这解释了为什么不能把两侧的主值直接接起来。

\begin{definition}
\label{b3:def:01-elementary-log-branch}
设 \(U\subseteq\mathbb C\setminus\{0\}\) 是开集。若连续函数 \(L\colon U\rightarrow\mathbb C\) 满足 \(e^{L(z)}=z\)，则称 \(L\) 为 \(U\) 上的一个\term[duishu-fenzhi]{对数分支}[branch of the logarithm]。
\end{definition}

\begin{proposition}
\label{b3:prop:01-elementary-local-log}
每个非零点都有一个邻域，在该邻域上存在全纯的对数分支。任何连续的对数分支自动全纯，导数为 \(1/z\)；在连通开集上，两个分支之差是某个固定的 \(2\pi\mathrm{i}k\)，其中 \(k\in\mathbb Z\)。但整个 \(\mathbb C\setminus\{0\}\) 上不存在连续的对数分支。
\end{proposition}

\begin{proof}
固定 \(z_0\neq0\)，选取 \(\ell_0\) 使 \(e^{\ell_0}=z_0\)。当 \(|z-z_0|<|z_0|\) 时，\(z/z_0\) 的实部为正，故
\[
L_0(z)=\ell_0+\operatorname{Log}(z/z_0)
\]
有定义，且是所需的局部分支。若 \(L\) 是同一邻域上的任意连续分支，则 \((L-L_0)/(2\pi\mathrm{i})\) 是整数值连续函数，在较小的连通邻域上恒定。因此 \(L\) 局部全纯且 \(L'=1/z\)。同一理由说明连通开集上两个分支的差恒定。

若整个去心平面上存在连续分支 \(L\)，令
\[
h(t)=\frac{L(e^{\mathrm{i}t})-\mathrm{i}t}{2\pi\mathrm{i}},
\quad 0\leq t\leq2\pi.
\]
由指数函数的周期刻画，\(h\) 是连续的整数值函数，因而恒定。但 \(e^{2\pi\mathrm{i}}=1\) 使 \(h(2\pi)=h(0)-1\)，矛盾。
\end{proof}

有了一个分支 \(L\)，便可在其定义域上令 \(z^{\alpha}=\exp(\alpha L(z))\)。改变分支会把这个值乘以 \(e^{2\pi\mathrm{i}k\alpha}\)。整数幂与分支无关，平方根等非整数幂则必须交代所选分支。\subsecref{b3:subsec:010404}构造一个非零全纯函数的局部根时，只需在像靠近某个非零值的小邻域内使用这里的局部分支。

\subsection{三角函数}
\label{b3:subsec:010303}

\begin{definition}
\label{b3:def:01-elementary-trigonometric}
在复平面上定义\term[zhengxian-hanshu]{正弦函数}[sine function]与\term[yuxian-hanshu]{余弦函数}[cosine function]为
\[
\sin z\coloneqq\frac{e^{\mathrm{i}z}-e^{-\mathrm{i}z}}{2\mathrm{i}},
\quad \cos z\coloneqq\frac{e^{\mathrm{i}z}+e^{-\mathrm{i}z}}2.
\]
\symidx{\(\sin z\)：复正弦函数}\symidx{\(\cos z\)：复余弦函数}
\end{definition}

由 Euler 公式，这些定义在实轴上还原通常的正弦与余弦。求导及相乘给出
\[
(\sin z)'=\cos z,
\quad (\cos z)'=-\sin z,
\quad \sin^2z+\cos^2z=1.
\]
例如最后一个等式来自
\[
\frac{(e^{\mathrm{i}z}+e^{-\mathrm{i}z})^2
-(e^{\mathrm{i}z}-e^{-\mathrm{i}z})^2}{4}=1.
\]
把指数的加法公式代入定义，还得到
\[
\begin{aligned}
\sin(z+w)&=\sin z\cdot\cos w+\cos z\cdot\sin w,\\
\cos(z+w)&=\cos z\cdot\cos w-\sin z\cdot\sin w.
\end{aligned}
\]
同样，由指数级数的奇偶项分别得到
\[
\sin z=\sum_{n=0}^{\infty}\frac{(-1)^nz^{2n+1}}{(2n+1)!},
\quad
\cos z=\sum_{n=0}^{\infty}\frac{(-1)^nz^{2n}}{(2n)!}.
\]
它们都是整个复平面上的全纯函数。由指数函数的周期刻画，\(\sin z=0\) 恰在 \(z\in\pi\mathbb Z\) 时发生，\(\cos z=0\) 恰在 \(z\in\pi/2+\pi\mathbb Z\) 时发生；例如前者等价于 \(e^{2\mathrm{i}z}=1\)，后者等价于 \(e^{2\mathrm{i}z}=-1\)。

\ProseParagraphBreak
复三角函数保留上述恒等式，却不保留实轴上的有界性。沿虚轴有
\[
\sin(\mathrm{i}y)=\frac{\mathrm{i}}2(e^y-e^{-y}),
\quad \cos(\mathrm{i}y)=\frac12(e^y+e^{-y}),
\quad y\in\mathbb R.
\]
当 \(y\rightarrow+\infty\) 时，两者的模都以指数速度增长。恒等式 \(\sin^2z+\cos^2z=1\) 中的平方没有复共轭，因而不能推出 \(|\sin z|^2+|\cos z|^2=1\)。计算复函数的大小时，代数恒等式与含模估计必须分别处理。

\subsection{Möbius 变换}
\label{b3:subsec:010304}

分式线性函数可能在一个有限点没有有限值。把这个值解释为无穷远点以后，它们可以作为整个扩充平面上的可逆映射来讨论。

\begin{definition}
\label{b3:def:01-elementary-extended-plane}
在复平面上添加一个点 \(\infty\)，得到\term[kuochong-fu-pingmian]{扩充复平面}[extended complex plane]
\[
% 零宽核避免 XeTeX 对 AMS 单字母加 unicode-math 重音时的字体错误。
\widehat{\mathbb C\kern0pt}\coloneqq\mathbb C\cup\{\infty\}.
\]
\symidx{\(\widehat{\mathbb C\kern0pt}\)：扩充复平面}
无穷远点的基本邻域取为 \(\{\infty\}\cup\{z:|z|>R\}\)，其中 \(R>0\)。在这样的邻域内使用坐标 \(\zeta=1/z\)，并规定无穷远点对应 \(\zeta=0\)。
\end{definition}

讨论映射在无穷远点附近的连续性或全纯性，是指改用上述坐标以后讨论相应的普通函数；若函数值也为无穷远，则目标同样使用倒数坐标。例如 \(z\mapsto1/z\) 在扩充平面上交换 \(0\) 与 \(\infty\)，在相应坐标中就是恒等函数。本节只需这两张坐标图，不预先使用一般曲面的理论。

\begin{definition}
\label{b3:def:01-elementary-mobius}
设 \(a,b,c,d\in\mathbb C\) 且 \(ad-bc\neq0\)。在分母不为零处定义
\[
T(z)=\frac{az+b}{cz+d}.
\]
若 \(c\neq0\)，补定义 \(T(-d/c)=\infty\)、\(T(\infty)=a/c\)；若 \(c=0\)，补定义 \(T(\infty)=\infty\)。这样得到的扩充平面上的映射称为\term[mobius-bianhuan]{Möbius 变换}[Möbius transformation]，也称分式线性变换。\footnote{“分式线性变换”见王晓光《复变函数》（高等教育出版社，2025）第十三章的\href{https://xuanshu.hep.com.cn/front/h5Mobile/bookDetails?bookId=68598849e119ac9729fe438b}{出版社目录}。}\termsee[fenshi-xianxing-bianhuan]{分式线性变换}{Möbius 变换}
\end{definition}

\begin{proposition}
\label{b3:prop:01-elementary-mobius-inverse}
每个 Möbius 变换都是扩充平面上的连续双射，其逆变换为
\[
T^{-1}(w)=\frac{dw-b}{-cw+a},
\]
其中无穷值按定义中的规则解释。在有限且函数值有限的点处，
\[
T'(z)=\frac{ad-bc}{(cz+d)^2}\neq0.
\]
在源与目标的其他坐标中，它也全纯且导数不为零。
\end{proposition}

\begin{proof}
令 \(\Delta=ad-bc\)。从 \(w(cz+d)=az+b\) 解出 \(z\)，得到逆变换公式；其系数行列式也是 \(\Delta\neq0\)，且代入补定义的特殊点仍互为逆。因此映射是双射，普通点处的连续性和导数公式由商法则得到。

若 \(c\neq0\)，在 \(z_0=-d/c\) 附近给目标使用倒数坐标，所得函数为
\[
\frac1{T(z)}=\frac{cz+d}{az+b}.
\]
分母在 \(z_0\) 处等于 \(-\Delta/c\neq0\)，导数在该点为 \(-c^2/\Delta\neq0\)。在源点 \(\infty\) 附近令 \(z=1/\zeta\)，则
\[
T(1/\zeta)=\frac{a+b\zeta}{c+d\zeta},
\]
右边在 \(\zeta=0\) 处全纯，导数为 \(-\Delta/c^2\neq0\)。若 \(c=0\)，则 \(a,d\neq0\)，在源与目标同时使用倒数坐标，得到
\[
\frac1{T(1/\zeta)}=\frac{d\zeta}{a+b\zeta},
\]
其在零点处的导数为 \(d/a\neq0\)。这些式子同时证明所补无穷值处的连续性及局部全纯性。
\end{proof}

下面把圆与直线放在一起讨论。圆取通常的有限半径圆，直线则补入无穷远点；这一约定使 Möbius 变换的几何性质能够统一表述。

\begin{proposition}
\label{b3:prop:01-elementary-mobius-circles}
Möbius 变换把每个圆或直线双射到一个圆或直线。
\end{proposition}

\begin{proof}
圆或直线的有限部分统一写成
\[
A|z|^2+Bz+\overline B\,\overline z+C=0,
\quad A,C\in\mathbb R,
\quad |B|^2-AC>0.
\]
若 \(A\neq0\)，配方后这是圆心为 \(-\overline B/A\)、半径为 \(\sqrt{|B|^2-AC}/|A|\) 的圆；若 \(A=0\)，则 \(B\neq0\)，方程给出一条直线，并按约定补入 \(\infty\)。

平移和非零复数乘法把圆变为圆、直线变为直线：例如代入 \(z=(w-\beta)/\alpha\)、乘以 \(|\alpha|^2\)，所得方程仍具同一形式，且新系数的判别量为 \(|\alpha|^2\cdot\bigl(|B|^2-AC\bigr)>0\)。对于倒数变换 \(w=1/z\)，在 \(w\neq0\) 处代入并乘以 \(|w|^2\)，得到
\[
C|w|^2+\overline B w+B\overline w+A=0.
\]
判别量没有改变。原方程含有 \(0\) 当且仅当 \(C=0\)，恰与像方程表示直线、因而含有 \(\infty\) 对应；原方程表示直线当且仅当 \(A=0\)，恰与像中含有 \(0\) 对应。所以特殊点也包含在这个对应中。

当 \(c=0\) 时，\(T\) 本身是非退化的仿射变换；当 \(c\neq0\) 时，可以分解为
\[
T(z)=\frac ac-\frac{ad-bc}{c^2}\cdot\frac1{z+d/c}.
\]
这是平移、倒数变换、非零复数乘法与平移的复合。前面的计算证明其像是圆或直线，逆变换又保证映射到的是整个圆或整条直线。
\end{proof}

一个常用例子是
\[
T(z)=\frac{z-\mathrm{i}}{z+\mathrm{i}}.
\]
当 \(\operatorname{Im}z>0\) 时，\(|z+\mathrm{i}|^2-|z-\mathrm{i}|^2=4\operatorname{Im}z>0\)，故 \(|T(z)|<1\)。反过来，由 \(z=\mathrm{i}(1+w)/(1-w)\) 算出
\[
\operatorname{Im}z=\frac{1-|w|^2}{|1-w|^2}>0
\quad (|w|<1).
\]
因此这个变换把上半平面双射到单位圆盘；它把扩充实轴送到单位圆周，无穷远点对应圆周上的 \(1\)。借助这个映射，可以将半平面上的问题化为圆盘上的问题，或反过来处理。

% !TeX root = ../../Book3.tex
\section{局部结构}
\label{b3:sec:0104}

局部幂级数把函数在一点附近的行为分成两部分：最先出现的非零项确定主要形状，其后的项给出可控制的修正。先用这种分解研究零点，再把局部相等沿连通区域传开；最后选择新的局部坐标，将非常值全纯函数化为一个幂映射。这样既能解释导数不为零时的可逆性，也能看清导数为零时原像怎样合并。

\subsection{孤立零点}
\label{b3:subsec:010401}

\Needspace{6\baselineskip}
\begin{definition}\label{b3:def:01-zero-order}
设 \(\Omega\subseteq\mathbb C\) 为开集，\(a\in\Omega\)，\(f:\Omega\to\mathbb C\) 全纯且 \(f(a)=0\)。若存在正整数 \(m\)，使
\[
f(a)=f'(a)=\cdots=f^{(m-1)}(a)=0,
\quad f^{(m)}(a)\neq0,
\]
则称 \(a\) 是 \(f\) 的 \(m\) 阶零点，\(m\) 称为\term[lingdian-jieshu]{零点阶数}[order of a zero]，记为 \(\operatorname{ord}_a f\)。\(m=1\) 时称为\term[dan-lingdian]{单零点}[simple zero]。若 \(a\) 的某邻域内没有 \(f\) 的其他零点，则称 \(a\) 为\term[guli-lingdian]{孤立零点}[isolated zero]。
\end{definition}
\symidx{\(\operatorname{ord}_a f\)：全纯函数在 \(a\) 处的零点阶数}

\Needspace{6\baselineskip}
\begin{proposition}\label{b3:prop:01-zero-factorization}
设 \(\Omega\subseteq\mathbb C\) 为开集，\(a\in\Omega\)，\(f:\Omega\to\mathbb C\) 全纯且 \(f(a)=0\)。则恰有以下两种情形之一：\(f\) 在 \(a\) 的一个邻域内恒为零；或存在唯一正整数 \(m\)、圆盘 \(D(a,R)\subseteq\Omega\) 及其上的全纯函数 \(g\)，使
\begin{equation}\label{b3:eq:01-zero-factorization}
f(z)=(z-a)^m\cdot g(z),
\quad g(a)\neq0.
\end{equation}
后一种情形中，\(a\) 是孤立零点，且阶数为 \(m\)；在所取圆盘上，函数 \(g\) 也由 \(f\) 唯一确定。
\end{proposition}
\begin{proof}
取以 \(a\) 为中心、包含于 \(\Omega\) 的圆盘，在其中使用\eqref{b3:eq:01-holomorphic-taylor}。若所有系数都为零，函数就在该圆盘内恒为零。否则，非零系数的指标组成非空的非负整数集，有最小元；又因 \(f(a)=0\)，这个最小元 \(m\) 至少为一。提出 \((z-a)^m\)，余下级数
\[
g(z)=\sum_{k=0}^\infty\frac{f^{(m+k)}(a)}{(m+k)!}(z-a)^k
\]
在原圆盘内部的每个较小闭圆盘上绝对一致收敛，因而全纯，并且 \(g(a)=f^{(m)}(a)/m!\neq0\)。由连续性，在更小的圆盘内有 \(|g(z)|\geq|g(a)|/2>0\)，所以其中只有 \(a\) 一个零点。系数唯一性确定 \(m\)；在 \(z\neq a\) 处 \(g\) 由商式唯一确定，在 \(a\) 处再由连续性确定。
\end{proof}

例如 \(e^z-1\) 在原点的一阶系数为 \(1\)，所以原点是单零点；\(1-\cos z=z^2/2+\cdots\) 在原点的零点阶数为 \(2\)。若两个全纯函数 \(f,g\) 在 \(a\) 的零点阶数分别为 \(m,n\)，相乘其局部分解便有
\[
\operatorname{ord}_a(fg)=m+n.
\]
求和则可能发生首项抵消。例如 \(e^z-1\) 和 \(-z\) 都在原点一阶消失，它们的和 \(e^z-1-z\) 却在原点二阶消失。阶数记录的是实际留下的第一个非零项，不能只根据两个输入阶数决定和的阶数。

\ProseParagraphBreak

阶数还给出邻域内的双边估计。若 \(f(z)=(z-a)^m\cdot g(z)\)、\(g(a)\ne0\)，由连续性缩小半径，可使
\[
 \frac{|g(a)|}{2}\cdot|z-a|^m\le |f(z)|
 \le\frac{3|g(a)|}{2}\cdot|z-a|^m.
\]
常数对这个圆盘中的所有点同时有效，所以有限阶零点沿任意方向都以同一整数次幂消失。以 \(e^z-1-z\) 为例，除以 \(z^2\) 后的全纯因子在原点取值 \(1/2\)，因而其模在充分小的复圆盘内与 \(|z|^2\) 双边可比。

\ProseParagraphBreak

反过来，设 \(f\) 在 \(a\) 附近全纯，\(m\) 为非负整数，\(\alpha>0\)。若存在常数 \(C>0\)，使在趋于 \(a\) 的一列不同点 \(z_n\ne a\) 上有
\[
 |f(z_n)|\le C\cdot|z_n-a|^{m+\alpha},
\]
则连续性先给 \(f(a)=0\)。若它的有限零点阶数为 \(k\le m\)，写 \(f(z)=(z-a)^k\cdot g(z)\)、\(g(a)\ne0\)，上面的下界会给出
\[
 0<\frac{|g(a)|}{2}
 \le C\cdot|z_n-a|^{m+\alpha-k}\longrightarrow0,
\]
矛盾。因此 \(f\) 的阶数至少为 \(m+1\)，或它局部恒零。特别地，对已经知道全纯的函数，沿实轴趋近取得这样的估计也足够；这里的依据正是全纯性事先给出的邻域分解。对一般实可微函数则不能这样推断，例如 \(z\mapsto\operatorname{Im}z\) 在实轴恒为零，在虚方向却只一阶消失。单条路径上的信息与整个邻域的统一结论之间，必须有相应的结构保证。

\subsection{恒等定理}
\label{b3:subsec:010402}

一个零点的附近可以由幂级数控制；要把局部信息传遍定义域，还需要连通性。\term[hengdeng-dingli]{恒等定理}[identity theorem]，也称唯一性定理或惟一性定理，\footnote{王晓光《复变函数》（高等教育出版社，2025）第 11.3 节的\href{https://xuanshu.hep.com.cn/front/h5Mobile/bookDetails?bookId=68598849e119ac9729fe438b}{出版社目录}用“唯一性定理”；钟玉泉《复变函数论》（第五版，高等教育出版社，2021）第四章第 4 节的\href{https://xuanshu.hep.com.cn/front/book/findBookDetails?bookId=60202746966feeeca73340f0}{出版社目录}用“惟一性定理”。}\termsee[weiyixing-dingli]{唯一性定理}{恒等定理}\termsee[weiyixing-dingli-old]{惟一性定理}{恒等定理}说明，在区域内部积累的一列相等函数值，已经足以确定整个全纯函数。

\Needspace{6\baselineskip}
\begin{theorem}[恒等定理]\label{b3:thm:01-identity}
设 \(\Omega\subseteq\mathbb C\) 为区域，\(f,g:\Omega\to\mathbb C\) 全纯。若集合 \(\{z\in\Omega:f(z)=g(z)\}\) 在 \(\Omega\) 内有一个聚点，则 \(f=g\) 于整个 \(\Omega\)。
\end{theorem}
\begin{proof}
令 \(h=f-g\)，设所述聚点为 \(a\in\Omega\)。由连续性，\(h(a)=0\)。\Cref{b3:prop:01-zero-factorization}排除了有限阶零点，因为那会使 \(a\) 附近没有其他零点。因此 \(h\) 在 \(a\) 的某邻域内恒为零。

令 \(E\) 为 \(\Omega\) 中那些具有恒零邻域的点组成的集合。它非空且开；由 Taylor 展开，亦可写为
\[
E=\bigcap_{k=0}^\infty\{z\in\Omega:h^{(k)}(z)=0\}.
\]
确实，恒零邻域使所有导数消失，反过来所有 Taylor 系数为零又给出恒零邻域。各阶导数连续，所以 \(E\) 在 \(\Omega\) 中闭。区域连通性使非空的开闭子集 \(E\) 等于 \(\Omega\)，结论得证。
\end{proof}

\Needspace{6\baselineskip}
\begin{corollary}\label{b3:cor:01-discrete-zeros}
设 \(\Omega\subseteq\mathbb C\) 为区域，\(f:\Omega\to\mathbb C\) 全纯且不恒为零。则 \(f\) 的每个零点都孤立，并且在每个紧集 \(K\subseteq\Omega\) 内只有有限个零点。
\end{corollary}
\begin{proof}
若有零点不是有限阶，\cref{b3:prop:01-zero-factorization}给出恒零邻域，再由恒等定理得到全域恒零，矛盾。所以每个零点都孤立。若紧集内有无穷多个不同零点，紧致性给出其中一列的聚点；聚点仍在 \(K\subseteq\Omega\) 内，再次与恒等定理矛盾。
\end{proof}

聚点位于区域内部这一条件不能删去。函数 \(\sin(\pi/z)\) 在 \(\mathbb C\setminus\{0\}\) 上全纯，具有零点 \(1/n\)、\(n\geq1\)，却不恒为零；这些零点的聚点 \(0\) 不在定义域内。连通性也有自己的作用：在两个不交圆盘上分别取常数 \(0\) 和 \(1\)，得到一个全纯函数，它在一个圆盘内恒为零，却不能因此确定另一个圆盘上的值。

\ProseParagraphBreak

一个常用后果是：若两个全纯函数在区域内的一段非退化实区间上相等，则它们在整个区域相等。实区间在复平面中没有内部，但它拥有位于定义域内的聚点。实际判断时甚至不需要整段区间：若整函数 \(f\) 满足 \(f(1/n)=e^{1/n}\) 对所有正整数 \(n\) 成立，聚点 \(0\) 位于共同定义域内，便有 \(f(z)=e^z\) 于全平面。相比之下，只给 \(f(n)=e^n\) 并不能唯一确定 \(f\)，因为 \(e^z+\sin\pi z\) 也满足同样条件，而正整数在平面内没有聚点。

\subsection{局部逆函数定理}
\label{b3:subsec:010403}

复导数不为零时，微分是可逆的旋转伸缩。局部\term[ni-hanshu-dingli]{逆函数定理}[inverse function theorem]把这一点上的可逆性提升为邻域内函数的可逆性。本书沿用第一、二卷的“逆函数”主称。

\Needspace{6\baselineskip}
\begin{theorem}[局部逆函数]\label{b3:thm:01-local-inverse}
设 \(\Omega\subseteq\mathbb C\) 为开集，\(a\in\Omega\)，\(f:\Omega\to\mathbb C\) 全纯且 \(f'(a)\neq0\)。则存在开邻域 \(U\subseteq\Omega\)、\(V\subseteq\mathbb C\)，分别包含 \(a\)、\(f(a)\)，使 \(f:U\to V\) 为双射，\(f'\) 在 \(U\) 上不为零，其逆函数 \(h:V\to U\) 全纯，并满足
\[
h'(w)=\frac1{f'\bigl(h(w)\bigr)}\quad(w\in V).
\]
\end{theorem}
\begin{proof}
由\cref{b3:thm:01-holomorphic-analytic}，将 \(f\) 看作 \(\mathbb R^2\) 间的映射时，它为 \(C^1\)。其在 \(a\) 处的实微分为乘以 \(f'(a)\) 的复线性映射，Jacobi 行列式为 \(|f'(a)|^2>0\)。第一卷引理 18.1.2 给出的实逆函数定理于是提供邻域间的 \(C^1\) 逆映射。由 \(f'\) 的连续性，可在所得源邻域内取更小的开邻域 \(U\)，使 \(f'\) 处处非零。实逆映射连续，故 \(V=f(U)\) 是原目标邻域中的开集；限制映射后仍得到邻域间的双射及连续逆映射 \(h\)。

固定 \(w_0\in V\)，记 \(z_0=h(w_0)\)。当 \(w\neq w_0\) 且 \(w\to w_0\) 时，单射性与逆映射连续性分别给出 \(z=h(w)\ne z_0\) 及 \(z\to z_0\)，因此
\[
\frac{h(w)-h(w_0)}{w-w_0}
=\left(\frac{f(z)-f(z_0)}{z-z_0}\right)^{-1}
\longrightarrow\frac1{f'(z_0)}.
\]
所以逆映射在每个点复可微，得到结论。
\end{proof}
% 跨卷引用：Book1/Part03/Chapter18/Section1801.tex，b1:lem:level-set-local-inverse，当前引理18.1.2；完整实逆函数证明不在本卷重复。

例如 \(f(z)=z^2\) 在 \(a=1\) 的导数为 \(2\)。目标点 \(w=1\) 附近选择
\[
 h(w)=\exp\bigl(\tfrac12\operatorname{Log}w\bigr),
\]
便有 \(h(1)=1\)、\(h(w)^2=w\) 和 \(h'(1)=1/2\)。另一个解 \(-h(w)\) 靠近 \(-1\)，不属于这里选定的源邻域。所谓局部逆函数，既要指定目标点附近的范围，也要指定逆值靠近哪个原像。

\ProseParagraphBreak

指数函数处处导数非零，因此每一点附近都可求全纯逆函数；但 \(e^{z+2\pi\mathrm{i}}=e^z\)，它在整个复平面上并非单射。若固定 \(e^{z_0}=w_0\ne0\)，则
\[
 h(w)=z_0+\operatorname{Log}(w/w_0)
 \quad\bigl(|w-w_0|<|w_0|\bigr)
\]
给出满足 \(h(w_0)=z_0\) 的逆分支。这里 \(w/w_0\) 落在右半平面，主对数有定义，且 \(e^{h(w)}=w\)。把 \(z_0\) 换成 \(z_0+2\pi\mathrm{i}k\) 就得到其他分支；局部求逆不会消除全局的周期。

\subsection{开映射定理}
\label{b3:subsec:010404}

在导数为零的点，局部逆函数定理不再适用。零点分解仍给出
\(f(z)-f(a)=(z-a)^m\cdot g(z)\)，其中 \(g\) 在小邻域内无零。若能把这个非零因子写成 \(q^m\)，便可令 \(\phi=(z-a)q\)，把所有高阶修正吸收到一个新坐标中。需要求根的是接近非零常数的 \(g\)，因此前节的局部对数已经足够。

\Needspace{6\baselineskip}
\begin{proposition}[局部标准形]\label{b3:prop:01-local-normal-form}
设 \(\Omega\subseteq\mathbb C\) 为开集，\(a\in\Omega\)，\(f:\Omega\to\mathbb C\) 全纯，且 \(f-f(a)\) 不在 \(a\) 的任何邻域内恒为零。令 \(m\geq1\) 为 \(f-f(a)\) 在 \(a\) 的零点阶数。则存在开邻域 \(U\subseteq\Omega\)、\(a\in U\)，数 \(r>0\) 及全纯双射 \(\phi:U\to D(0,r)\)，其逆函数也全纯，且 \(\phi(a)=0\)、\(\phi'\) 在 \(U\) 上不为零，并有
\begin{equation}\label{b3:eq:01-local-normal-form}
f(z)=f(a)+\phi(z)^m\quad(z\in U).
\end{equation}
邻域 \(U\) 可以取得包含于任意预先给定的 \(a\) 的开邻域。
\end{proposition}
\begin{proof}
由零点分解，在一个足够小的圆盘 \(D(a,R)\subseteq\Omega\) 上有 \(f(z)-f(a)=(z-a)^m\cdot g(z)\)，其中 \(g\) 全纯且 \(g(a)\neq0\)。若预先给定了 \(a\) 的开邻域，开始时便把这个圆盘取在其中。取 \(c\in\mathbb C\) 满足 \(c^m=g(a)\)；由非零复数的极坐标表示可取得这样的 \(c\)，且 \(c\ne0\)。继续缩小 \(R\)，使
\[
 \left|\frac{g(z)}{g(a)}-1\right|<\frac12
 \quad\bigl(z\in D(a,R)\bigr).
\]
因此 \(g(z)/g(a)\) 始终落在主对数的定义域内。利用\cref{b3:prop:01-elementary-principal-log}，在这个圆盘上定义
\[
 q(z)=c\cdot\exp\left(\frac1m\operatorname{Log}\frac{g(z)}{g(a)}\right).
\]
指数和对数的恒等式给出 \(q(a)=c\)、\(q(z)^m=g(z)\)。于是全纯函数 \(\phi(z)=(z-a)\cdot q(z)\) 满足 \(\phi(a)=0\)、\(\phi'(a)=c\ne0\) 及所需的幂关系。

对 \(\phi\) 应用\cref{b3:thm:01-local-inverse}，取得 \(a\) 的开邻域 \(U_0\) 和 \(0\) 的开邻域 \(V_0\)，使 \(\phi:U_0\to V_0\) 及其逆函数 \(\psi\) 均全纯，且 \(\phi'\) 在 \(U_0\) 无零。选 \(r>0\)，使 \(D(0,r)\subseteq V_0\)，再令
\[
 U=\{z\in U_0:\phi(z)\in D(0,r)\}.
\]
连续性使 \(U\) 开，且它包含 \(a\)；限制后的 \(\phi\) 把 \(U\) 双射到整个 \(D(0,r)\)，逆函数为 \(\psi\) 在此圆盘上的限制。这样取得的源邻域 \(U\) 未必是几何圆盘，但它在 \(\phi\) 坐标中恰是一个圆盘，幂关系在其上成立。
\end{proof}

当 \(m=1\) 时，原函数本身已在邻域内可逆；当 \(m\geq2\) 时，新坐标中的函数为 \(\xi\mapsto f(a)+\xi^m\)。构造中开根只针对始终靠近非零值的局部因子，没有要求含零点的函数存在全纯根，也没有要求任意区域上的非零函数都有全纯对数。

\Needspace{6\baselineskip}
\begin{theorem}[开映射]\label{b3:thm:01-open-mapping}
设 \(\Omega\subseteq\mathbb C\) 为区域，\(f:\Omega\to\mathbb C\) 为非常值全纯函数。则 \(f\) 是\term[kai-yingshe]{开映射}[open map]：对任意开集 \(O\subseteq\Omega\)，像 \(f(O)\) 在 \(\mathbb C\) 中开。这个结论称为\term[kai-yingshe-dingli]{开映射定理}[open mapping theorem]。
\end{theorem}
\begin{proof}
空集的像为空集。若 \(O\ne\varnothing\)，任取 \(a\in O\)。恒等定理说明 \(f-f(a)\) 不会在 \(a\) 附近恒零，否则 \(f\) 将在整个 \(\Omega\) 上为常数。将\cref{b3:prop:01-local-normal-form}的邻域取在 \(O\) 内，得到 \(f(z)=f(a)+\phi(z)^m\)，且 \(\phi(U)=D(0,r)\)。幂映射把 \(D(0,r)\) 满射到 \(D(0,r^m)\)：任意非零数 \(w\) 的 \(m\) 次根均有模 \(|w|^{1/m}\)，零的原像为零。因此
\[
D\bigl(f(a),r^m\bigr)=f(U)\subseteq f(O).
\]
像中每个点都是内点，故 \(f(O)\) 开。
\end{proof}

\Needspace{6\baselineskip}
\begin{corollary}\label{b3:cor:01-local-injectivity}
设 \(\Omega\subseteq\mathbb C\) 为开集，\(a\in\Omega\)，\(f:\Omega\to\mathbb C\) 全纯。则 \(f\) 在 \(a\) 的某邻域内单射，当且仅当 \(f'(a)\ne0\)。更一般地，若 \(f-f(a)\) 在 \(a\) 具有有限零点阶数 \(m\ge1\)，\(U\subseteq\Omega\) 是 \(a\) 的开邻域，且全纯双射 \(\phi:U\to D(0,r)\)、\(r>0\)，满足 \(\phi(a)=0\)、\(f=f(a)+\phi^m\)，则每个满足 \(0<|w-f(a)|<r^m\) 的值 \(w\) 恰有 \(m\) 个位于 \(U\) 的不同原像。
\end{corollary}
\begin{proof}
先数标准形邻域内的原像。写 \(w-f(a)=\rho e^{\mathrm{i}\theta}\)，其中 \(0<\rho<r^m\)。方程 \(f(z)=w\) 等价于 \(\phi(z)^m=w-f(a)\)，其 \(m\) 个根为
\[
 \xi_j=\rho^{1/m}\exp\left(\frac{\mathrm{i}(\theta+2\pi j)}m\right),
 \quad j=0,\ldots,m-1.
\]
这些根两两不同，均落在 \(D(0,r)\)，且由模和辐角的表示不存在其他根。\(\phi\) 的双射性给出恰好 \(m\) 个不同原像。

若 \(f'(a)=0\) 而 \(f\) 不局部恒定，则其首个非零 Taylor 项的次数 \(m\ge2\)。对任意给定的 \(a\) 的开邻域，都可在其中取标准形邻域，再选择 \(0<|w-f(a)|<r^m\)，得到多个原像，所以该邻域上不能单射；局部恒定时也不能在任何 \(a\) 的邻域内单射。反方向由局部逆函数定理得到。
\end{proof}

函数 \(z\mapsto z^m\) 在 \(m\ge2\) 时提供最直接的例子：原点处导数为零，但每个小圆盘仍映为一个圆盘。开映射要求像含有邻域，局部单射要求不同点不重合。对实可微函数 \(z\mapsto|z|^2\)，小圆盘的像只落在实轴上，开映射结论便不成立。

\ProseParagraphBreak

上述原像计数还可以写成具体的局部逆分支。令 \(b=f(a)\)，\(\psi=\phi^{-1}:D(0,r)\to U\)，并在目标圆盘中删去由 \(b\) 向负实方向伸出的半径，得到
\[
 W=\{w:0<|w-b|<r^m,\ w-b\notin(-\infty,0]\}.
\]
对 \(j=0,\ldots,m-1\)，定义
\[
 h_j(w)=\psi\left(e^{2\pi\mathrm{i}j/m}\cdot
          \exp\bigl(\tfrac1m\operatorname{Log}(w-b)\bigr)\right),
 \quad w\in W.
\]
括号中数的模为 \(|w-b|^{1/m}<r\)，所以它在 \(\psi\) 的定义域内。各 \(h_j\) 全纯，满足 \(f(h_j(w))=w\)，并且对每个 \(w\in W\) 恰好给出 \(U\) 内全部原像。这里删去半径是为了在同一目标开集上连续标定这些根；对单个非零目标值计数时，不需要作任何分支选择。

\ProseParagraphBreak

对开集上的全纯函数，导数为零的点称为\term[linjie-dian]{临界点}[critical point]，其像称为\term[linjie-zhi]{临界值}[critical value]。在上述标准形中，\(\phi'\) 处处非零，而
\[
 f'(z)=m\phi(z)^{m-1}\phi'(z).
\]
故 \(m\ge2\) 时，\(a\) 是 \(U\) 内唯一的临界点，\(b\) 是这部分映射的唯一临界值；对 \(w\ne b\)，那些不同原像处的导数均不为零，每一个原像附近都有通常的局部逆。其他源邻域是否也映到这些值，需要另行研究，不能由这一局部计数决定。

\ProseParagraphBreak

具体来看 \(f(z)=z^2+z^3=z^2(1+z)\)。在 \(|z|<1/2\) 上，\(1+z\) 落在右半平面，可以定义
\[
 \phi(z)=z\exp\bigl(\tfrac12\operatorname{Log}(1+z)\bigr).
\]
直接求值得 \(\phi(0)=0\)、\(\phi'(0)=1\)，并且 \(f=\phi^2\)。按局部标准形的做法，进一步限制到 \(\phi:U\to D(0,r)\) 可逆的邻域，记逆函数为 \(\psi\)。在 \(0<|w|<r^2\) 且删去负实半径的目标区域中，两条逆分支正是
\[
 z_+(w)=\psi\bigl(\sqrt w\bigr),
 \quad z_-(w)=\psi\bigl(-\sqrt w\bigr),
 \quad \sqrt w=\exp\bigl(\tfrac12\operatorname{Log}w\bigr).
\]
三次项通过坐标 \(\phi\) 保留下来。它会改变根的位置，却不改变足够小的 \(U\) 内恰有两个不同原像这一事实。

\ProseParagraphBreak

还可以算出位置修正。由 Taylor 展开，\(\phi(z)=z+z^2/2+O(|z|^3)\)。因 \(\psi'(0)=1\)，写 \(\psi(\xi)=\xi+A\xi^2+O(|\xi|^3)\)，代入 \(\phi(\psi(\xi))=\xi\)，比较二次项得 \(A=-1/2\)。因此
\[
\begin{aligned}
 z_+(w)&=\sqrt w-\tfrac12w+O(|w|^{3/2}),\\
 z_-(w)&=-\sqrt w-\tfrac12w+O(|w|^{3/2}).
\end{aligned}
\]
余项估计来自 \(\psi\) 在一个更小闭圆盘上的 Taylor 余项，所以在所取割圆盘中对所有趋近方向同时有效。若沿 \(w(t)=\rho e^{\mathrm{i}t}\)、\(0\le t\le2\pi\)，连续选根从 \(\sqrt\rho\) 出发，便得到 \(\sqrt\rho e^{\mathrm{i}t/2}\)：任意连续选根与这个表达式的比只能取 \(1\) 或 \(-1\)，连续性使该比保持常数。终点为 \(-\sqrt\rho\)，两条局部反解交换，因而不能各自成为整个去心目标圆盘上的连续单值逆函数。

\ProseParagraphBreak

本例中 \(f'(z)=z(2+3z)\)，两个临界点为 \(0\)、\(-2/3\)，对应临界值为 \(0\)、\(4/27\)。在第二个点令 \(z=-2/3+\zeta\)，则
\[
 f(-2/3+\zeta)=\frac4{27}-\zeta^2+\zeta^3,
\]
所以那里同样是二次的局部模型。另一方面，\(f(-1)=0\)、\(f'(-1)=1\)，局部逆函数定理还给出靠近 \(-1\) 的一个分支 \(z_0(w)=-1+w+O(|w|^2)\)。它与靠近原点的两个分支位于不交源邻域。由此，充分小的非零 \(w\) 在全平面有三个不同原像，而在原点的标准形邻域内只有两个；次数为三的多项式至多有三个不同根，因而这里已经列尽全部原像。局部重数记录的是原像在哪个点聚合，不是整个定义域中的总原像数。

\ProseParagraphBreak

\chapref{b3:ch:02}将讨论一般曲线上的复积分，以及 Cauchy 积分公式给出的估计与全局结论。开映射也将用来证明内部最大模原理。本章得到的局部幂级数、零点分解和标准形，随后还会用于奇点分类、零点计数与解析延拓。

% 本章小结（不计入编号）
\begin{chaptersummary}
复可微把二维实微分限制为复数乘法。偏导数满足 Cauchy--Riemann 方程只是其中一部分；若没有实可微的统一余项控制，复差商仍可能不存在。反过来，函数在开集上处处复可微，则由局部积分公式得到收敛的 Taylor 展开，各阶导数的存在与连续性随之成立。

\ProseParagraphBreak

幂级数逐项运算的依据是收敛圆内部的统一估计。选取两个有间隔的半径，可用几何衰减同时控制尾项和各阶导数；把内半径推到外半径时，原有估计可能失效，因此边界需要另行讨论。指数与三角函数可在全平面构造，对数却需要选定分支。局部存在全纯的对数或逆函数，并不保证能在整个定义域连续选取；指数函数的周期和去心平面上不存在连续单值对数的事实说明了这一区别。

\ProseParagraphBreak

非恒零全纯函数的零点由最先出现的非零 Taylor 项控制。恒等定理通过连通性把局部相等传到整个区域；局部标准形进一步把非常值函数化成一个幂映射。标准坐标中，导数非零对应一次幂，较高阶零点对应多次幂；附近非中心值的原像数目便由这个次数确定。这里计数的是选定邻域内的原像，其他位置仍可能另有逆分支。两种情形都把邻域映为邻域。下一章将在本章局部积分结果上继续建立一般 Cauchy 理论。
\end{chaptersummary}

\BookExercises

\Cref{b3:ex:01-reconstruct-real-part}与\cref{b3:ex:01-power-series-ode}训练方程和幂级数的具体计算；\cref{b3:ex:01-mean-value-failure}检验实变量中值定理的适用边界。\Cref{b3:ex:01-disk-distance-ratio}考察圆盘变换的展开与参数极限，\cref{b3:ex:01-equal-modulus}与\cref{b3:ex:01-local-involution}组合使用零点、恒等定理和局部求逆。两道带星号的题进一步讨论逆函数的展开半径和幂级数跨越收敛圆周的障碍，所需工具均已在本章建立。

% 自拟题组；保留既有训练方向，解答使用本章正文工具，不复刻第三方原题。
% 实质增量：实部重建、中值等式反例、整函数常微分方程、圆盘变换的参数极限、
% 等模函数共同零点、局部对合线性化、逆函数半径与阶乘缺项的边界障碍。

\Needspace{6\baselineskip}
\begin{exercise}\label{b3:ex:01-reconstruct-real-part}
求出所有满足
\[
  \operatorname{Re}f(x+\mathrm{i}y)=x^3-3xy^2+2xy
  \quad(x,y\in\mathbb R)
\]
的整函数 \(f:\mathbb C\rightarrow\mathbb C\)，并说明给定条件为什么只留下一个实常数的自由度。
\end{exercise}
\begin{solution}
记题给的实部为 \(u\)，虚部为 \(v\)。全纯性保证实部、虚部的偏导数存在并满足 Cauchy--Riemann 方程，因而
\[
  v_y=u_x=3x^2-3y^2+2y,
  \quad v_x=-u_y=6xy-2x.
\]
固定 \(x\)，对第一个等式关于 \(y\) 积分，得到
\[
  v(x,y)=3x^2y-y^3+y^2+h(x),
  \quad h(x)=v(x,0).
\]
这里的积分常数可以依赖 \(x\)。由于 \(h\) 可微，再代入第二个等式，有
\(6xy+h'(x)=6xy-2x\)，所以 \(h(x)=-x^2+C\)，其中 \(C\in\mathbb R\)。因此
\[
  f(z)=z^3-\mathrm{i}z^2+\mathrm{i}C,
  \quad C\in\mathbb R.
\]
反向核对时，\(\operatorname{Re}(z^3)=x^3-3xy^2\)、\(\operatorname{Re}(-\mathrm{i}z^2)=2xy\)，所以这些整函数都符合题设。计算已从必要条件确定了全部可能的虚部，故没有遗漏。固定实部后只能给虚部加一个实常数；若给 \(f\) 加上实部非零的常数，就会改变题给实部。
\end{solution}

\Needspace{6\baselineskip}
\begin{exercise}\label{b3:ex:01-mean-value-failure}
设 \(U\subseteq\mathbb C\) 为开集，\(f:U\rightarrow\mathbb C\) 全纯，\(a,b\in U\)、\(a\ne b\)，且整条闭线段 \([a,b]\) 包含于 \(U\)。证明
\[
  f(b)-f(a)=(b-a)\int_0^1 f'\bigl(a+t(b-a)\bigr)\,\mathrm{d}t
\]
及相应的模估计。再构造一个整函数，说明一般不存在 \(t\in(0,1)\) 使
\[
  f(b)-f(a)=(b-a)\cdot f'\bigl(a+t(b-a)\bigr).
\]
\end{exercise}
\begin{solution}
由全纯与解析的等价性，\(f'\) 连续。函数
\(F(t)=f\bigl(a+t(b-a)\bigr)\) 在 \([0,1]\) 上是复值 \(C^1\) 函数，并满足
\[
  F'(t)=(b-a)\cdot f'\bigl(a+t(b-a)\bigr).
\]
将实变量微积分基本定理分别用于 \(F\) 的实部和虚部，合并后便得到题给的积分等式。三角不等式进一步给出
\[
\begin{aligned}
  |f(b)-f(a)|
  &\le |b-a|\cdot
       \int_0^1\left|f'\bigl(a+t(b-a)\bigr)\right|\,\mathrm{d}t\\
  &\le |b-a|\cdot
       \max_{0\le t\le1}\left|f'\bigl(a+t(b-a)\bigr)\right|.
\end{aligned}
\]
最后一个最大值存在，因为被取模的函数在紧区间上连续。

取 \(f(z)=e^z\)、\(a=0\)、\(b=2\pi\mathrm{i}\)。这时 \(f(b)-f(a)=0\)，而
\[
  (b-a)\cdot f'\bigl(a+t(b-a)\bigr)
  =2\pi\mathrm{i}e^{2\pi\mathrm{i}t}\ne0
\]
对所有实数 \(t\) 都成立。因此，实变量中值定理中的单点等式不能照搬。复导数沿线段的方向可以相互抵消，使积分为零，而没有任何一个导数值等于零。
\end{solution}

\Needspace{6\baselineskip}
\begin{exercise}\label{b3:ex:01-power-series-ode}
求所有满足
\[
  f'(z)=z\cdot f(z)\quad(z\in\mathbb C),
  \quad f(0)=1
\]
的整函数 \(f:\mathbb C\rightarrow\mathbb C\)。先由幂级数求出系数并核对收敛半径，再证明解的唯一性。
\end{exercise}
\begin{solution}
任何可能的解在原点附近都具有展开 \(f(z)=\sum_{n\ge0}a_nz^n\)。逐项求导后比较同次幂系数，得到
\[
  a_0=1,\quad a_1=0,
  \quad (n+1)a_{n+1}=a_{n-1}\quad(n\ge1).
\]
这个递推把奇数次与偶数次系数分开。所有奇数次系数均为零，而
\[
  a_{2k}=\frac{a_{2k-2}}{2k}
         =\frac1{2^k k!}\quad(k\ge1).
\]
因此唯一的候选级数为
\[
  f(z)=\sum_{k=0}^{\infty}\frac{z^{2k}}{2^k k!}
       =e^{z^2/2}.
\]
对任意固定的 \(r>0\)，在 \(|z|\le r\) 上，绝对值级数由
\(\sum_{k\ge0}r^{2k}/(2^k k!)\) 控制。其相邻项比值为
\(r^2/(2(k+1))\)，趋于零，故控制级数收敛。所求幂级数在每个闭圆盘上一致绝对收敛，半径为无穷大；逐项求导和上述递推证明其确实满足微分方程与初值。

若 \(g\) 是另一整函数解，则同样的递推使它与 \(f\) 在原点附近的全部 Taylor 系数相同。因此它们在一个开圆盘内相等，再由恒等定理在连通的 \(\mathbb C\) 上相等。也可以直接计算积分因子：
\[
  \bigl(e^{-z^2/2}\cdot g(z)\bigr)'
  =e^{-z^2/2}\cdot\bigl(g'(z)-z\cdot g(z)\bigr)=0.
\]
其值于是局部恒定，再由连通性为同一个常数；原点处的值为 \(1\)，再次得到 \(g(z)=e^{z^2/2}\)。
\end{solution}

\Needspace{6\baselineskip}
\begin{exercise}\label{b3:ex:01-disk-distance-ratio}
设 \(a\in\mathbb C\)、\(0<|a|<1\)，在 \(\mathbb C\setminus\{1/\overline a\}\) 上定义
\[
  \phi_a(z)=\frac{z-a}{1-\overline a z}.
\]
求它在原点的 Taylor 级数、精确收敛半径，以及截到 \(z^N\) 后在 \(|z|\le r\) 上的余项界，其中 \(N\ge0\) 为整数，\(0<r<1/|a|\)。再固定 \(\zeta\in\mathbb C\)、\(|\zeta|=1\)，令 \(a=t\zeta\)，其中 \(0<t<1\)。证明当 \(t\uparrow1\) 时，\(\phi_{t\zeta}\) 在每个固定的闭圆盘 \(\overline{D(0,r)}\)、\(0<r<1\) 上一致趋于常数 \(-\zeta\)，却不在闭单位圆盘上一致趋于该常数。
\end{exercise}
\begin{solution}
先把常数项分离：
\[
  \phi_a(z)=-a+(1-|a|^2)\frac{z}{1-\overline a z}.
\]
在 \(|z|<1/|a|\) 上应用几何级数，得
\[
  \phi_a(z)=-a+(1-|a|^2)
       \sum_{n=1}^{\infty}\overline a^{\,n-1}z^n.
\]
对 \(n\ge1\)，系数模的 \(n\) 次根为
\[
  (1-|a|^2)^{1/n}|a|^{(n-1)/n}\longrightarrow|a|.
\]
所以 Cauchy--Hadamard 公式给出的半径恰为 \(1/|a|\)。由于 \(a\ne0\) 且 \(|a|<1\)，这些高次系数并不消失。

设 \(P_N\) 为保留到 \(z^N\) 的多项式，其中 \(P_0=-a\)。对余下的几何级数求和，得到
\[
  \phi_a(z)-P_N(z)
  =\frac{(1-|a|^2)\overline a^{\,N}z^{N+1}}
         {1-\overline a z}.
\]
在 \(|z|\le r<1/|a|\) 上有
\(|1-\overline a z|\ge1-|a|\cdot|z|\ge1-|a|r>0\)，因而
\[
  \sup_{|z|\le r}|\phi_a(z)-P_N(z)|
  \le\frac{(1-|a|^2)|a|^N r^{N+1}}{1-|a|r}.
\]
这个界对 \(N=0\) 同样成立。

最后直接计算
\[
  \phi_{t\zeta}(z)+\zeta
  =\frac{(1-t)(z+\zeta)}{1-t\overline\zeta z}.
\]
在固定的 \(|z|\le r<1\) 上，其模不超过
\((1-t)(1+r)/(1-tr)\)。分子趋于零，分母至少为 \(1-r>0\)，所以得到所需一致收敛。对于每个 \(t<1\)，分母在闭单位圆盘上也不为零，故边界点值均有定义；但
\[
  \phi_{t\zeta}(\zeta)=\frac{(1-t)\zeta}{1-t}=\zeta,
  \quad |\phi_{t\zeta}(\zeta)+\zeta|=2.
\]
于是闭单位圆盘上的一致收敛失败。固定内圆盘的正边界间隔正是前一估计能够统一控制分母的原因。
\end{solution}

\Needspace{6\baselineskip}
\begin{exercise}\label{b3:ex:01-equal-modulus}
设 \(U\subseteq\mathbb C\) 为非空连通开集，\(f,g:U\rightarrow\mathbb C\) 全纯，且
\(|f(z)|=|g(z)|\) 对每个 \(z\in U\) 成立。证明存在常数 \(c\in\mathbb C\)、\(|c|=1\)，使 \(f=cg\) 于 \(U\)。解答须处理共同零点，不能直接假设 \(f/g\) 处处有定义。
\end{exercise}
\begin{solution}
若 \(g\equiv0\)，等模条件使 \(f\equiv0\)，取 \(c=1\) 即可。以下设 \(g\not\equiv0\)，因而 \(f\not\equiv0\)。它们的零点集合相同，且由恒等定理和零点分解，每个共同零点都有有限阶。

任取一个共同零点 \(a\)。在它的某个圆盘邻域内写
\[
  f(z)=(z-a)^m\cdot F(z),
  \quad g(z)=(z-a)^n\cdot G(z),
  \quad F(a)\ne0,\quad G(a)\ne0.
\]
这里 \(m,n\) 为正整数，\(F,G\) 全纯。缩小该圆盘，使 \(F,G\) 均无零点。对其中的 \(z\ne a\)，等模条件化为
\[
  |z-a|^{m-n}\cdot\left|\frac{F(z)}{G(z)}\right|=1.
\]
当 \(z\to a\) 时，商的模趋于有限正数 \(|F(a)/G(a)|\)。如果 \(m>n\)，左边趋于零；如果 \(m<n\)，左边趋于无穷大。因此只能有 \(m=n\)。这时在去心圆盘上 \(f/g=F/G\)，而 \(F/G\) 已在整个圆盘上全纯，故可用其在 \(a\) 的值补定义商。

对每个共同零点作这样的补定义，并在 \(g\ne0\) 处仍取通常的商，得到 \(U\) 上的函数 \(h\)。它在每点附近都等于一个全纯函数，所以全纯；各局部定义在交集上相容，因为它们在非零点处都等于 \(f/g\)，在零点处的值再由连续性唯一确定。由非零点的等模条件以及零点处的连续性，\(|h|=1\) 在整个 \(U\) 上成立。

若 \(h\) 非常值，开映射定理将使 \(h(U)\) 成为非空开集；但 \(h(U)\) 包含于单位圆，后者不含任何开圆盘。这是矛盾。因此 \(h=c\)，其中 \(|c|=1\)。在非零点处有 \(f=cg\)，共同零点处两边均为零，故所求等式在全域成立。
\end{solution}

\Needspace{6\baselineskip}
\begin{exercise}\label{b3:ex:01-local-involution}
设 \(U\subseteq\mathbb C\) 是包含原点的区域，\(f:U\rightarrow U\) 全纯，并满足
\[
  f(0)=0,\quad f\bigl(f(z)\bigr)=z\quad(z\in U).
\]
证明 \(f'(0)\) 只能为 \(1\) 或 \(-1\)。若为 \(1\)，证明 \(f(z)=z\) 于 \(U\)；若为 \(-1\)，构造一个在原点附近具有全纯逆的函数 \(h\)，使
\(h\bigl(f(z)\bigr)=-h(z)\)，并说明如何选取原点邻域，使这一等式确实把 \(f\) 化成新坐标中的 \(w\mapsto-w\)。
\end{exercise}
\begin{solution}
在原点对复合恒等式使用链式法则，得到
\(f'(f(0))\cdot f'(0)=1\)，即 \(f'(0)^2=1\)。

先设 \(f'(0)=1\)。若 \(f\) 不恒等于恒等映射，则恒等定理排除了 \(f(z)-z\) 在原点附近恒零的可能。由 Taylor 展开，存在最小的整数 \(m\ge2\) 及 \(a\ne0\)，使
\[
  f(z)=z+az^m+R(z),\quad R(z)=O(|z|^{m+1}).
\]
为了核对复合的首个非零项，先注意 \(f(z)=O(|z|)\)，故
\(R(f(z))=O(|z|^{m+1})\)。又由
\(f(z)=z\bigl(1+O(|z|^{m-1})\bigr)\) 得
\[
  f(z)^m=z^m+O(|z|^{2m-1})
         =z^m+O(|z|^{m+1}),
\]
最后一步用了 \(m\ge2\)。因此
\[
\begin{aligned}
  f\bigl(f(z)\bigr)
  &=f(z)+a\cdot f(z)^m+R(f(z))\\
  &=z+2az^m+O(|z|^{m+1}).
\end{aligned}
\]
它不可能恒等于 \(z\)：否则减去 \(z\) 后除以 \(z^m\)，令 \(z\to0\)，便有 \(2a=0\)。这个矛盾证明 \(f(z)=z\) 于整个 \(U\)。

再设 \(f'(0)=-1\)。在 \(U\) 上令
\[
  h(z)=\frac{z-f(z)}2.
\]
则 \(h(0)=0\)、\(h'(0)=1\)，且对所有 \(z\in U\)，
\[
  h\bigl(f(z)\bigr)
  =\frac{f(z)-f(f(z))}2
  =-h(z).
\]
由局部逆函数定理，可取原点邻域 \(V_0\subseteq U\)，使 \(h\) 在 \(V_0\) 上为具有全纯逆的双射。还须保证所选邻域在 \(f\) 下保持不变。题设 \(f\circ f=\operatorname{id}\) 说明 \(f\) 是自己的连续逆，故
\[
  V_1=V_0\cap f(V_0)
\]
是包含原点的开集，并满足 \(f(V_1)=V_1\)。限制后的 \(h:V_1\rightarrow h(V_1)\) 仍为全纯双射，像是原点的开邻域。取 \(r>0\) 使 \(D(0,r)\subseteq h(V_1)\)，并令
\[
  V=(h|_{V_1})^{-1}(D(0,r)).
\]
若 \(z\in V\)，则 \(f(z)\in V_1\)，并由 \(h(f(z))=-h(z)\in D(0,r)\) 得 \(f(z)\in V\)。再次应用 \(f\) 可知 \(f(V)=V\)。因此 \(h:V\rightarrow D(0,r)\) 确实给出同一邻域内的坐标变换，并且
\[
  (h\circ f\circ h^{-1})(w)=-w
  \quad(|w|<r).
\]
\end{solution}

\Needspace{6\baselineskip}
\begin{optionalexercise}\label{b3:ex:01-inverse-radius}
设 \(F:\mathbb C\rightarrow\mathbb C\)、\(F(z)=z+z^2\)。求满足 \(G(0)=0\) 的全纯局部逆 \(G\) 在原点的前四个非零 Taylor 系数，并证明该 Taylor 级数的收敛半径恰为 \(1/4\)。这里 \(G\) 的初始定义域只要求是原点的一个开邻域。
\end{optionalexercise}
\begin{solution}
因为 \(F'(0)=1\)，局部逆函数定理保证所需 \(G\) 存在；限制到充分小的邻域后，它也由 \(G(0)=0\) 唯一确定。写
\[
  G(w)=b_1w+b_2w^2+b_3w^3+b_4w^4+O(|w|^5).
\]
代入 \(G+G^2=w\)，比较前四次幂的系数，依次得到
\[
\begin{aligned}
  b_1&=1,&b_2+b_1^2&=0,\\
  b_3+2b_1b_2&=0,&b_4+2b_1b_3+b_2^2&=0.
\end{aligned}
\]
因此
\[
  G(w)=w-w^2+2w^3-5w^4+O(|w|^5).
\]

接着证明收敛半径的下界。若 \(|w|<1/4\)，则
\(\operatorname{Re}(1+4w)>0\)，主对数有定义。函数
\[
  H(w)=\frac{\exp\bigl(\tfrac12\operatorname{Log}(1+4w)\bigr)-1}{2}
\]
在整个 \(D(0,1/4)\) 上全纯。令分子中的指数值为 \(S(w)\)，则
\(S(w)^2=1+4w\)、\(S(0)=1\)，从而 \(H(0)=0\) 且
\[
  H(w)+H(w)^2=\frac{S(w)^2-1}{4}=w.
\]
在充分小的原点邻域内，\(H(w)\) 落在 \(F\) 的局部单射邻域，故 \(F(H(w))=w\) 迫使 \(H=G\)。于是 \(H\) 的原点 Taylor 系数就是所求局部逆的系数。由圆盘上的全纯函数具有覆盖整个圆盘的 Taylor 展开，所求级数的收敛半径 \(R\ge1/4\)。

若 \(R>1/4\)，记该 Taylor 级数在 \(D(0,R)\) 上的和为 \(\widetilde G\)。它在原点附近与 \(G\) 相同，故
\(\widetilde G+\widetilde G^2-w\) 在那里为零。这一表达式在连通圆盘 \(D(0,R)\) 上全纯，恒等定理便使
\[
  \widetilde G(w)+\widetilde G(w)^2=w
  \quad(|w|<R)
\]
在全盘成立，而不只在最初的逆函数邻域成立。特别地，\(w_0=-1/4\) 属于此圆盘，代入得到
\((\widetilde G(w_0)+1/2)^2=0\)，所以 \(\widetilde G(w_0)=-1/2\)。但对全盘恒等式求导有
\[
  \bigl(1+2\widetilde G(w)\bigr)\cdot\widetilde G'(w)=1.
\]
在 \(w=w_0\) 处左边为零，矛盾。因此 \(R=1/4\)。这里出现的障碍来自 \(F'(-1/2)=0\)，而 \(F(-1/2)=-1/4\) 正好位于所求收敛圆周上。
\end{solution}

\Needspace{6\baselineskip}
\begin{optionalexercise}\label{b3:ex:01-factorial-gaps}
在单位圆盘 \(D(0,1)\) 上考虑
\[
  f(z)=\sum_{n=1}^{\infty}z^{n!}.
\]
证明这确实给出一个全纯函数，且其幂级数的收敛半径为 \(1\)。再证明：对每个 \(\zeta\in\mathbb C\)、\(|\zeta|=1\)，都不存在 \(\varepsilon>0\) 及 \(V=D(\zeta,\varepsilon)\) 上的全纯函数 \(F:V\rightarrow\mathbb C\)，使 \(F=f\) 在 \(V\cap D(0,1)\) 上成立。
\end{optionalexercise}
\begin{solution}
把缺少的幂次系数补成零，可写成 \(f(z)=\sum_{k=0}^{\infty}c_kz^k\)。由于 \(1!,2!,3!,\ldots\) 严格递增，系数 \(c_k\) 在 \(k=n!\) 时为 \(1\)，其余为 \(0\)。所以
\[
  \limsup_{k\to\infty}|c_k|^{1/k}=1.
\]
Cauchy--Hadamard 公式给出半径 \(1\)，并保证在每个闭子圆盘上一致绝对收敛；由幂级数逐项求导定理，其和在单位圆盘上全纯。

先考虑一个单位根 \(\eta\)，即存在正整数 \(q\) 使 \(\eta^q=1\)。取整数 \(N\ge q\)。当 \(n\ge N\) 时，\(q\) 整除 \(n!\)，所以 \(\eta^{n!}=1\)。对 \(0<r<1\)，有
\[
  f(r\eta)
  =\sum_{n<N}(r\eta)^{n!}+T(r),
  \quad T(r)=\sum_{n\ge N}r^{n!}.
\]
前一有限和的模不超过 \(N-1\)。后一和的每项都是非负实数；任意固定正整数 \(L\) 后，
\[
  T(r)\ge\sum_{n=N}^{N+L-1}r^{n!}
       \longrightarrow L\quad(r\uparrow1).
\]
因此 \(\liminf_{r\uparrow1}T(r)\ge L\) 对每个 \(L\) 成立，故 \(T(r)\to+\infty\)。于是
\[
  \operatorname{Re}f(r\eta)\ge T(r)-(N-1)\longrightarrow+\infty.
\]
这说明 \(f\) 沿趋向任意单位根的半径都没有有限极限，因而不可能等于该单位根邻域内某个连续函数的限制，更不可能全纯延拓到那里。

单位根 \(e^{2\pi\mathrm{i}p/q}\) 在单位圆上稠密：这是有理数在实轴上的稠密性与映射 \(t\mapsto e^{2\pi\mathrm{i}t}\) 的连续性给出的。若题设的 \(V,F\) 对某个 \(\zeta\) 存在，可以在开圆盘 \(V\) 中选取一个单位根 \(\eta\)。当 \(r<1\) 足够接近 \(1\) 时，\(r\eta\in V\cap D(0,1)\)，故
\[
  f(r\eta)=F(r\eta)\longrightarrow F(\eta).
\]
右边有限，与刚得到的实部趋于无穷大矛盾。由此排除了跨越单位圆上任意一点的所述局部延拓。
\end{solution}

% !TeX root = ../../Book3.tex
\chapter{Cauchy 理论}
\label{b3:ch:02}
\BookChapterNumber{b3:ch:02}
\begin{chapterabstract}
曲线积分记住行走的方向与次数，原函数则使它只依赖端点。本章从参数计算出发，把小圆盘上的原函数沿连续路径连接起来，再用同伦方格的边界相消建立单连通区域上的 Cauchy 定理。对一般区域边界，明确使用实变量 Green 公式，并通过挖去小圆盘恢复内部点值。导数公式与 Cauchy 估计继而给出整函数增长、零点存在、最大模和圆盘映射的刚性。
\end{chapterabstract}

\begin{learninggoals}
\begin{itemize}
\item\label{b3:item:02-goal-1}能从参数化计算曲线积分，区分曲线与像集，并在误差估计中同时控制函数大小和曲线长度。
\item\label{b3:item:02-goal-2}能解释连续路径积分的良定义和同伦不变性，用闭路积分判断原函数是否存在，并使用 Morera 定理完成粘合。
\item\label{b3:item:02-goal-3}能核对积分公式的区域、边界及方向条件，保留孔洞贡献，说明导数为何可以移入积分。
\item\label{b3:item:02-goal-4}能选择 Cauchy 估计中的圆心与半径，控制导数和截断误差，并利用增长、最大模和 Schwarz 引理判断刚性。
\end{itemize}
\end{learninggoals}

沿单位圆积分 \(z^2\) 得到零，沿同一圆积分 \(1/z\) 却得到 \(2\pi\mathrm{i}\)。两者在圆周附近都全纯，但第二个函数无法穿过中心全纯延拓。第一章的圆盘理论已经说明局部行为；现在的问题是，沿许多局部圆盘积累起来的端点差，何时能够在闭路上相消，又何时会留下绕过孔洞的贡献。

\ProseParagraphBreak

本章使用第一章的三角形 Goursat 定理、圆盘 Cauchy 公式和全纯函数的解析性，以及实变量积分和紧性的有限覆盖。前两节从局部原函数建立同伦积分与路径独立性；第三节的一般 Jordan 边界公式另使用多元微积分中的平面 Green 公式，并在使用处写清适用版本。它不参与第一章全纯即解析的证明。单连通性的收缩定义在本章给出，更系统的同伦与基本群理论留到\chapref{b3:ch:09}。

\ProseParagraphBreak

John B. Conway 的《Functions of One Complex Variable I》\cite{Conway1978Functions}可用于对照 Cauchy 理论与幂级数的章节组织。已有来源核验只覆盖其公开样章的幂级数部分，本章不附未经实读核实的积分定理页码。

% !TeX root = ../../Book3.tex
\section{复曲线积分}
\label{b3:sec:0201}

圆周上的积分已在\subsecref{b3:subsec:010204}用于恢复函数。一般曲线可以绕过孔洞，也可以反复经过同一点；积分必须记住行走方向与次数。下面先把参数曲线、积分估计和端点公式分开处理。

\subsection{分段光滑曲线}
\label{b3:subsec:020101}

设 \(\gamma:[a,b]\rightarrow\mathbb C\) 连续。若存在有限分点，使它在每个闭子区间上为 \(C^1\)，就称它为\term[fenduan-guanghua-quxian]{分段光滑曲线}[piecewise smooth curve]。允许某些区间上的导数为零，不要求曲线单射。曲线的长度定义为
\[
 L(\gamma)=\int_a^b|\gamma'(t)|\,\mathrm{d}t.
\]
反向曲线 \(\gamma^-(t)=\gamma(a+b-t)\) 的长度相同；首尾相接的曲线可以拼接，其长度相加。后文的闭曲线包括自交曲线，不能将“闭曲线”与“简单闭曲线”混为一谈。

\ProseParagraphBreak

圆周 \(a+r e^{\mathrm{i}t}\)、\(0\le t\le2\pi\)，规定正向为逆时针。把参数变为 \(2t\) 会把圆周走两遍；它与原曲线有相同的像，却有不同的积分。对非单射曲线，像集不能代替参数化。

\subsection{围道积分}
\label{b3:subsec:020102}

设 \(f\) 在曲线像上连续，沿 \(\gamma\) 的积分为
\[
 \int_\gamma f(z)\,\mathrm{d}z=\int_a^b f\bigl(\gamma(t)\bigr)\cdot\gamma'(t)\,\mathrm{d}t.
\]
右端按实部与虚部分别积分，有限个分点不影响数值。保持方向的分段 \(C^1\) 换参由实变量换元公式保持积分；反向换参改变符号。闭曲线上的积分也称\term[weidao-jifen]{围道积分}[contour integral]。
\begin{proposition}\label{b3:prop:02-ml}
设 \(f\) 在分段光滑曲线 \(\gamma\) 的像上连续。则
\[
 \left|\int_\gamma f(z)\,\mathrm{d}z\right|\le L(\gamma)\cdot\max_{z\in\gamma}|f(z)|.
\]
积分对被积函数线性，对曲线拼接可加，且 \(\int_{\gamma^-}f\,\mathrm{d}z=-\int_\gamma f\,\mathrm{d}z\)。
\end{proposition}
\begin{proof}
由实积分的三角不等式，积分模不超过 \(\int_a^b|f(\gamma(t))|\cdot|\gamma'(t)|\,\mathrm{d}t\)。取出连续函数在紧曲线像上的最大值即得估计。线性和拼接从有限区间积分的相应性质得到；反向时令 \(s=a+b-t\)，导数与积分上下限的变化给出负号。
\end{proof}
在半径 \(r\) 的正向圆周上，直接参数计算给出
\[
 \int_{|z-a|=r}(z-a)^n\,\mathrm{d}z=
 \begin{cases}2\pi\mathrm{i},&n=-1,\\0,&n\in\mathbb Z\setminus\{-1\}.\end{cases}
\]
这项计算不要求负幂在圆盘中心有定义。它将成为 Laurent 系数提取的依据。

\ProseParagraphBreak

复积分并不是把实积分中的变量符号换成 \(z\)。对任意分段光滑曲线，\(\int_\gamma\overline z\,\mathrm{d}z\) 一般不能写成端点的某个复函数之差。若 \(\gamma\) 是正向简单闭曲线，由 \(z=x+\mathrm{i}y\) 展开，
\[
 \overline z\,\mathrm{d}z=x\,\mathrm{d}x+y\,\mathrm{d}y+\mathrm{i}(x\,\mathrm{d}y-y\,\mathrm{d}x).
\]
前两项是 \(\tfrac12\mathrm{d}(x^2+y^2)\)，闭路上为零；后两项由平面 Green 公式等于内部面积的 \(2\mathrm{i}\) 倍。若暂不使用 Green 公式，对圆周直接计算已可看到面积贡献。被积函数不全纯时，内部积分无法用 Cauchy 消边法消去。

\ProseParagraphBreak

长度估计在曲线族随参数变化时也有用，但必须同时控制长度。例如圆 \(|z|=r\) 上的 \(z^{-2}\) 模为 \(r^{-2}\)，长度为 \(2\pi r\)，粗估计随 \(r\downarrow0\) 增大；实际积分却恒为零。这表明长度估计给上界，不自动反映不同相位的抵消。反之，若一段长为 \(\ell_n\) 的曲线上误差上界为 \(\varepsilon_n\)，仅 \(\varepsilon_n\rightarrow0\) 也未必使积分误差趋于0，还需 \(\ell_n\varepsilon_n\rightarrow0\)。

计算路径积分时，先检查被积式有没有单值原函数，常比直接代入参数更有效。例如沿任意从 \(0\) 到 \(1+\mathrm{i}\) 的分段光滑曲线，\(z e^{z^2}\) 的积分均为 \((e^{2\mathrm{i}}-1)/2\)。路径形状可以很复杂，但原函数 \(e^{z^2}/2\) 已把形状信息消去。相反，\(\overline z\) 的积分不能这样处理。在以 \(0,1,1+\mathrm{i},\mathrm{i}\) 为顶点的正向方形上，逐边参数化给积分 \(2\mathrm{i}\)。这个数来自区域的面积，说明仅有连续性与光滑性都不足以保证闭路积分为零。

\ProseParagraphBreak
参数方向也必须与被积式分开判断。把同一曲线反向走一遍，积分变号；把参数换成一个严格递增的分段光滑参数，积分不变。若参数曲线将某段往返走过，去程与回程相消。以后把区域切开时出现的两条切口边，正是利用这条相消规则。边界上的点值无需处处有导数；分段光滑条件只要求每段参数有连续导数，使积分能够逐段定义。尖角不产生额外项。

\subsection{原函数}
\label{b3:subsec:020103}

\begin{proposition}[原函数的端点公式]\label{b3:prop:02-primitive-endpoints}
设 \(U\) 为开集，\(F\) 在 \(U\) 上全纯，\(F'=f\)。若 \(\gamma\) 是 \(U\) 内由 \(p\) 到 \(q\) 的分段光滑曲线，则
\[
 \int_\gamma f(z)\,\mathrm{d}z=F(q)-F(p).
\]
\end{proposition}
\begin{proof}
全纯函数的导数连续已由\cref{b3:thm:01-holomorphic-analytic}得到。在每个光滑参数区间上，链式法则给出 \((F\circ\gamma)'=f(\gamma)\cdot\gamma'\)。分别对实部和虚部应用微积分基本定理，再把有限个区间的端点差相加，中间端点相消。
\end{proof}
例如 \(z^n\) 在 \(\mathbb C\) 上有原函数 \(z^{n+1}/(n+1)\)，其中 \(n\ge0\)。函数 \(1/z\) 在任何避开零点的小圆盘上都有对数原函数，却在整个 \(\mathbb C\setminus\{0\}\) 上没有原函数：单位圆上的积分为 \(2\pi\mathrm{i}\)，违反端点公式。局部存在与全局存在的差别，来自闭曲线能否收缩而不是函数是否局部光滑。

\ProseParagraphBreak

为判断路径独立性，可以先在一个具体区域计算闭路周期。例如在右半平面 \(\operatorname{Re}z>0\)，第一章的对数主值满足 \((\operatorname{Log}z)'=1/z\)，所以由 \(p\) 到 \(q\) 的每条路径都给 \(\operatorname{Log}q-\operatorname{Log}p\)。在去心平面，沿从1到 \(-1\) 的上半圆与下半圆积分分别为 \(\pi\mathrm{i}\) 和 \(-\pi\mathrm{i}\)；拼合得到一个单位圆周期。两个端点相同而积分不同，正是不存在单值原函数的具体表现。

\ProseParagraphBreak

若定义域不连通，原函数在每个连通分支上可以任意加一个常数，不存在跨分支的端点公式。连通分支在平面开集中仍然是开集；构造原函数时固定的起点应逐分支选择。后文说“原函数相差常数”时，常数的作用域也必须是一个区域，而不是任意开集。

% !TeX root = ../../Book3.tex
\section{Cauchy 积分定理}
\label{b3:sec:0202}

局部 Taylor 展开给出小圆盘内的原函数。若曲线穿过许多圆盘，不同原函数的端点差能否接起来？若曲线围住一个孔洞，它们又为什么可能留下非零总和？本节先把积分沿连续路径定义清楚，再用参数方格说明连续变形怎样保持积分，最后把路径独立性转回原函数的存在。

\subsection{三角形积分与局部原函数}
\label{b3:subsec:020201}

若 \(f\) 在包含一个闭三角形 \(T\) 的开集上全纯，\cref{b3:lem:01-goursat-local}已经证明 \(\int_{\partial T}f(z)\,\mathrm{d}z=0\)。现在全纯即解析已经建立，可以直接逐项积分，得到小圆盘上的原函数。

\Needspace{6\baselineskip}
\begin{lemma}[局部原函数]\label{b3:lem:02-local-primitive}
设 \(a\in\mathbb C\)、\(r>0\)，\(f:D(a,r)\to\mathbb C\) 全纯。则存在全纯函数 \(F:D(a,r)\to\mathbb C\)，满足 \(F'=f\)；还可要求 \(F(a)=0\)，此时 \(F\) 唯一。
\end{lemma}
\begin{proof}
由\cref{b3:thm:01-holomorphic-analytic}，有
\[
 f(z)=\sum_{n=0}^{\infty}c_n(z-a)^n
 \quad(|z-a|<r).
\]
令
\[
 F(z)=\sum_{n=0}^{\infty}\frac{c_n}{n+1}(z-a)^{n+1}.
\]
对每个 \(0<\rho<r\)，这个级数的绝对值在 \(|z-a|\le\rho\) 上由
\(\rho\sum_{n\ge0}|c_n|\rho^n\) 控制，因而在每个较小闭圆盘上一致绝对收敛。\Cref{b3:thm:01-power-differentiation}允许逐项求导，得到 \(F'=f\)，而 \(F(a)=0\) 直接来自级数。两个这样的原函数之差导数为零；圆盘连通，故其差为常数，再在 \(a\) 处求值即得唯一性。
\end{proof}

在一般开集 \(U\) 内，每个点都有闭包仍包含于 \(U\) 的小圆盘，所以局部原函数总能取得。不过，相邻圆盘上的原函数未必以同一常数归一化。积分只使用端点差，常数在每段内都会消失；这使我们不必先解决全局拼合，就能沿一条路径继续计算。

\subsection{连续路径上的局部积分}
\label{b3:subsec:020202}

连续变形产生的中间路径未必分段光滑，甚至可能没有有限长度。此时不能把不存在的导数放进参数积分。但是只要被积函数在每个小圆盘有原函数，沿一段路径的贡献已经由端点决定。这给出一种不对路径求导的定义。

\ProseParagraphBreak

先说明怎样选出有限分段。紧度量空间 \(K\) 的任意开覆盖都有一个数 \(\delta>0\)，使直径小于 \(\delta\) 的每个非空子集落在一个覆盖成员中。这样的数称为这个覆盖的\term[lebesgue-shu]{Lebesgue 数}[Lebesgue number]。为证明它，对每个 \(x\in K\) 选 \(r_x>0\)，使相对球 \(B_K(x,2r_x)\) 落在某个覆盖成员内；由紧性，从 \(B_K(x,r_x)\) 中取有限子覆盖，并令 \(\delta\) 为所取有限个半径的最小值。若小集合含有一点 \(y\in B_K(x,r_x)\)，则其中任一点到 \(x\) 的距离小于 \(\delta+r_x\le2r_x\)，故整个集合属于同一个覆盖成员。下面分别把 \(K\) 取成参数区间与参数方块。

\Needspace{6\baselineskip}
\begin{lemma}[局部积分与连续路径]\label{b3:lem:02-continuous-path-integral}
设 \(U\subseteq\mathbb C\) 为开集，\(f:U\to\mathbb C\) 全纯，\(\gamma:[0,1]\to U\) 连续。可以选择分割 \(0=t_0<\cdots<t_N=1\) 以及圆盘 \(D_k\subset U\)，使 \(\gamma([t_{k-1},t_k])\subset D_k\)。在每个 \(D_k\) 上选择 \(f\) 的一个原函数 \(F_k\)，令
\[
 I_f(\gamma)\coloneqq\sum_{k=1}^N
 \bigl(F_k(\gamma(t_k))-F_k(\gamma(t_{k-1}))\bigr).
\]
这个数与分割、圆盘和原函数的选择无关。它在严格递增连续双射换参下不变，在路径拼接下可加，在反向时变号；若 \(\gamma\) 分段光滑，则 \(I_f(\gamma)=\int_\gamma f(z)\,\mathrm{d}z\)。
\end{lemma}
\begin{proof}
对紧像集中的每一点选一个包含于 \(U\) 的圆盘。圆盘的逆像组成 \([0,1]\) 的开覆盖，取其 Lebesgue 数 \(\delta>0\)，再选择每段长度小于 \(\delta\) 的有限分割。每段路径的像就完整地落在一个圆盘中；前一引理提供该圆盘上的原函数。

固定一套圆盘和原函数时，加入分点只是把一个端点差拆成两个，所以不改变和。现在比较同一参数段所能使用的两个圆盘 \(D,E\)，并设它们的原函数为 \(F,G\)。整段路径的像包含在 \(D\cap E\) 中。两个圆盘的非空交是凸集，因而连通；\((F-G)'=0\)，所以 \(F-G\) 在交集中为常数。于是
\[
 F(\gamma(t_k))-F(\gamma(t_{k-1}))
 =G(\gamma(t_k))-G(\gamma(t_{k-1})).
\]
对任意两套选择，把两套分割的全部分点合并。细分后的每一小段同时处于原来两套所选的圆盘中，逐段应用这个等式，再相加，就得到相同的总和。

对严格递增连续双射换参，把原分点换成其逆像，所有端点及次序不变，故总和不变。反向使每段的两个端点交换，拼接使两组端点差依次排列，因此分别变号与相加。若路径分段光滑，把其光滑分点加入分割；每段的端点差由\cref{b3:prop:02-primitive-endpoints}等于通常积分，所以两种定义一致。
\end{proof}

以下仍将这个数写成 \(\int_\gamma f(z)\,\mathrm{d}z\)。它针对具有局部原函数的被积函数定义，不是在任意不可求长路径上积分任意连续函数的方法。对这里的全纯函数，每一小段还可替换成连接其端点的直线段：圆盘凸性使线段留在原圆盘内，端点公式保证积分相同。这样得到一条有限折线路径，适合实际计算。

\ProseParagraphBreak

不同圆盘的端点差不能未经核对就全部望远镜相消。以 \(1/z\) 沿单位圆为例，每个短圆弧都可选连续辐角，并用相应对数作原函数；沿逆时针方向依次衔接后，辐角从 \(0\) 增加到 \(2\pi\)，端点差之和为 \(2\pi\mathrm{i}\)。若存在同一个全局原函数，起点与终点相同才会迫使总和为零。下一小节说明何种路径变形可以保证这种总和消失。

\subsection{单连通区域}
\label{b3:subsec:020203}

用连续映射 \(H:[0,1]^2\to U\) 表示路径的连续变形：固定第二个参数 \(t\) 后，\(s\mapsto H(s,t)\) 是变形到时刻 \(t\) 的路径。这种变形称为\term[tonglun]{同伦}[homotopy]。若所有路径具有固定的起点和终点，便有 \(H(0,t)=p\)、\(H(1,t)=q\)；若只要求每条路径闭合，则要求 \(H(0,t)=H(1,t)\)，共同端点可以移动。

\Needspace{6\baselineskip}
\begin{proposition}[同伦下的积分不变性]\label{b3:prop:02-homotopy-invariance}
设 \(U\subseteq\mathbb C\) 为开集，\(f:U\to\mathbb C\) 全纯，\(H:[0,1]^2\to U\) 连续。定义四条路径
\[
 \gamma_j(s)=H(s,j),
 \quad \eta_j(t)=H(j,t)
 \quad(j=0,1).
\]
沿连续路径的积分按\cref{b3:lem:02-continuous-path-integral}定义，则
\[
 \int_{\gamma_0}f(z)\,\mathrm{d}z
 +\int_{\eta_1}f(z)\,\mathrm{d}z
 -\int_{\gamma_1}f(z)\,\mathrm{d}z
 -\int_{\eta_0}f(z)\,\mathrm{d}z=0.
\]
特别地，若 \(H(0,t)\)、\(H(1,t)\) 分别与 \(t\) 无关，或对所有 \(t\) 有 \(H(0,t)=H(1,t)\)，则
\(\int_{\gamma_0}f(z)\,\mathrm{d}z=\int_{\gamma_1}f(z)\,\mathrm{d}z\)。端路径分段光滑时，这些积分等于通常的曲线积分。
\end{proposition}
\begin{proof}
用包含于 \(U\) 的原函数圆盘覆盖紧集 \(H([0,1]^2)\)。其逆像是参数方块的开覆盖，取 Lebesgue 数 \(\delta>0\)。选择整数 \(N\) 使 \(\sqrt2/N<\delta\)，将方块分成边长 \(1/N\) 的 \(N^2\) 个闭小方格。每格的直径小于 \(\delta\)，故整格的像位于一个原函数圆盘内。

取一个小格，按逆时针顺序记其四个顶点为 \(v_0,v_1,v_2,v_3\)，并记 \(v_4=v_0\)。若 \(F\) 是覆盖该格像的原函数，则四条边经 \(H\) 映成连续路径，其积分总和是
\[
 \sum_{k=0}^3\bigl(F(H(v_{k+1}))-F(H(v_k))\bigr)=0.
\]
这个计算不要求 \(H\) 可微或单射，也不要求格子的像围成一个区域。

把所有小格的等式相加。每条内部边恰出现两次，一次从某个顶点走向另一个顶点，一次反向行走；其像是同一条参数路径的两个方向。即使相邻格子使用不同原函数，前一引理的良定义性也保证两次积分互为相反数，所以它们相消。余下外边界的走向依次为：底边从左到右，即 \(\gamma_0\)；右边从下到上，即 \(\eta_1\)；顶边从右到左，即 \(\gamma_1^-\)；左边从上到下，即 \(\eta_0^-\)。这恰好给出所写等式。

定端点时，两条竖边映为常路径，积分都是零。闭路径同伦时，两条竖边映为同一连续路径，故两个积分抵消。两种情形都剩下 \(\gamma_0\) 与 \(\gamma_1\) 的积分相等。
\end{proof}

这里的紧性用于找到一个统一的细分尺度；局部原函数用于计算每格的四个端点差；反向与拼接规则用于消去公共边。无须对 \(H\) 求偏导，也无须把弯曲或重叠的格子像看作可以使用平面面积公式的区域。

\Needspace{6\baselineskip}
\begin{definition}\label{b3:def:02-simply-connected}
区域 \(U\subseteq\mathbb C\) 称为\term[danliantong]{单连通}[simply connected]，是指对每条闭连续路径 \(\gamma:[0,1]\to U\)，都存在点 \(p\in U\) 及连续映射 \(H:[0,1]^2\to U\)，满足
\[
 H(s,0)=\gamma(s),
 \quad H(s,1)=p,
 \quad H(0,t)=H(1,t)
 \quad(0\le s,t\le1).
\]
也就是说，每条闭路径都能在 \(U\) 内经过闭路径连续收缩到一点。
\end{definition}

每个非空凸开集都单连通：固定 \(p\in U\)，取
\(H(s,t)=(1-t)\gamma(s)+tp\)，凸性保证所有变形点仍在 \(U\) 内，闭合关系由 \(\gamma(0)=\gamma(1)\) 保留。圆盘、半平面都属此例。在一般区域，是否存在这样的收缩就是积分定理所需的拓扑条件。

\ProseParagraphBreak

\term[cauchy-jifen-dingli]{Cauchy 积分定理}[Cauchy's integral theorem]现在由同伦不变性直接得到。

\Needspace{6\baselineskip}
\begin{theorem}[Cauchy，单连通区域]\label{b3:thm:02-cauchy-simply-connected}
设 \(U\subseteq\mathbb C\) 是单连通区域，\(f:U\to\mathbb C\) 全纯。则 \(U\) 内每条闭分段光滑曲线 \(\gamma\) 都满足
\[
 \int_\gamma f(z)\,\mathrm{d}z=0.
\]
\end{theorem}
\begin{proof}
按定义，把 \(\gamma\) 在 \(U\) 内通过闭路径连续收缩到常路径。由\cref{b3:prop:02-homotopy-invariance}，其积分等于常路径的积分，后者为零。连续路径积分与 \(\gamma\) 的通常曲线积分一致，所以得到结论。
\end{proof}

\begin{figure}[htbp]
\centering
\begin{tikzpicture}[>=stealth,scale=0.95,font=\small]
 \draw[step=0.75,gray!65] (0,0) grid (3,3);
 \fill[gray!20] (0.75,0.75) rectangle (1.5,1.5);
 \draw (0,0) rectangle (3,3);
 \draw[->,thick] (0.7,0)--(1.5,0);
 \draw[->,thick] (3,0.8)--(3,1.5);
 \draw[->,thick] (2.3,3)--(1.5,3);
 \draw[->,thick] (0,2.2)--(0,1.5);
 \draw[->] (1.44,0.87)--(1.44,1.37);
 \draw[->] (1.56,1.37)--(1.56,0.87);
 \node[below] at (1.5,0) {\(H(s,0)=\gamma(s)\)};
 \node[above] at (1.5,3) {\(H(s,1)=p\)};
 \node[rotate=90,above] at (0,1.5) {\(H(0,t)\)};
 \node[rotate=90,below] at (3,1.5) {\(H(1,t)\)};
 \node at (1.12,1.12) {\(Q\)};
 \draw[->,thick] (3.6,1.6)--(4.7,1.6) node[midway,above] {\(H\)};
 \draw (6,1.6) circle (1.1);
 \draw[fill=gray!20] (5.45,1.35)..controls (5.7,2.05) and (6.45,2.2)..(6.55,1.6)
   ..controls (6.4,1.05) and (5.75,0.95)..cycle;
 \node at (6,1.55) {\(H(Q)\)};
 \node[above] at (6,2.7) {原函数圆盘 \(D_j\)};
\end{tikzpicture}
\caption[同伦方格中的边界相消]{细分的是参数方块，不是 \(U\)。每格的像包含于某个原函数圆盘；右图仅示意包含关系，像不必是简单区域。相邻小格的公共边方向相反，外侧两竖边的像相同而方向相反，顶边的像为一点，最后只剩底边 \(\gamma\)。}
\label{b3:fig:02-homotopy-grid}
\end{figure}

\Cref{b3:fig:02-homotopy-grid}画出收缩同伦的情形。细分发生在参数方块，灰格的像只需位于一个原函数圆盘内；它不必是简单图形。闭合关系使外侧两竖边的积分抵消，常值顶边没有贡献，所以最后底边积分为零。

\ProseParagraphBreak

去心平面不是单连通的，因为若单位圆能在其中收缩，对全纯函数 \(1/z\) 使用同伦不变性，就会得到 \(2\pi\mathrm{i}=0\)。同样，在圆环 \(r<|z|<R\) 内，绕中心的圆周不能收缩到点，但半径介于 \(r\) 和 \(R\) 的任意两个正向圆周可以相互变形：
\[
 H(s,t)=\bigl((1-t)\rho_0+t\rho_1\bigr)e^{2\pi\mathrm{i}s},
 \quad r<\rho_0,\rho_1<R.
\]
所以任意在圆环内全纯的函数沿这两条圆周的积分相同。积分不变并不总意味着积分为零；只有能变形成常路径时，才取得零值。

\subsection{原函数判别与 Morera 定理}
\label{b3:subsec:020204}

若所有闭路积分都为零，从一个固定起点出发积分，就不会因路径选择不同而得到两个值。这给出了寻找全局原函数的自然构造。所需的被积函数条件可以放宽为连续，全纯性会成为构造的结果。

\Needspace{6\baselineskip}
\begin{proposition}[原函数判别准则]\label{b3:prop:02-primitive-criterion}
设 \(U\subseteq\mathbb C\) 为区域，\(f:U\to\mathbb C\) 连续。则 \(f\) 在 \(U\) 上有全纯原函数，当且仅当它沿 \(U\) 内每条闭分段光滑曲线的积分为零。上述条件还等价于：任意两条起点相同、终点相同的分段光滑路径给出相同的积分。
\end{proposition}
\begin{proof}
若有原函数，端点公式立即给出路径独立性，也给出闭路积分为零。路径独立性反过来蕴含闭路积分为零，因为闭路可以与同一起终点的常路径比较。以下由闭路积分消失构造原函数。

固定 \(a\in U\)。先验证每一点都可从 \(a\) 沿有限折线路径到达。记这些可达点组成的集合为 \(E\)，则 \(a\in E\)。对任意 \(z\in U\)，取小圆盘 \(D(z,r)\subset U\)。若圆盘内有一个可达点，就可以从该点续接线段到圆盘内任意点，所以整个圆盘都属于 \(E\)；若没有可达点，整个圆盘属于补集。因此 \(E\) 与 \(U\setminus E\) 都在 \(U\) 中开，由连通性得到 \(E=U\)。

对 \(z\in U\)，选一条由 \(a\) 到 \(z\) 的折线路径，定义 \(F(z)\) 为其上 \(f\) 的积分。两条这样的路径与其中一条的反向拼成闭路，闭路积分为零，所以定义与选择无关。固定 \(z\) 并取 \(D(z,r)\subset U\)，当 \(0<|h|<r\) 时，可在到达 \(z\) 后续接线段到 \(z+h\)，故
\[
 \frac{F(z+h)-F(z)}h
 =\int_0^1f(z+th)\,\mathrm{d}t.
\]
右边与 \(f(z)\) 之差的模不超过
\(\sup_{|w-z|\le|h|}|f(w)-f(z)|\)，由连续性趋于零。因此 \(F'(z)=f(z)\)，\(F\) 在 \(U\) 全纯，构造完成。
\end{proof}

结合 Cauchy 积分定理，单连通区域上的每个全纯函数都有原函数。但“某一个函数有原函数”不能反推区域单连通：例如 \(f=1\) 在任意区域都有原函数 \(F(z)=z\)。去心平面上的 \(1/z^2\) 也有原函数 \(-1/z\)，而 \(1/z\) 没有。孔洞只提供可能的障碍，是否产生非零闭路积分还取决于被积函数。

\ProseParagraphBreak

原函数构造还能局部使用。如果只知道所有小三角形的边界积分为零，便在圆盘内以线段连接固定基点；三角形条件恰好保证改变终点时可以续接一条短线段。这就是\term[morera-dingli]{Morera 定理}[Morera's theorem]的证明方法。

\Needspace{6\baselineskip}
\begin{theorem}[Morera]\label{b3:thm:02-morera}
设 \(U\subseteq\mathbb C\) 为开集，\(f:U\to\mathbb C\) 连续。若对每个闭三角形 \(T\subset U\)，都有
\(\int_{\partial T}f(z)\,\mathrm{d}z=0\)，则 \(f\) 在 \(U\) 上全纯。三角形边界取正向。
\end{theorem}
\begin{proof}
任取 \(a\in U\) 及圆盘 \(D(a,r)\subset U\)，对 \(z\) 在圆盘内定义
\(F(z)=\int_{[a,z]}f(\zeta)\,\mathrm{d}\zeta\)。当 \(z+h\) 仍在圆盘内时，凸性保证以 \(a,z,z+h\) 为顶点的三角形位于圆盘中。它的边界积分为零，因而
\[
 F(z+h)-F(z)=\int_{[z,z+h]}f(\zeta)\,\mathrm{d}\zeta.
\]
三点共线时，同一式子直接由有向线段积分的可加性得到。于是对 \(h\ne0\)，差商仍为 \(\int_0^1f(z+th)\,\mathrm{d}t\)，由连续性趋于 \(f(z)\)。所以 \(F\) 全纯且 \(F'=f\)。由第一章已经证明的全纯即解析，\(F\) 的导数也全纯，从而 \(f\) 在这个圆盘上全纯。每一点都能选择这样的圆盘，故结论在 \(U\) 上成立。
\end{proof}

连续性在最后的差商极限中不可省略。例如在复平面上令 \(f(0)=1\)，其余点取 \(f(z)=0\)。任意非退化三角形的各条边最多在一个参数点遇到原点，改变有限个参数点的值不改变实积分，所以所有三角形边界积分仍为零；但 \(f\) 在原点不连续，更不全纯。积分条件需要连续性才能恢复原来的点值。

\ProseParagraphBreak

接下来把上下两侧的全纯函数跨过一条直线粘合。这里不能对交界处的未知导数直接使用 Cauchy--Riemann 方程；改用 Morera，只须证明切割后的边界积分能在极限中相接。

\Needspace{6\baselineskip}
\begin{proposition}[跨直线粘合]\label{b3:prop:02-line-gluing}
设 \(U\subseteq\mathbb C\) 为开集，\(f:U\to\mathbb C\) 连续，并在 \(U\setminus\mathbb R\) 上全纯。则 \(f\) 在 \(U\) 上全纯。
\end{proposition}
\begin{proof}
由 Morera 定理，只需证明任意闭三角形 \(T\subset U\) 的边界积分为零。先记录连续性对短距离移动线段的作用。设 \(a_m\to a\)、\(b_m\to b\)，且全部线段都位于紧集 \(T\) 内，取线性参数化
\(\sigma_m(t)=a_m+t(b_m-a_m)\)、\(0\le t\le1\)。它一致趋于 \(\sigma(t)=a+t(b-a)\)，导数常值也趋于 \(b-a\)。由于 \(f\) 在 \(T\) 上一致连续且有界，
\[
 \int_{[a_m,b_m]}f(z)\,\mathrm{d}z
 =(b_m-a_m)\int_0^1 f(\sigma_m(t))\,\mathrm{d}t
 \longrightarrow\int_{[a,b]}f(z)\,\mathrm{d}z.
\]
若 \(a=b\)，右端按长度为零的常线段解释，结论仍成立。

令 \(T_+=T\cap\{\operatorname{Im}z\ge0\}\)，先假定它有非空内部。对充分小的 \(\varepsilon>0\)，令
\[
 P_\varepsilon=T\cap\{\operatorname{Im}z\ge\varepsilon\}.
\]
这是一个非退化凸多边形，闭包包含于 \(U\setminus\mathbb R\)。从其中一个顶点向不相邻顶点连线，可将它剖成有限个三角形；\cref{b3:lem:01-goursat-local}使每个三角形的边界积分为零，内部公共边相消，得到 \(\int_{\partial P_\varepsilon}f(z)\,\mathrm{d}z=0\)。

当 \(\varepsilon\downarrow0\) 时，\(P_\varepsilon\) 的各条边由原三角形边的截取部分及一条水平截线组成。各端点沿固定线段连续趋向 \(T_+\) 的相应边端点；某条边若缩成一点，其积分趋于零。新增水平截线趋于 \(T_+\) 在实轴上的底边，或者趋于一个顶点。对这有限条线段分别应用刚证的端点连续性，并保留逆时针方向，得到
\(\int_{\partial T_+}f(z)\,\mathrm{d}z=0\)。若 \(T_+\) 完全位于上半平面，则充分小时截取不再改变它，同一结论更直接。

对 \(T_-=T\cap\{\operatorname{Im}z\le0\}\)，用水平线 \(\operatorname{Im}z=-\varepsilon\) 从下侧截取，得到同样的零积分结论。空的部分不作贡献；若某一部分退化成线段或点，其正反两向积分相消。最后将上下两部分的边界积分相加，公共实线段分别向左、向右行走，因而完全抵消；其余边正好组成 \(\partial T\)。于是 \(\int_{\partial T}f(z)\,\mathrm{d}z=0\)，Morera 定理给出全纯性。
\end{proof}

交界处函数值的连续性使两侧的积分贡献相消。全纯性成立以后，所有导数的相容性才由解析性得到。若上侧取常数 \(1\)、下侧取常数 \(0\)，两侧都全纯，却无法连续粘合，恰好缺少命题的假设。对任意其他直线，先作平移和非零复数旋转把它变成实轴，再应用同一结论即可；后面的反射原理将用到这种局部粘合。

% !TeX root = ../../Book3.tex
\section{Cauchy 积分公式}
\label{b3:sec:0203}

积分定理消去全纯被积函数的闭路积分。要从边界恢复内部一点 \(z\) 的函数值，就把被积函数改成 \(f(\zeta)/(\zeta-z)\)：它在 \(z\) 处留下一个不能消去的贡献。这个贡献可以压缩到任意小的圆周，再由连续性算出。本节先说明一般边界的消边步骤，随后研究对积分中的内部变量求导，以及由此产生的定量估计。

\subsection{基本公式}
\label{b3:subsec:020301}

前节的同伦证明适用于连续变形，不必预先给区域剖分。这里处理实际的区域边界时，使用多元微积分中的平面 Green 公式。所需版本如下：若 \(W\) 是有界平面区域，边界由有限条互不相交的分段光滑 Jordan 曲线组成，实函数 \(P,Q\) 在 \(\overline W\) 的开邻域为 \(C^1\)，则
\[
 \int_{\partial W}P\,\mathrm{d}x+Q\,\mathrm{d}y
 =\iint_W(Q_x-P_y)\,\mathrm{d}x\,\mathrm{d}y.
\]
这里 Jordan 曲线指除起终点相同外没有自交的闭曲线。边界取正向，即沿边界行走时区域在左侧；因此外边界逆时针，孔洞边界顺时针。第二卷定理 1.2.4 已给出 \(C^1\) 边界版本；本节采用上面允许分段光滑 Jordan 边界的经典 Green 公式作为实变量微积分先修，下面完整写出它在复积分中的应用。第一章的 Goursat 定理和全纯即解析并未依赖这个公式；此时全纯函数已经具有连续偏导数，才可使用这一版本。
% 跨卷背景定位：Book2/Part01/Chapter01/Section0102.tex，b2:thm:01-vector-green-plane，当前定理1.2.4；该处只承担C1边界版，不冒称已证明本节全部分段光滑Jordan边界版。

\Needspace{6\baselineskip}
\begin{lemma}[区域边界上的 Cauchy 定理]\label{b3:lem:02-cauchy-boundary}
设 \(W\subseteq\mathbb C\) 是有界区域，边界由有限条互不相交的分段光滑 Jordan 曲线组成，均按区域在左侧定向。若 \(h\) 在 \(\overline W\) 的一个开邻域上全纯，则
\[
 \int_{\partial W}h(\zeta)\,\mathrm{d}\zeta=0,
\]
其中左端表示沿所有边界分支的积分之和。
\end{lemma}
\begin{proof}
写 \(h=u+\mathrm{i}v\)、\(\zeta=x+\mathrm{i}y\)。第一章的解析性定理保证 \(u,v\) 为 \(C^1\)，而 Cauchy--Riemann 方程给出 \(u_x=v_y\)、\(u_y=-v_x\)。展开微分式并分别应用实 Green 公式，有
\[
 \begin{aligned}
 \int_{\partial W}h(\zeta)\,\mathrm{d}\zeta
 &=\int_{\partial W}(u\,\mathrm{d}x-v\,\mathrm{d}y)
   +\mathrm{i}\int_{\partial W}(v\,\mathrm{d}x+u\,\mathrm{d}y)\\
 &=\iint_W(-v_x-u_y)\,\mathrm{d}x\,\mathrm{d}y
   +\mathrm{i}\iint_W(u_x-v_y)\,\mathrm{d}x\,\mathrm{d}y=0.
 \end{aligned}
\]
两个面积积分的被积函数分别为零。孔洞的作用已经包含在正向边界的约定中，不需要把内边界遗漏或改成与外边界同向。
\end{proof}

这一引理允许直接在一个区域内挖去小圆盘。挖去的圆盘应当闭包含于原区域，函数则须在留下的闭区域附近全纯；原来的奇点可以落在挖去部分。

\Needspace{6\baselineskip}
\begin{theorem}[Cauchy 积分公式]\label{b3:thm:02-cauchy-formula}
设 \(V\subseteq\mathbb C\) 是有界单连通区域，其边界是一条正向分段光滑 Jordan 曲线，\(f\) 在 \(\overline V\) 的一个开邻域上全纯。则对每个 \(z\in V\)，
\[
 f(z)=\frac1{2\pi\mathrm{i}}\int_{\partial V}
       \frac{f(\zeta)}{\zeta-z}\,\mathrm{d}\zeta.
\]
\end{theorem}
\begin{proof}
固定 \(z\in V\)，选 \(\varepsilon>0\) 使 \(\overline{D(z,\varepsilon)}\subset V\)。函数 \(h(\zeta)=f(\zeta)/(\zeta-z)\) 在挖去这个圆盘后留下的闭区域附近全纯。前一引理给出
\[
 \int_{\partial V}\frac{f(\zeta)}{\zeta-z}\,\mathrm{d}\zeta
 =\int_{|\zeta-z|=\varepsilon}
       \frac{f(\zeta)}{\zeta-z}\,\mathrm{d}\zeta,
\]
右侧的小圆此时取逆时针方向。负号已经在将顺时针内边界移到等式右侧时消去。

小圆的参数式为 \(\zeta=z+\varepsilon e^{\mathrm{i}t}\)、\(0\le t\le2\pi\)，所以右侧等于
\[
 \mathrm{i}\int_0^{2\pi}f(z+\varepsilon e^{\mathrm{i}t})\,\mathrm{d}t.
\]
它与 \(2\pi\mathrm{i}f(z)\) 的差的模不超过
\[
 2\pi\max_{|\zeta-z|=\varepsilon}|f(\zeta)-f(z)|,
\]
由 \(f\) 在 \(z\) 连续而趋于零。左侧不依赖 \(\varepsilon\)，令 \(\varepsilon\downarrow0\) 即得公式。
\end{proof}

核的大小在小圆上为 \(1/\varepsilon\)，圆长却为 \(2\pi\varepsilon\)，两者抵消后留下有限的点值贡献。因此不能仅因小圆周长趋零就把其积分删去。若将核换成在 \(z\) 附近有界的函数，长度估计才会使这项积分趋于零。

\ProseParagraphBreak

“在闭区域的邻域全纯”保证边界积分有定义，也保证挖洞后可以应用 Green 公式。若只知道 \(f\) 在 \(V\) 内全纯，仍可在闭包含于 \(V\) 的小圆盘上使用第一章的公式；在没有边界值或边界长度不受控制时，不能直接把围道推到 \(\partial V\)。例如 \(1/(1-z)\) 在单位圆盘全纯，但单位圆上的被积式在 \(z=1\) 已经没有通常的连续曲线积分意义。

\ProseParagraphBreak

外边界公式也能计算不以所求点为中心的围道。例如在正向矩形
\(-2<\operatorname{Re}\zeta<2\)、\(-1<\operatorname{Im}\zeta<1\) 的边界上，
\[
 \int\frac{e^\zeta}{\zeta-\tfrac12}\,\mathrm{d}\zeta
 =2\pi\mathrm{i}e^{1/2}.
\]
这里无需分别计算四条边；被积式在矩形内只有分母指定的一个异常点。若把分母改成 \(\zeta-3\)，被积函数就在整个闭矩形附近全纯，边界积分为零。点是否位于围道内部是公式的一项条件，不能从分母的形式省略。

\Needspace{6\baselineskip}
\begin{proposition}[圆环上的 Cauchy 公式]\label{b3:prop:02-annulus-cauchy}
设 \(a\in\mathbb C\)、\(0<r<R<\infty\)，\(f\) 在闭圆环
\(\{\,\zeta\in\mathbb C:r\le|\zeta-a|\le R\,\}\) 的开邻域上全纯。若 \(r<|z-a|<R\)，则
\[
 f(z)=\frac1{2\pi\mathrm{i}}\int_{|\zeta-a|=R}
       \frac{f(\zeta)}{\zeta-z}\,\mathrm{d}\zeta
       -\frac1{2\pi\mathrm{i}}\int_{|\zeta-a|=r}
       \frac{f(\zeta)}{\zeta-z}\,\mathrm{d}\zeta,
\]
其中两个圆周都取逆时针方向。
\end{proposition}
\begin{proof}
在圆环内部挖去闭包含于其中的 \(D(z,\varepsilon)\)。对剩余区域使用\cref{b3:lem:02-cauchy-boundary}，外圆的积分减去内圆的积分，等于绕 \(z\) 的逆时针小圆积分。后一积分按前一定理的参数计算趋于 \(2\pi\mathrm{i}f(z)\)，即得结论。
\end{proof}

取 \(a=0\)、\(f(\zeta)=1/\zeta\)，可以把孔洞项算得很清楚。由于 \(z\ne0\)，
\[
 \frac1{\zeta(\zeta-z)}
 =\frac1z\left(\frac1{\zeta-z}-\frac1\zeta\right).
\]
在外圆上，括号中两项的积分都是 \(2\pi\mathrm{i}\)，故整个外圆积分为零。在内圆上，第一项在圆盘附近全纯，积分为零；第二项的积分为 \(2\pi\mathrm{i}\)，故整个内圆积分为 \(-2\pi\mathrm{i}/z\)。公式恰好由减去内圆贡献恢复 \(1/z\)。若只保留外圆，将错误地得到零。

\ProseParagraphBreak

圆环公式将在下一章产生 Laurent 展开：外圆适合按 \(z-a\) 的正幂展开核，内圆适合按它的负幂展开核。这里先把两条边界及其符号固定下来，后面的级数计算才有正确起点。

\subsection{导数公式}
\label{b3:subsec:020302}

积分公式把内部变量放在一个简单的有理核中。只要内部点与边界保持正距离，核的差商误差便可在整条边界上一致控制；这正是将求导移入积分号的依据。

\Needspace{6\baselineskip}
\begin{proposition}[Cauchy 导数公式]\label{b3:prop:02-cauchy-derivatives}
设 \(V\subseteq\mathbb C\) 是有界单连通区域，其边界是一条正向分段光滑 Jordan 曲线，\(f\) 在 \(\overline V\) 的开邻域上全纯。则对每个整数 \(n\ge0\) 和 \(z\in V\)，
\[
 f^{(n)}(z)=\frac{n!}{2\pi\mathrm{i}}\int_{\partial V}
        \frac{f(\zeta)}{(\zeta-z)^{n+1}}\,\mathrm{d}\zeta.
\]
\end{proposition}
\begin{proof}
固定非空紧集 \(K\subset V\)，记 \(d=\operatorname{dist}(K,\partial V)>0\)。当 \(z\in K\)、\(|h|<d/2\) 时，
\[
 \frac1h\left(\frac1{\zeta-z-h}-\frac1{\zeta-z}\right)
 =\frac1{(\zeta-z-h)(\zeta-z)}.
\]
右端与 \((\zeta-z)^{-2}\) 的差为
\[
 \frac{h}{(\zeta-z-h)(\zeta-z)^2},
\]
其模至多为 \(2|h|/d^3\)，对所有 \(\zeta\in\partial V\)、\(z\in K\) 同时成立。因此将差商移入积分后，误差的模不超过
\[
 \frac{L(\partial V)\cdot\max_{\partial V}|f|}{2\pi}
       \frac{2|h|}{d^3}\longrightarrow0.
\]
这证明一阶公式。

为使归纳步骤同样明确，令 \(m\ge1\)、\(w=\zeta-z\)。把实变量的二阶积分余项公式用于 \(t\mapsto(w-th)^{-m}\)，得到
\[
 (w-h)^{-m}-w^{-m}-mh\,w^{-m-1}
 =m(m+1)h^2\int_0^1(1-t)(w-th)^{-m-2}\,\mathrm{d}t.
\]
由 \(|w-th|\ge d/2\)，右端的模至多为
\[
 m(m+1)2^{m+1}\frac{|h|^2}{d^{m+2}}.
\]
除以 \(h\) 后误差一致趋零，所以核的导数可以再次移入积分。对 \(m=1,2,\ldots\) 依次进行，每步产生正因子 \(m\)，第 \(n\) 阶的总因子为 \(n!\)。任一内部点都可放入这样的紧集，结论因而成立于整个 \(V\)。
\end{proof}

求导的正号来自内部变量 \(z\)：对 \((\zeta-z)^{-1}\) 求导时，幂函数的负号与内层 \(-z\) 的负号相消。例如对包围 \(a\) 的正向圆周，只要 \(f\) 在闭盘附近全纯，
\[
 \int\frac{f(\zeta)}{(\zeta-a)^3}\,\mathrm{d}\zeta
 =\pi\mathrm{i}f''(a).
\]
取 \(f(\zeta)=\zeta e^\zeta\)，有 \(f''(a)=(a+2)e^a\)，于是积分等于 \(\pi\mathrm{i}(a+2)e^a\)。分母的三次幂对应二阶导数及因子 \(2!\)，不是三阶导数。

\ProseParagraphBreak

同一估计还说明积分可以用来构造全纯函数。设 \(\gamma:[\alpha,\beta]\to\mathbb C\) 为分段光滑曲线，\(g\) 在其像上连续。在开集 \(W=\mathbb C\setminus\gamma([\alpha,\beta])\) 上定义
\[
 H(z)=\int_\gamma\frac{g(\zeta)}{\zeta-z}\,\mathrm{d}\zeta.
\]
对每个非空紧集 \(K\subset W\)，距离 \(d=\operatorname{dist}(K,\gamma([\alpha,\beta]))\) 为正，刚才的差商估计仍然适用，故
\[
 H^{(n)}(z)=n!\int_\gamma
          \frac{g(\zeta)}{(\zeta-z)^{n+1}}\,\mathrm{d}\zeta.
\]
这里 \(g\) 只需沿曲线连续；所求导的变量出现在核中，并没有要求 \(g\) 延拓为全纯函数。

\ProseParagraphBreak

对正向单位圆与 \(g\equiv1\)，\(H\) 在圆内恒为 \(2\pi\mathrm{i}\)，在圆外恒为零。圆内的值来自积分公式；圆外的值来自被积函数在闭单位圆盘附近全纯。两个连通分支上的全纯性并不提供跨越单位圆的连续性。若在曲线上把数据 \(g\) 改为 \(\widetilde g\)，相应函数之差却有直接控制：
\[
 \sup_{z\in K}|H^{(n)}(z)-\widetilde H^{(n)}(z)|
 \le \frac{n!L(\gamma)}{d^{n+1}}
       \max_{\zeta\in\gamma}|g(\zeta)-\widetilde g(\zeta)|.
\]
因此边界数据的小误差能控制远离曲线处的各阶导数；离曲线越近或求导阶数越高，所需的误差控制越强。

\subsection{Cauchy 估计}
\label{b3:subsec:020303}

\term[cauchy-guji]{Cauchy 估计}[Cauchy estimates]把边界上的函数大小变成内部的导数大小。圆周特别适合这一步，因为中心到各边界点的距离相同。

\Needspace{6\baselineskip}
\begin{theorem}[Cauchy 估计]\label{b3:thm:02-cauchy-estimate}
设 \(a\in\mathbb C\)、\(r>0\)，\(f\) 在 \(\overline{D(a,r)}\) 的开邻域上全纯，记 \(M_r=\max_{|z-a|=r}|f(z)|\)。则对每个整数 \(n\ge0\)，
\[
 |f^{(n)}(a)|\le\frac{n!M_r}{r^n}.
\]
\end{theorem}
\begin{proof}
在中心的导数公式中，分母的模为 \(r^{n+1}\)，圆长为 \(2\pi r\)。长度估计给出
\[
 |f^{(n)}(a)|
 \le\frac{n!}{2\pi}\cdot2\pi r\cdot\frac{M_r}{r^{n+1}}
 =\frac{n!M_r}{r^n}.
\]
\end{proof}

若 \(f\) 只在 \(D(a,R)\) 上全纯且 \(|f|\le M\)，对任意 \(r<R\) 应用估计，再令 \(r\uparrow R\)，仍得
\[
 |f^{(n)}(a)|\le\frac{n!M}{R^n}.
\]
这个极限只涉及已经有共同上界的数值，不需要先定义最外圆周上的函数值。估计中的常数不能普遍减小：\(f(z)=M((z-a)/r)^n\) 在取样圆上模恒为 \(M\)，中心的 \(n\) 阶导数恰好达到上界。

\ProseParagraphBreak

把圆心移到待估计的点，就能得到紧集上的统一控制。具体地，设 \(K\) 为非空紧集，\(\delta>0\)，而 \(f\) 在
\[
 K_\delta=\{\,z\in\mathbb C:\operatorname{dist}(z,K)<\delta\,\}
\]
全纯且 \(|f|\le M\)。每个 \(a\in K\) 都有 \(D(a,\delta)\subset K_\delta\)，故 \(|f'(a)|\le M/\delta\)。还要估计邻近两点的函数差时，应确保连接它们的线段也处于可控范围。若 \(a,b\in K\)、\(|a-b|<\delta/2\)，线段上的每个点到 \(K\) 的距离小于 \(\delta/2\)，以它为中心的半径 \(\delta/2\) 圆盘包含于 \(K_\delta\)。于是沿线段有 \(|f'|\le2M/\delta\)，由端点公式，
\[
 |f(b)-f(a)|\le\frac{2M}{\delta}|b-a|.
\]
这说明一族函数若在同一邻域有共同上界，就在 \(K\) 上具有共同的连续性控制。第四章的紧性论证将使用这一机制。

\ProseParagraphBreak

半径有时也需要选择。设整函数满足 \(|f(z)|\le A e^{B|z|}\)，其中 \(A,B>0\)，写 \(f(z)=\sum_{n\ge0}c_nz^n\)。对每个 \(r>0\)，
\[
 |c_n|\le A\frac{e^{Br}}{r^n}.
\]
固定整数 \(n\ge1\)，右端的对数除常数外为 \(Br-n\log r\)。它的导数是 \(B-n/r\)，在 \(r<n/B\) 时为负，在 \(r>n/B\) 时为正，故最佳选择为 \(r=n/B\)，得到
\[
 |c_n|\le A\left(\frac{eB}{n}\right)^n.
\]
这里先固定求导阶数，再为这一阶选择圆周。对 \(f=e^z\) 得到 \(1/n!\le(e/n)^n\)：估计未恢复精确系数，却已经显示高阶系数快速衰减。后面研究整函数增长时，同样的优化会反复出现。

\Needspace{6\baselineskip}
\begin{proposition}[Taylor 截断误差]\label{b3:prop:02-taylor-remainder}
设 \(a\in\mathbb C\)、\(R>0\)、\(M\ge0\)，\(f\) 在闭圆盘 \(\overline{D(a,R)}\) 的邻域全纯，且边界上 \(|f|\le M\)。对整数 \(N\ge0\)，令 \(T_N(z)=\sum_{n=0}^N f^{(n)}(a)(z-a)^n/n!\)。若 \(0\le r<R\)，则
\[
 \sup_{|z-a|\le r}|f(z)-T_N(z)|
 \le \frac{M(r/R)^{N+1}}{1-r/R}.
\]
\end{proposition}
\begin{proof}
第一章保证 \(f\) 在圆盘内等于中心 Taylor 级数。Cauchy 估计给出 \(|f^{(n)}(a)|/n!\le M/R^n\)，故截断之后的尾项模不超过 \(M\sum_{n>N}(r/R)^n\)。求几何级数的和即得结论；\(r=0\) 时余项本身为零。
\end{proof}

这个界与第一章的误差估计一致，现在可以用来反推所需的截断次数。设 \(M>0\)、\(0<r<R\)，记 \(q=r/R\)，希望误差不超过 \(\varepsilon>0\)。充分条件是
\[
 N+1\ge
 \frac{\log\bigl(M/(\varepsilon(1-q))\bigr)}{\log(1/q)}.
\]
若右端为负，任意非负截断次数都满足该充分条件；否则选择满足不等式的整数即可。例如 \(r=R/2\)、\(M=1\)，余项界为 \(2^{-N}\)，取 \(N=10\) 就使误差小于 \(10^{-3}\)。若 \(r\) 接近 \(R\)，相同次数的保证会明显变差，因为几何衰减的比值接近一。

\ProseParagraphBreak

取 \(N=0\)，还得到 \(|f(z)-f(a)|\le M|z-a|/(R-|z-a|)\)。这与先估计导数再沿线段积分是同一局部控制的两种用法。无论采用哪一种，都要明确函数上界在哪个圆盘成立，以及待估计点与取样边界之间还留下多少距离。

% !TeX root = ../../Book3.tex
\section{基本推论}
\label{b3:sec:0204}

Cauchy 估计中可以选择的半径，能把局部导数变成全局增长的探测工具。对整函数把圆周不断扩大，增长过慢会迫使高阶导数消失；在有界区域内，把圆周平均与模的上界比较，又会得到内部极值的限制。本节沿这两个方向证明 Liouville 定理、代数基本定理与最大模原理，最后在单位圆盘中加入中心归一化，得到 Schwarz 引理。一般两点之间的距离估计留待\chapref{b3:ch:06}。

\subsection{Liouville 定理}
\label{b3:subsec:020401}

实轴上的有界光滑函数有很多，例如正弦函数。若同时要求在整个复平面全纯，有界性却会使所有变化消失。\term[liouville-dingli]{Liouville 定理}[Liouville's theorem]的关键是围道可以取得任意大，而函数的模始终受同一个常数控制。

\Needspace{6\baselineskip}
\begin{theorem}[Liouville]\label{b3:thm:02-liouville}
设 \(f:\mathbb C\to\mathbb C\) 为整函数。若存在常数 \(M\ge0\)，使 \(|f(z)|\le M\) 对所有 \(z\in\mathbb C\) 成立，则 \(f\) 为常函数。
\end{theorem}
\begin{proof}
固定任意 \(a\in\mathbb C\)。因为 \(f\) 在整个平面全纯，对每个 \(r>0\) 都能在圆周 \(|z-a|=r\) 上使用\cref{b3:thm:02-cauchy-estimate}，得到
\[
 |f'(a)|\le\frac{M}{r}.
\]
左端与半径无关，令 \(r\to\infty\) 得 \(f'(a)=0\)。\(a\) 任意，故 \(f'\) 在整个平面为零。对任意两点 \(p,q\)，沿连接它们的线段使用\cref{b3:prop:02-primitive-endpoints}，有
\[
 f(q)-f(p)=\int_{[p,q]}f'(z)\,\mathrm{d}z=0.
\]
于是 \(f\) 为常数。
\end{proof}

整性与全平面的共同上界各自不可缺少。函数 \(z\) 在单位圆盘有界且全纯，却不能在该圆盘内把取样半径无限扩大。整函数 \(e^{\mathrm{i}z}\) 在实轴上模恒为一，但 \(|e^{\mathrm{i}(x+\mathrm{i}y)}|=e^{-y}\)，沿负虚轴无界；一条直线上的界不能控制大圆周。Liouville 定理限制的是函数在全平面的同时增长。

\ProseParagraphBreak

同一证明还说明，允许多项式速度的增长，只会留下有限多个 Taylor 系数。这里要固定求导阶数，再把半径放大；不同阶数允许分别取极限。

\Needspace{6\baselineskip}
\begin{proposition}[多项式增长判别]\label{b3:prop:02-polynomial-growth}
设 \(f:\mathbb C\to\mathbb C\) 为整函数，\(m\ge0\) 为整数。若存在常数 \(C\ge0\)，使
\[
 |f(z)|\le C(1+|z|)^m\quad(z\in\mathbb C),
\]
则 \(f\) 是次数不超过 \(m\) 的多项式，允许为零多项式。反之，每个次数不超过 \(m\) 的多项式都满足这样的估计。
\end{proposition}
\begin{proof}
固定整数 \(n>m\)。在以零为中心、半径为 \(r\) 的圆周上，Cauchy 估计给
\[
 \frac{|f^{(n)}(0)|}{n!}
 \le\frac{C(1+r)^m}{r^n}
 =C\left(1+\frac1r\right)^m r^{m-n}\longrightarrow0
 \quad(r\to\infty).
\]
所以所有阶数大于 \(m\) 的 Taylor 系数为零。\(f\) 的 Taylor 展开在每个有限圆盘内都等于 \(f\)，而这些圆盘覆盖整个平面，故
\[
 f(z)=\sum_{n=0}^m\frac{f^{(n)}(0)}{n!}z^n
 \quad(z\in\mathbb C).
\]
反之，若 \(f(z)=\sum_{n=0}^m a_nz^n\)，则 \(|z|^n\le(1+|z|)^m\) 对 \(0\le n\le m\) 成立，从而取 \(C=\sum_{n=0}^m|a_n|\) 即得所需上界。
\end{proof}

这一判别把增长条件变成有限个未知系数。例如设整函数 \(f\) 满足 \(|f(z)|\le C(1+|z|)^2\)，且 \(f(0)=1\)、\(f(1)=f(-1)=2\)。先由命题写成 \(f(z)=a+bz+cz^2\)，再代入三点，得到
\[
 a=1,\quad b+c=1,\quad -b+c=1.
\]
故 \(f(z)=1+z^2\)。若没有增长条件，有限个点值无法唯一决定整函数；在这个解上加 \(z(z^2-1)e^z\)，三个指定点值仍相同。这里真正减少自由度的是全平面的增长控制。

\subsection{代数基本定理}
\label{b3:subsec:020402}

要证明一个多项式具有零点，可以反过来考察假设无零时的倒数。多项式在无穷远处变大，使倒数在外部有界；无零假设则使倒数在剩余紧集上有界。两部分合在一起，便能使用 Liouville 定理。这给出\term[daishu-jiben-dingli]{代数基本定理}[fundamental theorem of algebra]的分析证明。

\Needspace{6\baselineskip}
\begin{theorem}[代数基本定理]\label{b3:thm:02-fundamental-algebra}
设 \(n\ge1\) 为整数，\(P(z)=a_nz^n+\cdots+a_0\) 为复系数多项式，且 \(a_n\ne0\)。则存在复数 \(\alpha_1,\ldots,\alpha_n\)，允许重复，使
\[
 P(z)=a_n\prod_{j=1}^n(z-\alpha_j).
\]
因此 \(P\) 在复平面中计重数恰有 \(n\) 个零点。
\end{theorem}
\begin{proof}
先证明至少有一个零点。反设 \(P(z)\ne0\) 于整个平面，则 \(1/P\) 是整函数。对 \(z\ne0\)，写成
\[
 P(z)=a_nz^n\left(1+\sum_{k=0}^{n-1}\frac{a_k}{a_n}z^{k-n}\right).
\]
括号内有限和的每一项都在 \(|z|\to\infty\) 时趋于零。故可选 \(R>0\)，使
\[
 \sum_{k=0}^{n-1}\left|\frac{a_k}{a_n}\right|\cdot|z|^{k-n}
 \le\frac12\quad(|z|\ge R).
\]
于是外圆之外有
\[
 |P(z)|\ge\frac{|a_n|}{2}\cdot|z|^n,
 \quad\left|\frac1{P(z)}\right|
 \le\frac{2}{|a_n|\cdot R^n}.
\]
在闭圆盘 \(\overline{D(0,R)}\) 上，\(|P|\) 是连续的正函数；紧性使它取得严格正的最小值 \(\delta\)。因而这个圆盘上 \(|1/P|\le1/\delta\)。内外两界合起来说明 \(1/P\) 在整个平面有界。由 Liouville 定理，它为常数；该常数不为零，故 \(P\) 也为常数，与 \(n\ge1\) 矛盾。

既然已有零点 \(\alpha_1\)，多项式除法给出
\(P(z)=(z-\alpha_1)Q(z)\)，其中 \(Q\) 的次数为 \(n-1\)，最高次系数仍为 \(a_n\)。这一步也可逐项使用 \(z^k-\alpha_1^k\) 被 \(z-\alpha_1\) 整除的恒等式，因为余项 \(P(\alpha_1)\) 为零。对 \(Q\) 重复上述存在性结论，有限次之后得到所列的 \(n\) 个一次因子。

最后核对重数。若某个值 \(\alpha\) 在这些因子中出现 \(l\) 次，将它们合并可写成 \(P(z)=(z-\alpha)^lQ_\alpha(z)\)，且 \(Q_\alpha(\alpha)\ne0\)。由\cref{b3:prop:01-zero-factorization}，\(l\) 正是全纯零点的阶数。乘积表示也表明不存在未列出的零点，所以计重数总共为 \(n\)。
\end{proof}

例如 \((z-1)^2(z+2)\) 只有两个不同零点，但在 \(1\) 与 \(-2\) 的阶数分别为二和一，总重数仍为三。定理中的计数不要求各根不同，也不保证它们能用根式写出。

\ProseParagraphBreak

对任意固定 \(w\in\mathbb C\)，多项式 \(P-w\) 仍具有相同的最高次项。定理因此还说明：非常值复多项式把平面满射到平面，对每个值恰有 \(n\) 个计重数的原像。这个结论包含一项全局信息，不能只从局部开映射推出。例如指数函数在每点附近都可逆，其像也开放，却始终遗漏零；多项式的无穷远增长在这里排除了遗漏值。

\subsection{最大模原理}
\label{b3:subsec:020403}

在有界区域内，通常不能把半径推向无穷；可以利用的是圆周平均。全纯函数在中心等于周围值的平均，如果中心的模已经最大，平均中的各个值就没有偏离它的余地。\term[zuidamo-yuanli]{最大模原理}[maximum modulus principle]由此把内部极值与边界控制联系起来。

\Needspace{6\baselineskip}
\begin{theorem}[最大模原理]\label{b3:thm:02-maximum-modulus}
设 \(U\subseteq\mathbb C\) 为区域，\(f:U\to\mathbb C\) 全纯。若 \(|f|\) 在 \(U\) 内某点取得局部最大值，则 \(f\) 为常数。此外，设 \(V\subseteq\mathbb C\) 为非空有界区域，\(F:\overline V\to\mathbb C\) 连续且在 \(V\) 上全纯，则
\[
 \max_{z\in\overline V}|F(z)|
 =\max_{z\in\partial V}|F(z)|.
\]
\end{theorem}
\begin{proof}
设 \(a\in U\) 是 \(|f|\) 的局部最大点，记 \(M=|f(a)|\)。取 \(\varepsilon>0\)，使 \(D(a,\varepsilon)\subseteq U\)，且该圆盘上 \(|f|\le M\)。若 \(M=0\)，则 \(f\) 已在这个圆盘恒为零，由恒等定理在 \(U\) 上也恒为零。

若 \(M>0\)，取单位复数 \(\lambda=\overline{f(a)}/M\)，令 \(g=\lambda f\)。这样 \(g(a)=M\)，并且 \(|g|\le M\) 于 \(D(a,\varepsilon)\)。对每个 \(0<r<\varepsilon\)，圆盘内的 Cauchy 公式给出中心平均
\[
 M=\frac1{2\pi}\int_0^{2\pi}g(a+re^{\mathrm{i}t})\,\mathrm{d}t.
\]
取实部后，连续非负函数 \(M-\operatorname{Re}g(a+re^{\mathrm{i}t})\) 的积分为零，所以它处处为零。圆周上 \(\operatorname{Re}g=M\)，同时 \(|g|\le M\)，从
\(M^2+(\operatorname{Im}g)^2=|g|^2\le M^2\)
可知虚部也为零。因此 \(g=M\) 于这条圆周。让 \(r\) 取遍 \((0,\varepsilon)\)，得到 \(f\) 在整个小圆盘上为常数，再由恒等定理将常值延至 \(U\)。

对于第二项，\(\overline V\) 紧，连续函数 \(|F|\) 在其中取得最大值。如果最大点位于边界，等式立即成立；如果最大点位于内部，第一项使 \(F\) 在 \(V\) 为常数。\(V\) 非空且有界，其边界非空，连续性把同一个常值延续到边界，等式仍成立。
\end{proof}

第一章的开映射定理也解释了这一结论：非常值全纯函数把一个内点邻域映成开集，而任何包含 \(f(a)\) 的小圆盘都有模比 \(|f(a)|\) 更大的点。上面的平均值证明进一步指出取等为何迫使整个小圆周上的函数值相同。

\ProseParagraphBreak

边界最大模结论不要求区域单连通，也不要求边界光滑。例如在圆环 \(V=\{1<|z|<2\}\) 上，\(F(z)=z+1/z\) 全纯并连续到闭包。内圆周上 \(|F|\le2\)，外圆周上 \(|F|\le2+1/2=5/2\)，所以整个闭环上都有 \(|F|\le5/2\)。在 \(z=2\) 时达到 \(5/2\)，故这正是最大值。两条边界都要比较；最大模原理本身不替我们决定哪一条边界控制结果。

\ProseParagraphBreak

区域有界这一条件不能直接删去。在上半平面，函数 \(F(z)=e^{-\mathrm{i}z}\) 连续达到实轴，在实轴上模恒为一，但对 \(z=x+\mathrm{i}y\)、\(y>0\)，有 \(|F(z)|=e^y>1\)。它的增长出现在无穷远处，有限边界上的上界不足以控制内部。另一方面，对有界区域也须具有陈述中的连续边界值，才能在紧闭包上取到最大值。

\ProseParagraphBreak

最大模原理给出边界值的唯一性与稳定性。设 \(V\) 为有界区域，\(F,G\) 在 \(\overline V\) 上连续、在 \(V\) 上全纯，且边界上 \(|F-G|\le\eta\)，其中 \(\eta\ge0\)。将定理用于 \(F-G\)，便有
\[
 \max_{\overline V}|F-G|\le\eta.
\]
特别地，相同的连续边界值至多对应一个全纯函数。然而连续边界数据未必都有全纯解。若单位圆上的数据 \(\overline z\) 是某个 \(h\in C(\overline{\mathbb D})\)、在 \(\mathbb D\) 全纯的函数的边界值，则 \(zh(z)-1\) 在圆周恒为零。唯一性迫使 \(zh(z)-1\) 在圆盘恒零，但代入 \(z=0\) 得到 \(-1=0\)，矛盾。

\ProseParagraphBreak

内部极值还限制插值。例如 \(f\in C(\overline{\mathbb D})\) 在 \(\mathbb D\) 全纯，且圆周上 \(|f|\le1\)。若再指定 \(f(0)=1\)，则 \(f\equiv1\)，不能同时要求 \(f(1/2)=0\)。若只指定 \(f(0)=1/2\)，则没有这项刚性；函数 \((1+z)/2\) 就满足条件且非常值。起作用的是内点已经取到允许的最大模。

\ProseParagraphBreak

最小模的情形需要先排除所考察点为零点。设 \(f\) 在区域 \(U\) 全纯，\(|f|\) 于 \(a\in U\) 取得局部最小值且 \(f(a)\ne0\)。由连续性，\(f\) 在 \(a\) 的一个小圆盘上无零，因此 \(1/f\) 在该圆盘全纯，并在 \(a\) 取得局部最大模。最大模原理使 \(1/f\)、从而 \(f\) 在该圆盘常值，再由恒等定理得 \(f\) 在 \(U\) 常值。若允许 \(f(a)=0\)，函数 \(f(z)=z\) 就在原点取得真正的最小模而不恒定。这项条件仅在使用倒数的邻域内需要，不能用不存在的 \(1/f(a)\) 代替论证。

\subsection{Schwarz 引理初步}
\label{b3:subsec:020404}

对圆盘中的全纯函数，边界上界控制内部模；若还知道中心是一个零点，就应当能把这个零点带来的衰减提出来。\term[schwarz-yinli]{Schwarz 引理}[Schwarz lemma]正是对商 \(f(z)/z\) 使用最大模原理。困难只在于函数未必连续延至单位圆周，因此先在内部半径上估计，再逐渐扩大半径。

\Needspace{6\baselineskip}
\begin{theorem}[Schwarz]\label{b3:thm:02-schwarz}
记 \(\mathbb D=D(0,1)\)。若 \(f:\mathbb D\to\mathbb D\) 全纯且 \(f(0)=0\)，则
\[
 |f(z)|\le|z|\quad(z\in\mathbb D),
 \quad |f'(0)|\le1.
\]
若存在 \(z_0\in\mathbb D\setminus\{0\}\) 使 \(|f(z_0)|=|z_0|\)，或 \(|f'(0)|=1\)，则存在常数 \(c\in\mathbb C\)、\(|c|=1\)，使 \(f(z)=cz\) 于整个 \(\mathbb D\)。反之，每个这样的旋转都取等。
\end{theorem}
\begin{proof}
因为 \(f(0)=0\)，其 Taylor 展开在原点没有常数项，所以
\[
 g(z)=
 \begin{cases}
 f(z)/z,&z\ne0,\\
 f'(0),&z=0
 \end{cases}
\]
在整个 \(\mathbb D\) 全纯：原点附近的级数为 \(f'(0)+f''(0)z/2!+\cdots\)，与去心圆盘上的商相同。固定 \(0<r<1\)，在 \(|z|=r\) 上有 \(|g(z)|\le1/r\)。\(g\) 在闭圆盘 \(\overline{D(0,r)}\) 的邻域全纯，故最大模原理给出
\[
 |g(z)|\le\frac1r\quad(|z|\le r).
\]
现在固定 \(z\in\mathbb D\)，在 \(|z|<r<1\) 中令 \(r\uparrow1\)，得到 \(|g(z)|\le1\)。在 \(z\ne0\) 处乘以 \(|z|\) 得第一项估计，在原点用 \(g(0)=f'(0)\) 得第二项。

若有非零点满足 \(|f(z_0)|=|z_0|\)，则 \(|g(z_0)|=1\)；若 \(|f'(0)|=1\)，则 \(|g(0)|=1\)。两种情形都使 \(g\) 在内部取到其最大模一，所以 \(g\) 是单位模常数 \(c\)，即 \(f(z)=cz\)。反过来，这些函数把 \(\mathbb D\) 映到自身，且两个估计均取等。
\end{proof}

等号条件中的非零点必须保留，因为 \(f(0)=0\) 已使原点的模估计对每个候选函数都取等。如果 \(f\) 不是旋转，证明还给出 \(|f'(0)|<1\) 和 \(|f(z)|<|z|\) 对所有非零 \(z\in\mathbb D\) 成立。不过严格不等号并不自动给出一个在整个圆盘内有效的比例常数 \(q<1\)。例如
\[
 f(z)=\frac{z+z^2}{2}
\]
满足 \(|f(z)|\le |z|\cdot(1+|z|)/2<|z|\) 于 \(0<|z|<1\)，且 \(f'(0)=1/2\)。但当实数 \(r\uparrow1\) 时，\(|f(r)|/r=(1+r)/2\to1\)，所以不存在统一的 \(q<1\) 使整个圆盘上 \(|f(z)|\le q\cdot|z|\)。

\ProseParagraphBreak

中心条件同样有实质作用。常函数 \(f\equiv1/2\) 把单位圆盘映入自身，却不满足 \(|f(z)|\le|z|\) 于原点附近。要处理不固定中心的映射，必须先选取合适的源、目标坐标，把指定点送到原点。第六章将用圆盘自同构完成这一步，得到两点之间的估计。

\ProseParagraphBreak

对于圆心处的比较，仅用平移和缩放即可恢复尺度。设 \(a,b\in\mathbb C\)、\(R,S>0\)，\(F:D(a,R)\to D(b,S)\) 全纯且 \(F(a)=b\)。令
\[
 h(\zeta)=\frac{F(a+R\zeta)-b}{S}\quad(\zeta\in\mathbb D).
\]
它满足 Schwarz 引理的全部条件，故
\[
 |F(z)-b|\le\frac SR\cdot|z-a|,
 \quad |F'(a)|\le\frac SR.
\]
若其中的点值估计在 \(z\ne a\) 取等，或导数估计取等，则 \(h(\zeta)=c\zeta\)、\(|c|=1\)，换回原变量得到
\[
 F(z)=b+c\frac SR(z-a).
\]
因此取等映射正是把两个圆盘作旋转和比例伸缩的仿射映射。

\ProseParagraphBreak

Schwarz 引理还能确定固定圆心的全纯双射。设 \(F:\mathbb D\to\mathbb D\) 全纯且双射，\(F(0)=0\)。由\cref{b3:cor:01-local-injectivity}，\(F'\) 在每点非零，局部逆函数组成全局全纯逆映射 \(H=F^{-1}\)；且 \(H(0)=0\)。分别对 \(F,H\) 使用 Schwarz 引理，有
\[
 |F(z)|\le|z|=|H(F(z))|\le|F(z)|.
\]
所以每个非零点都取等，\(F(z)=cz\)、\(|c|=1\)。归一化后只剩旋转的自由，这将用于共形映射的唯一性；任意圆盘自同构的完整表达式将在后文建立。


\begin{chaptersummary}
曲线积分的局部依据是原函数的端点公式。把连续路径分成落在小圆盘中的有限段，端点差便不依赖分割；再细分同伦参数方块，内部边相消，得到同伦不变性。闭路积分全部为零时，路径积分能够定义全局原函数；局部三角形积分消失则通过 Morera 定理给出全纯性。

\ProseParagraphBreak

积分公式保留一个简单核的点值贡献。小圆的长度与核的增长相互抵消，孔洞边界则按相反方向进入公式。与边界的正距离控制核的差商，因而可以逐次求导，并将函数上界转换为导数及 Taylor 余项的上界。估计之前先选择可用的圆心与半径，极限才具有明确的适用范围。

\ProseParagraphBreak

大圆上的增长控制使高阶 Taylor 系数消失，导向 Liouville 定理与多项式的零点存在。小圆上的平均取等带来最大模刚性；除去原点的零因子，再让取样圆扩大到单位圆盘，则得到 Schwarz 引理。下一章将把这里的内外边界公式展开成 Laurent 级数，研究被挖去的点究竟具有怎样的奇性。
\end{chaptersummary}

\BookExercises

前两题练习参数化与 Cauchy 公式，并辨认变动曲线和变动参数带来的条件变化。\Cref{b3:ex:02-3}与\cref{b3:ex:02-4}把增长控制转成整函数的刚性；\cref{b3:ex:02-5}用两个孔洞检验全局原函数。最后三题分别训练带零点的模比较、高阶消失的 Schwarz 估计，以及最小模原理的边界。

% 本章题组自拟，以正文已建立的工具设计计算、反例、构造与综合训练。
% 保留原8题标签与训练方向；不调用第3章的留数、绕数或Laurent展开。

\Needspace{6\baselineskip}
\begin{exercise}\label{b3:ex:02-1}
设 \(\gamma(t)=2e^{\mathrm{i}t}\)，\(0\le t\le4\pi\)。计算
\[
  \int_\gamma\frac{\mathrm{d}z}{z},
  \quad\int_\gamma\overline z\,\mathrm{d}z.
\]
再对每个正整数 \(N\)，令 \(\gamma_N(t)=2e^{\mathrm{i}t}\)，\(0\le t\le2\pi N\)。在圆周 \(|z|=2\) 上构造连续函数 \(h_N\)，使 \(h_N\) 一致趋于零，但
\(\int_{\gamma_N}h_N(z)\,\mathrm{d}z\) 不趋于零。用曲线长度说明这不违背固定曲线上一致收敛函数可以交换极限与积分的结论。
\end{exercise}
\begin{solution}
由 \(\gamma'(t)=2\mathrm{i}e^{\mathrm{i}t}\)，两个参数积分的被积函数分别为
\[
  \frac{\gamma'(t)}{\gamma(t)}=\mathrm{i},
  \quad \overline{\gamma(t)}\cdot\gamma'(t)=4\mathrm{i}.
\]
因此
\[
  \int_\gamma\frac{\mathrm{d}z}{z}=4\pi\mathrm{i},
  \quad\int_\gamma\overline z\,\mathrm{d}z=16\pi\mathrm{i}.
\]
参数区间的长度是 \(4\pi\)，圆周被走了两遍；只记像集会遗漏这个因子。

在同一个圆周上取
\[
  h_N(z)=\frac1{Nz}.
\]
它连续，并且 \(\max_{|z|=2}|h_N(z)|=1/(2N)\to0\)。但代入 \(\gamma_N\) 得
\[
  \int_{\gamma_N}h_N(z)\,\mathrm{d}z
  =\int_0^{2\pi N}\frac{\mathrm{i}}N\,\mathrm{d}t
  =2\pi\mathrm{i}.
\]
这里 \(L(\gamma_N)=4\pi N\)，所以长度估计给出的误差上界为
\[
  L(\gamma_N)\cdot\max_{|z|=2}|h_N(z)|=2\pi,
\]
并不趋于零。若曲线固定，其有限长度是与 \(N\) 无关的常数，一致趋零就会使同一估计的右端趋零。本例中增长的是行走次数以及由参数化计算的长度，尽管每条曲线的像都相同。
\end{solution}

\Needspace{6\baselineskip}
\begin{exercise}\label{b3:ex:02-2}
对满足 \(|a|\ne2\) 的复数 \(a\)，定义
\[
  I(a)=\int_{|z|=2}\frac{e^z}{z^3(z-a)}\,\mathrm{d}z,
\]
其中圆周取逆时针方向。分别计算 \(0<|a|<2\)、\(|a|>2\) 以及 \(a=0\) 时的值。把圆内参数的答案写成一个在 \(a=0\) 也全纯的表达式，并核对 \(a\to0\) 时的极限。计算不得使用留数理论。
\end{exercise}
\begin{solution}
当 \(a\ne0\) 时，部分分式分解为
\[
  \frac1{z^3(z-a)}
  =-\frac1{a^3z}-\frac1{a^2z^2}
   -\frac1{az^3}+\frac1{a^3(z-a)}.
\]
可通过乘以 \(z^3(z-a)\) 核对恒等式，此时右端各项之和化为 \(1\)。前三项由中心为零的 Cauchy 导数公式给出
\[
  \int_{|z|=2}\frac{e^z}{z}\,\mathrm{d}z=2\pi\mathrm{i},
  \quad\int_{|z|=2}\frac{e^z}{z^2}\,\mathrm{d}z=2\pi\mathrm{i},
  \quad\int_{|z|=2}\frac{e^z}{z^3}\,\mathrm{d}z=\pi\mathrm{i}.
\]
若 \(|a|<2\)，Cauchy 公式给出
\[
  \int_{|z|=2}\frac{e^z}{z-a}\,\mathrm{d}z=2\pi\mathrm{i}e^a;
\]
若 \(|a|>2\)，函数 \(e^z/(z-a)\) 在闭圆盘的邻域全纯，故这一积分为零。乘上部分分式中的相应系数后得到
\[
  I(a)=
  \begin{cases}
    \displaystyle 2\pi\mathrm{i}\,
      \frac{e^a-1-a-a^2/2}{a^3},&0<|a|<2,\\[6pt]
    \displaystyle-2\pi\mathrm{i}\,
      \frac{1+a+a^2/2}{a^3},&|a|>2.
  \end{cases}
\]
两种参数范围的差别在于 \(a\) 是否位于积分圆内，而不是代数分解是否成立。

对圆内的答案，指数级数给出
\[
  \frac{e^a-1-a-a^2/2}{a^3}
  =\sum_{k=0}^{\infty}\frac{a^k}{(k+3)!}\quad(a\ne0).
\]
右端在整个复平面收敛，在零点的值为 \(1/6\)。因此圆内参数的统一表达式是
\[
  I(a)=2\pi\mathrm{i}\sum_{k=0}^{\infty}\frac{a^k}{(k+3)!}
  \quad(|a|<2),
\]
并有 \(\lim_{a\to0}I(a)=\pi\mathrm{i}/3\)。直接令 \(a=0\) 时，原积分为
\[
  \int_{|z|=2}\frac{e^z}{z^4}\,\mathrm{d}z
  =\frac{2\pi\mathrm{i}}{3!}\,(e^z)^{(3)}\big|_{z=0}
  =\frac{\pi\mathrm{i}}3,
\]
与极限相同。部分分式中的各个 \(a^{-1}\) 因子单独发散，却在圆内总和中消去；不能把其中某一项的行为当作整个含参积分的行为。
\end{solution}

\Needspace{6\baselineskip}
\begin{exercise}\label{b3:ex:02-3}
设 \(f:\mathbb C\rightarrow\mathbb C\) 为整函数。假设存在常数 \(C\ge0\) 以及正数列 \(R_j\to\infty\)，使
\[
  \frac1{2\pi}\int_0^{2\pi}|f(R_je^{\mathrm{i}t})|^2\,\mathrm{d}t
  \le C(1+R_j^3)\quad(j\ge1).
\]
证明 \(f\) 是次数不超过 \(1\) 的多项式，其中允许零多项式。说明证明为何不需要在每个大圆上控制最大模。再给出一个非仿射的整函数，它在实轴和虚轴上都有界，说明只在这两条直线上控制函数不能代替题给的圆周平均条件。
\end{exercise}
\begin{solution}
记 \(c_n=f^{(n)}(0)/n!\)。在半径 \(R_j\) 的圆周上使用 Cauchy 导数公式，并令 \(z=R_je^{\mathrm{i}t}\)，得到
\[
  c_n=\frac1{2\pi R_j^n}
        \int_0^{2\pi}f(R_je^{\mathrm{i}t})e^{-n\mathrm{i}t}\,\mathrm{d}t.
\]
这一步直接提取一个 Taylor 系数，没有把平均值换成最大值。由于 \(|e^{-n\mathrm{i}t}|=1\)，实积分的 Cauchy--Schwarz 不等式给出
\[
\begin{aligned}
  |c_n|^2
  &\le\frac1{R_j^{2n}}
       \left(\frac1{2\pi}\int_0^{2\pi}
                |f(R_je^{\mathrm{i}t})|^2\,\mathrm{d}t\right)\cdot
       \left(\frac1{2\pi}\int_0^{2\pi}1\,\mathrm{d}t\right)\\
  &\le C\left(R_j^{-2n}+R_j^{3-2n}\right).
\end{aligned}
\]
固定 \(n\ge2\) 后令 \(j\to\infty\)，右端趋于零，所以 \(c_n=0\)。整函数在任意有限圆盘上都等于以零为中心的 Taylor 级数，故全平面上
\[
  f(z)=c_0+c_1z.
\]
每一阶系数只需沿题给的半径列趋于无穷就能判定为零；圆周上个别点的最大模不在这项证明的假设中。

取 \(g(z)=e^{-z^4}\)。对于实数 \(x,y\)，有
\[
  |g(x)|=e^{-x^4}\le1,
  \quad |g(\mathrm{i}y)|=e^{-y^4}\le1.
\]
但 \(g'(0)=0\)，且 \(g(0)=1\ne e^{-1}=g(1)\)，所以它不是仿射函数：仿射函数若在零点导数为零，就应为常数。两条坐标轴上的有界性没有控制其余方向；例如在 \(z=re^{\pi\mathrm{i}/4}\) 上，有 \(|g(z)|=e^{r^4}\)。
\end{solution}

\Needspace{6\baselineskip}
\begin{exercise}\label{b3:ex:02-4}
设 \(f:\mathbb C\rightarrow\mathbb C\) 为整函数，并存在实常数 \(A\) 与 \(B\ge0\)，使
\[
  \operatorname{Re}f(z)\le A+B\log(1+|z|)
  \quad(z\in\mathbb C).
\]
证明 \(f\) 为常数。特别地，这包含实部只有一个全局上界的情形。若把右端改成 \(A+B|z|\)，结论是否仍成立？
\end{exercise}
\begin{solution}
令 \(g=e^f\)。指数函数与 \(f\) 的复合是整函数，并且处处非零；题给的单侧实部估计化为对整个模的估计
\[
  |g(z)|=e^{\operatorname{Re}f(z)}
  \le e^A(1+|z|)^B.
\]
选取非负整数 \(m\ge B\)，则 \(1+|z|\ge1\) 使
\(|g(z)|\le e^A(1+|z|)^m\)。由\cref{b3:prop:02-polynomial-growth}，\(g\) 是一个多项式。它又处处非零，代数基本定理便排除了正次数的可能；所以 \(g\) 为某个非零常数。

再求导得 \(g'=e^f\cdot f'=0\)。由于 \(e^f\) 没有零点，\(f'=0\) 于全平面；沿任意线段积分可知 \(f\) 为常数。这里从未声称 \(f\) 本身有界，而是把实部上界转成指数的多项式增长，再利用指数不消失这一额外性质。

线性增长的版本不成立。函数 \(f(z)=z\) 非常值，却满足
\(\operatorname{Re}f(z)\le|z|\)，即可以取 \(A=0\)、\(B=1\)。
\end{solution}

\Needspace{6\baselineskip}
\begin{exercise}\label{b3:ex:02-5}
设 \(A,B,C\in\mathbb C\)，在有两个去点的区域
\(U=\mathbb C\setminus\{0,1\}\) 上定义
\[
  h(z)=\frac A z+\frac B{z-1}+\frac C{(z-1)^2}.
\]
计算它沿圆周 \(|z|=1/3\)、\(|z-1|=1/3\) 和 \(|z|=2\) 的积分，三个圆周均取逆时针方向。由此确定 \(h\) 在 \(U\) 上有全局全纯原函数的充要条件，并在存在时写出一个原函数。说明为什么仅检查包围两个去点的外圆积分为零还不够。
\end{exercise}
\begin{solution}
在圆盘 \(|z|<1/2\) 上，后两项全纯，所以它们在 \(|z|=1/3\) 上的积分为零。第一项直接参数化或使用圆盘 Cauchy 公式，得到
\[
  \int_{|z|=1/3}h(z)\,\mathrm{d}z=2\pi\mathrm{i}A.
\]
在 \(|z-1|<1/2\) 上，第一项全纯；第二项贡献 \(2\pi\mathrm{i}B\)。第三项有单值原函数 \(-C/(z-1)\)，所以其闭曲线积分为零。因而
\[
  \int_{|z-1|=1/3}h(z)\,\mathrm{d}z=2\pi\mathrm{i}B.
\]
在外圆 \(|z|=2\) 上，分别对点 \(0\)、\(1\) 使用 Cauchy 公式，两项简单分母的积分为 \(2\pi\mathrm{i}A\) 与 \(2\pi\mathrm{i}B\)；第三项仍由其全局原函数给出零积分。所以
\[
  \int_{|z|=2}h(z)\,\mathrm{d}z=2\pi\mathrm{i}(A+B).
\]
这里每一项都由已建立的积分公式或端点公式计算，无需关于一般闭曲线的计数公式。

若 \(h\) 在 \(U\) 上有全局原函数，沿前两个小圆的积分都必须为零，故 \(A=B=0\)。反过来，若 \(A=B=0\)，则
\[
  H(z)=-\frac C{z-1}\quad(z\in U)
\]
在 \(U\) 上单值全纯，且 \(H'=h\)。因此充要条件正是 \(A=B=0\)，与 \(C\) 的值无关。

例如取 \(A=1\)、\(B=-1\)、\(C=0\)。外圆积分虽然为零，围绕零点的小圆积分却为 \(2\pi\mathrm{i}\)，所以不存在全局原函数。两个孔洞分别产生的积分在外圆计算中相消，外圆只检验了它们的和。
\end{solution}

\Needspace{6\baselineskip}
\begin{exercise}\label{b3:ex:02-6}
设 \(f\) 在闭单位圆盘 \(\overline{\mathbb D}\) 的某个开邻域上全纯，满足 \(f(0)=f'(0)=0\)，并且
\[
  |f(z)|\le|z^2(2+z)|\quad(|z|=1).
\]
证明这一不等式在整个闭单位圆盘上成立。刻画在某个 \(z_0\in\mathbb D\setminus\{0\}\) 取等的全部函数。最后构造反例，说明只保留 \(f(0)=0\) 而删去 \(f'(0)=0\) 时，边界不等式不足以推出内部不等式。
\end{exercise}
\begin{solution}
由原点的 Taylor 展开及两个消失条件，可以在零点附近写
\(f(z)=z^2u(z)\)，其中 \(u\) 全纯且 \(u(0)=f''(0)/2\)。在非零点处令 \(u(z)=f(z)/z^2\)，这个商与零点附近的幂级数一致，所以 \(u\) 在闭圆盘的一个邻域上全纯。

函数 \(2+z\) 在闭单位圆盘上不为零，在略小的邻域上也不为零。因此
\[
  q(z)=\frac{u(z)}{2+z}
\]
在闭圆盘的邻域上全纯，并在非零点处等于 \(f(z)/(z^2(2+z))\)。边界不等式给出 \(|q(z)|\le1\)；最大模原理遂使 \(|q(z)|\le1\) 于整个闭圆盘。乘回分母可得
\[
  |f(z)|\le|z^2(2+z)|.
\]
在 \(z=0\) 处，两边都为零，所以这一点也已包括在结论中。

若某个内部非零点 \(z_0\) 取等，因为分母在那里不为零，就有 \(|q(z_0)|=1\)。最大模原理的等号情形使 \(q=c\)，其中 \(|c|=1\)，从而
\[
  f(z)=c z^2(2+z)\quad(z\in\mathbb D).
\]
连续性使该等式也延续到边界。反过来，这些函数在所有点都取等。要求取等点非零不能省略，因为原点处 \(0=0\) 对全部题设函数成立，不提供任何刚性信息。

若只要求 \(f(0)=0\)，取 \(f(z)=z(2+z)\)。它在单位圆上满足与比较函数完全相同的模，但对 \(0<|z|<1\)，
\[
  \frac{|f(z)|}{|z^2(2+z)|}=\frac1{|z|}>1.
\]
此例与比较函数在圆盘内有相同的零点集合，却没有相同的最低消失阶数。补商所需的是零点重数条件，仅有共同零点不够。
\end{solution}

\Needspace{6\baselineskip}
\begin{exercise}\label{b3:ex:02-7}
设 \(m\ge1\) 为整数，\(f:\mathbb D\rightarrow\mathbb D\) 全纯，并满足
\[
  f(0)=f'(0)=\cdots=f^{(m-1)}(0)=0.
\]
证明
\[
  |f(z)|\le|z|^m\quad(z\in\mathbb D),
  \quad |f^{(m)}(0)|\le m!.
\]
刻画在某个非零点满足 \(|f(z)|=|z|^m\)，或满足 \(|f^{(m)}(0)|=m!\) 时的全部函数。说明系数界中的每个允许值都能出现。\(m=2\) 时即为具有二阶消失条件的情形。
\end{exercise}
\begin{solution}
原点的 Taylor 展开说明，对每个 \(0\le j\le m\)，函数
\[
  q_j(z)=\frac{f(z)}{z^j}\quad(z\ne0)
\]
都可在零点全纯补定义。这里 \(q_0=f\)，并且
\[
  q_j(0)=0\quad(0\le j<m),
  \quad q_m(0)=\frac{f^{(m)}(0)}{m!}.
\]
下面逐次使用 Schwarz 引理，同时检查每一步的值域。初始的 \(q_0\) 映入开单位圆盘且在零点为零，所以 Schwarz 引理给出 \(|q_1|\le1\)。若 \(1<m\)，则 \(q_1(0)=0\)；假使 \(q_1\) 在内部某点达到模 \(1\)，最大模原理将使它成为模为 \(1\) 的常数，与零点值矛盾。因此此时 \(q_1\) 实际映入开单位圆盘，可以再次使用 Schwarz 引理。

一般地，当 \(j<m\) 时，同一论证使 \(q_j\) 映入开单位圆盘，并由 \(q_j(0)=0\) 推出 \(|q_{j+1}|\le1\)。进行有限次后得到 \(|q_m|\le1\)。对 \(z\ne0\) 乘以 \(|z|^m\)，再在原点直接取值，得到第一项估计；在 \(q_m(0)\) 处取值则得到导数界。

若非零点的函数值估计取等，或中心的导数估计取等，就有 \(|q_m|\) 在圆盘内某点等于 \(1\)。最大模原理使 \(q_m=c\)，其中 \(|c|=1\)，所以
\[
  f(z)=c z^m.
\]
这类函数确实把开单位圆盘映入自身，并使两项估计均取等。

对于任意复数 \(b\) 满足 \(|b|\le m!\)，取
\[
  f_b(z)=\frac b{m!}z^m.
\]
它满足全部消失条件，且在 \(|z|<1\) 上始终有 \(|f_b(z)|<1\)，包括 \(|b|=m!\) 的端点情形。又 \(f_b^{(m)}(0)=b\)，故导数界中没有遗漏不能达到的值。
\end{solution}

\Needspace{6\baselineskip}
\begin{exercise}\label{b3:ex:02-8}
给定实数 \(a\in(0,1)\)，构造处处非零的全纯函数 \(f:\mathbb D\rightarrow\mathbb D\)，使
\[
  f(0)=a,\quad\inf_{z\in\mathbb D}|f(z)|=0.
\]
对你构造的函数，求出任意固定的 \(0<r<1\) 所对应的
\(\min_{|z|\le r}|f(z)|\) 与 \(\max_{|z|\le r}|f(z)|\)，并指出取等点。解释为什么全域没有正下界不违反最小模原理，也不违反每个紧子圆盘上连续非零函数有正最小值的结论。
\end{exercise}
\begin{solution}
要让指数的模变得任意小而又始终非零，可以让指数的实部趋于负无穷。取
\[
  f(z)=\exp\left(\log a\cdot\frac{1+z}{1-z}\right),
  \quad z\in\mathbb D,
\]
其中 \(\log a<0\) 是实对数。分母在开单位圆盘中不为零，因此 \(f\) 全纯；指数函数没有零点，所以 \(f\) 也没有零点。直接计算
\[
  \operatorname{Re}\frac{1+z}{1-z}
  =\frac{1-|z|^2}{|1-z|^2}>0,
\]
于是
\[
  |f(z)|=a^{(1-|z|^2)/|1-z|^2}\in(0,1),
  \quad f(0)=a.
\]
沿实半径 \(z=r\uparrow1\)，指数 \((1+r)/(1-r)\to+\infty\)，故
\(|f(r)|=a^{(1+r)/(1-r)}\to0\)。这证明全域模的下确界为零，而函数在任何内点都不取零值。

为计算闭子圆盘上的极值，写 \(z=\rho e^{\mathrm{i}t}\)、\(0\le\rho\le r\)。由
\(1-\rho\le|1-z|\le1+\rho\)，得到
\[
  \frac{1-\rho}{1+\rho}
  \le\frac{1-|z|^2}{|1-z|^2}
  \le\frac{1+\rho}{1-\rho}.
\]
左端的函数随 \(\rho\) 增大而减小，右端的函数随 \(\rho\) 增大而增大，因而在 \(\rho\le r\) 时
\[
  \frac{1-r}{1+r}
  \le\frac{1-|z|^2}{|1-z|^2}
  \le\frac{1+r}{1-r}.
\]
两端分别在 \(z=-r\) 与 \(z=r\) 达到。由于 \(a^s\) 关于实变量 \(s\) 严格递减，模的极值次序反过来，最终有
\[
  \min_{|z|\le r}|f(z)|=a^{(1+r)/(1-r)},
  \quad\max_{|z|\le r}|f(z)|=a^{(1-r)/(1+r)}.
\]
当 \(0<r<1\) 时，最小值在 \(z=r\) 取得，最大值在 \(z=-r\) 取得；两个点都位于所选子圆盘的边界。

最小模原理排除的是无零非常值全纯函数在其定义域内部取得局部最小模。本例在整个 \(\mathbb D\) 上的下确界 \(0\) 没有取得，闭子圆盘的最小点也只是该子圆盘的边界点，继续沿实半径前进会使模更小。每个固定 \(r<1\) 的最小值都严格为正，但这些正数随着 \(r\uparrow1\) 趋于零；紧子集上的正下界可以依赖这个紧子集。
\end{solution}

% !TeX root = ../../Book3.tex
\chapter{Laurent 级数、奇点与留数}
\label{b3:ch:03}
\BookChapterNumber{b3:ch:03}
\begin{chapterabstract}
Laurent 展开把环域中的全纯函数分成正幂与主部，并由主部分类孤立奇点。留数提取积分中的唯一局部贡献，辐角原理和 Rouché 定理则把围道积分变成稳定的零点计数。本章同时区分函数在无穷远的奇点与微分在无穷远的留数。
\end{chapterabstract}

\begin{learninggoals}
\begin{itemize}
\item\label{b3:item:03-goal-1}能在不同环域选择 Laurent 展开并求出系数。
\item\label{b3:item:03-goal-2}能用局部有界性、主部和倒数函数辨认三类孤立奇点。
\item\label{b3:item:03-goal-3}能计算有限点与无穷远留数，并证明实积分附加弧的极限。
\item\label{b3:item:03-goal-4}能以重数使用辐角原理与 Rouché 定理，并解释零点稳定性。
\end{itemize}
\end{learninggoals}

函数 \(1/(z(z-1))\) 在 \(0<|z|<1\) 与 \(|z|>1\) 有不同的展开。负幂并不是错误项，而是中心孔洞留下的信息；其中 \(-1\) 次系数又恰好决定绕中心的积分。这种局部展开与全局计数的关系，是本章从级数走向留数的原因。

\ProseParagraphBreak

本章需要\chapref{b3:ch:02}的积分公式与估计，以及\chapref{b3:ch:01}的孤立零点分解。John B. Conway 的《Functions of One Complex Variable I》\cite{Conway1978Functions}和 Elias M. Stein、Rami Shakarchi 的《Complex Analysis》\cite{Stein2003Complex}涵盖 Laurent 与留数主题，可用于延伸阅读；下面的奇点判别与计数推导在正文展开。

% !TeX root = ../../Book3.tex
\section{Laurent 级数}
\label{b3:sec:0301}

Taylor 级数要求中心点也在定义域内。去掉中心后，负幂可能出现，但它们并不破坏环域内部的一致收敛。把内外两条边界上的 Cauchy 核分别展开，就能同时得到正幂与负幂；两个部分有不同的几何来源。

\subsection{环域}
\label{b3:subsec:030101}

设 \(0\le r<R\le\infty\)，记
\[
 A(a;r,R)\coloneqq\{\,z\in\mathbb C\mid r<|z-a|<R\,\}.
\]
它称为以 \(a\) 为中心的\term[huanyu]{环域}[annulus]；\(r=0\) 时是去心圆盘。所谓 Laurent 级数的收敛，指正幂部分与负幂部分分别收敛，不能用正负指标混合的任意排列改变和。

\ProseParagraphBreak

负幂级数 \(\sum_{n\ge1}c_{-n}(z-a)^{-n}\) 在倒数变量 \(w=1/(z-a)\) 中是没有常数项的幂级数。若该幂级数半径为 \(S\)，它在 \(|z-a|>1/S\) 上收敛；正幂部分则在 \(|z-a|<R\) 上收敛。两者的公共范围因而天然是环域，可能为空。边界圆上的收敛须分别考察。

\ProseParagraphBreak

例如 \(1/(z(z-1))\) 在 \(0<|z|<1\) 时写成 \(-\sum_{n\ge0}z^{n-1}\)，在 \(|z|>1\) 时写成 \(\sum_{n\ge0}z^{-n-2}\)。展开中心相同，而穿过 \(|z|=1\) 时表达式改变，原因是圆周上的奇点 \(z=1\) 阻断了统一展开。

\ProseParagraphBreak

选择展开环域时，可以先将所有与中心的距离画在数轴上。例如函数 \(1/((z-1)(z-2))\) 的奇点到中心0的距离为1与2，因此有三个候选范围：\(|z|<1\)、\(1<|z|<2\) 和 \(|z|>2\)。中间范围中的 \(1/(z-1)\) 应展开成负幂，而 \(1/(z-2)\) 应展开成正幂。同一个有理式里既有正幂又有负幂，并不表示使用了互相矛盾的几何级数；它们的公共收敛域正好是指定环域。

\ProseParagraphBreak

Laurent 主部是否有限，也取决于中心在哪里。函数 \(1/(z-1)\) 在中心0的小圆盘是 Taylor 级数，在中心1则是只有一个负幂的 Laurent 级数，在无穷远坐标又具有可去奇点。分类和系数总要附一个中心，不能只说“这个函数有负幂”而略去展开范围。

\ProseParagraphBreak

若两个奇点到中心距离相同，例如 \(1/(z^2+1)\)，它们在同一分界圆上。尽管圆周上只有两个点坏掉，也不能把圆内与圆外的 Laurent 展开拼成一个穿过该圆的同心环域展开，因为一个 Laurent 级数在其开放收敛环域上处处全纯，会把这两个真极点也包括进去。

\subsection{Laurent 展开}
\label{b3:subsec:030102}

\term[laurent-jishu]{Laurent 级数}[Laurent series]允许整数幂指标。
\begin{theorem}[Laurent]\label{b3:thm:03-laurent}
若 \(f\) 在 \(A(a;r,R)\) 上全纯，则有唯一展开
\[
 f(z)=\sum_{n=-\infty}^{\infty}c_n(z-a)^n,
\]
其中正幂部分在 \(|z-a|<R\) 全纯，负幂部分在 \(|z-a|>r\) 全纯；级数在每个闭子环域上一致绝对收敛。
\end{theorem}
\begin{proof}
固定 \(r<\rho_1<|z-a|<\rho_2<R\)。沿环域的一条径向线切开，将内外圆的边界用两条相反方向的切口边连接；切口可以选择不经过 \(z\)。对切开区域使用 Cauchy 公式，再让两条切口边靠拢。核与函数在切口附近连续，两个切口积分相消，得到
\[
 f(z)=\frac1{2\pi\mathrm{i}}\int_{|\zeta-a|=\rho_2}\frac{f(\zeta)}{\zeta-z}\,\mathrm{d}\zeta-
 \frac1{2\pi\mathrm{i}}\int_{|\zeta-a|=\rho_1}\frac{f(\zeta)}{\zeta-z}\,\mathrm{d}\zeta,
\]
两圆均取逆时针方向。在外圆上，展开 \((\zeta-z)^{-1}=\sum_{n\ge0}(z-a)^n/(\zeta-a)^{n+1}\)；在内圆上，展开 \((\zeta-z)^{-1}=-\sum_{n\ge0}(\zeta-a)^n/(z-a)^{n+1}\)。若 \(z\) 限于严格夹在两圆之间的闭子环域，两个几何级数都一致绝对收敛，故可以逐项积分。

所得系数为 \(c_n=(2\pi\mathrm{i})^{-1}\int_{|\zeta-a|=\rho}f(\zeta)/(\zeta-a)^{n+1}\,\mathrm{d}\zeta\)。在两个容许半径之间，被积函数全纯；对其环形边界切开并用 Cauchy 定理，两个圆周积分相等。因而系数不依赖所选半径。取半径分别靠近 \(r\) 和 \(R\)，前述表示覆盖整个环域。对任意闭子环域，选更内、更外两圆，几何估计给出一致绝对收敛。

若另有一组 Laurent 系数，在某个中间圆周上逐项积分乘以 \((z-a)^{-k-1}\)，只有指标 \(k\) 的项留下 \(2\pi\mathrm{i}\)。于是两组系数逐个相等。
\end{proof}
这里切口不是对函数选分支，而是把有孔区域分解成可以使用 Cauchy 公式的区域。切口两侧的函数本来相同，故会相消；对多值表达式不能直接照用这一步。

\subsection{Laurent 系数}
\label{b3:subsec:030103}

\begin{proposition}[Laurent 系数公式]\label{b3:prop:03-laurent-coefficients}
在 Laurent 定理的条件下，对任意 \(r<\rho<R\) 与整数 \(n\)，
\[
 c_n=\frac1{2\pi\mathrm{i}}\int_{|z-a|=\rho}\frac{f(z)}{(z-a)^{n+1}}\,\mathrm{d}z,
 \qquad |c_n|\le M_\rho\rho^{-n},
\]
其中 \(M_\rho=\max_{|z-a|=\rho}|f(z)|\)。
\end{proposition}
\begin{proof}
公式已由前一定理的构造与唯一性得到。对积分使用长度估计，圆长 \(2\pi\rho\) 与分母 \(\rho^{n+1}\) 相消，即得对正负指标同时有效的估计。
\end{proof}
当 \(n=-m<0\) 时，估计成为 \(|c_{-m}|\le M_\rho\rho^m\)。若函数靠近中心有界，令 \(\rho\downarrow0\) 就迫使全部负幂系数为零。正幂系数由外侧增长控制，负幂系数由内侧增长控制；同一公式的两个方向将给出奇点分类。

\ProseParagraphBreak

级数仍可在闭子环域上逐项微分，因为两部分分别是在原变量与倒数变量中的幂级数；微分后 \((z-a)^n\) 产生 \(n(z-a)^{n-1}\)，包括负指标。逐项原函数则需额外检查 \(n=-1\) 项，它不能在整个环域上取单值原函数。

\ProseParagraphBreak

Laurent 展开中的系数可由对称性提前筛选。若 \(f(-z)=f(z)\) 且定义环域以0为中心，唯一性给 \(c_n=(-1)^n c_n\)，所以奇数指标系数为零，包括负的奇数指标。若 \(f(\overline z)=\overline{f(z)}\)，则把展开共轭并替换变量后，由唯一性得到全部系数为实数。这里的对称必须在整个环域成立；只在实轴知道函数值实，并且有内部聚点时，也可由恒等定理传播，但要先保证所比较的共轭反射函数全纯。

\ProseParagraphBreak

负幂项不能用普通原函数遗漏。若 \(f\) 在同心环域的留数为0，可以对全部其他项逐项积分，得到在该环域收敛的 Laurent 原函数。其正幂与负幂部分仍在闭子环域一致收敛，可用幂级数估计核对。若留数非零，绕中心圆周积分非零，端点公式立即排除原函数。因此对同心环域上的全纯函数，留数为零恰是存在单值原函数的条件；对于有多个孔洞的区域，要分别检查各孔的周期。

\begin{proposition}\label{b3:prop:03-annulus-primitive}
设 \(f\) 在环域 \(r<|z-a|<R\) 全纯，Laurent 系数为 \(c_n\)。它在该环域有单值全纯原函数，当且仅当 \(c_{-1}=0\)。
\end{proposition}
\begin{proof}
有原函数时任意中间圆周的积分为零，所以 \(2\pi\mathrm{i}c_{-1}=0\)。反之，去掉 \(n=-1\) 项后令
\[
 F(z)=\sum_{n\ne-1}\frac{c_n}{n+1}(z-a)^{n+1}.
\]
在任意闭子环域上，正幂与负幂部分均正规收敛；略扩大子环域后用几何估计可同时控制导数级数。因此逐项求导合法，且 \(F'=f\)。
\end{proof}
这里每个负幂项 \((z-a)^{-n}\)、\(n\ge2\) 都有原函数；只有 \((z-a)^{-1}\) 会产生对数的多值性。故“奇点存在”并不等于“没有原函数”。\(1/z^2\) 虽在0有极点，却在去心平面有原函数 \(-1/z\)。闭路积分测量的是留数这一特定系数，不能直接读出全部主部。

% !TeX root = ../../Book3.tex
\section{孤立奇点}
\label{b3:sec:0302}

奇点分类要先确认中心确实被一个完整的去心圆盘隔离。设 \(a\in\mathbb C\)、\(R>0\)，若 \(f\) 在 \(0<|z-a|<R\) 上全纯，而暂未在 \(a\) 定义，就称 \(a\) 为 \(f\) 的\term[guli-qidian]{孤立奇点}[isolated singularity]。若其他奇点不断积聚到 \(a\)，便不存在这样的去心圆盘，不能将它们全部当作中心的一个孤立奇点。

\ProseParagraphBreak

Laurent 展开把问题变成负幂系数的判别。没有负幂时，正幂级数会填入中心；只有有限个非零负幂时，乘上适当的中心距离幂就能消去分母；负幂无限多时，这两种修补都不再有效。下面分别把这三种情形转成函数值及局部因子的条件。

\subsection{可去奇点}
\label{b3:subsec:030201}

\Needspace{6\baselineskip}
\begin{definition}\label{b3:def:03-removable}
设 \(a\in\mathbb C\)、\(R>0\)，\(f:D(a,R)\setminus\{a\}\to\mathbb C\) 全纯。若存在全纯函数 \(\widetilde f:D(a,R)\to\mathbb C\)，满足 \(\widetilde f(z)=f(z)\) 对所有 \(z\ne a\) 成立，则称 \(a\) 为 \(f\) 的\term[kequ-qidian]{可去奇点}[removable singularity]。
\end{definition}

延拓只能在中心补入一个值，连续性迫使这个值等于 \(\lim_{z\to a}f(z)\)，所以若延拓存在，它必唯一。真正的判别不必事先知道这个极限，只须给出整个去心邻域的上界。

\Needspace{6\baselineskip}
\begin{theorem}[可去奇点判别]\label{b3:thm:03-removable}
设 \(a\in\mathbb C\)、\(R>0\)，\(f:D(a,R)\setminus\{a\}\to\mathbb C\) 全纯。则以下条件等价：\(a\) 是可去奇点；以 \(a\) 为中心的 Laurent 展开没有非零负幂系数；存在 \(0<r<R\)、\(M<\infty\)，使 \(|f(z)|\le M\) 对 \(0<|z-a|<r\) 成立；有限极限 \(\lim_{z\to a}f(z)\) 存在。成立时，延拓在 \(a\) 处的值为 Laurent 系数 \(c_0\)。
\end{theorem}
\begin{proof}
若 \(f\) 可去，延拓函数在 \(D(a,R)\) 内的 Taylor 展开，在去心圆盘上也是原函数的 Laurent 展开。唯一性说明所有负幂系数为零。反之，若负幂系数全为零，Laurent 展开只有正幂部分，它在整个 \(D(a,R)\) 全纯，中心值为 \(c_0\)，直接给出延拓。

若存在题设的局部上界，则对每个 \(m\ge1\) 及 \(0<\rho<r\)，\cref{b3:prop:03-laurent-coefficients}给出
\[
 |c_{-m}|\le M\rho^m.
\]
系数与半径无关，令 \(\rho\downarrow0\) 得 \(c_{-m}=0\)。这对所有 \(m\) 成立，因此有界性推出可去。可去时，延拓在一个较小闭圆盘上连续，故局部有界。最后，可去蕴含有限极限存在；有限极限又使函数在充分小的去心圆盘内有界。四个条件由此等价。
\end{proof}

例如在零点以外考虑
\[
 f(z)=\frac{e^z-1-z}{z^2}.
\]
指数级数给出 \(f(z)=1/2+z/6+z^2/24+\cdots\)，这个幂级数在零点也收敛，所以补入 \(f(0)=1/2\) 后全纯。若只沿某一条路径看极限，结论就不可靠：\(e^{-1/z^2}\) 沿实轴趋于零，但在 \(z=\mathrm{i}y\) 时等于 \(e^{1/y^2}\)，沿虚轴无界。可去判别中的上界必须同时覆盖所有充分接近中心的方向。

\ProseParagraphBreak

有时表达式本身的上界不易取得，却知道乘以 \(z-a\) 以后比常数还小。这也足以排除全部负幂。

\Needspace{6\baselineskip}
\begin{corollary}\label{b3:cor:03-removable-small-growth}
设 \(a\in\mathbb C\)、\(R>0\)，\(f:D(a,R)\setminus\{a\}\to\mathbb C\) 全纯。若
\[
 \lim_{z\to a}(z-a)f(z)=0,
\]
则 \(a\) 是 \(f\) 的可去奇点。
\end{corollary}
\begin{proof}
令 \(g(z)=(z-a)f(z)\)。有限极限使 \(g\) 可去，并使延拓在 \(a\) 处取值零。它的 Taylor 展开因而没有常数项，可写成 \(g(z)=(z-a)h(z)\)，其中 \(h\) 在整个 \(D(a,R)\) 全纯。在 \(z\ne a\) 时 \(h(z)=f(z)\)，所以 \(h\) 是所需延拓。
\end{proof}

这里极限必须等于零。函数 \(1/(z-a)\) 满足 \((z-a)f(z)=1\)，却不能补成连续函数。乘去一个分母后得到非零极限，描述的正是下一小节的一阶极点。

\subsection{极点}
\label{b3:subsec:030202}

极点可以看作全纯函数零点的倒数：去掉一个有限整数次的分母后，剩余因子在中心全纯且非零。这一非零条件确定了真正的阶数。

\Needspace{6\baselineskip}
\begin{definition}\label{b3:def:03-pole}
设 \(a\in\mathbb C\)、\(R>0\)，\(f:D(a,R)\setminus\{a\}\to\mathbb C\) 全纯。若存在整数 \(m\ge1\)、\(0<r<R\) 以及全纯函数 \(g:D(a,r)\to\mathbb C\)，满足 \(g(a)\ne0\) 且
\[
 f(z)=(z-a)^{-m}g(z)\quad(0<|z-a|<r),
\]
则称 \(a\) 为 \(f\) 的 \(m\) 阶\term[jidian]{极点}[pole]。
\end{definition}

\Needspace{6\baselineskip}
\begin{proposition}\label{b3:prop:03-pole-criterion}
设 \(a\in\mathbb C\)、\(R>0\)，\(f:D(a,R)\setminus\{a\}\to\mathbb C\) 全纯。则 \(a\) 是极点，当且仅当 \(|f(z)|\to\infty\)（\(z\to a\)）；也当且仅当 Laurent 展开只有有限个非零负幂系数且至少有一个。若最低非零指数为 \(-m\)，则极点阶数恰为 \(m\)，并有
\[
 \lim_{z\to a}(z-a)^m f(z)=c_{-m}\ne0.
\]
反之，若对某个正整数 \(m\)，上式左端存在且为非零有限值，则 \(a\) 是 \(m\) 阶极点。
\end{proposition}
\begin{proof}
若 \(f(z)=(z-a)^{-m}g(z)\)、\(g(a)\ne0\)，连续性允许缩小半径，使
\[
 \frac{|g(a)|}{2}\le|g(z)|\le\frac{3|g(a)|}{2}.
\]
因此
\[
 \frac{|g(a)|}{2|z-a|^m}\le|f(z)|
 \le\frac{3|g(a)|}{2|z-a|^m},
\]
从而 \(|f(z)|\to\infty\)。将 \(g\) 的 Taylor 展开除以 \((z-a)^m\)，最低非零指数是 \(-m\)，系数为 \(g(a)\)，其他负幂只可能有有限个，并得到所述极限。

反之，若 \(|f(z)|\to\infty\)，则在一个充分小的去心圆盘内有 \(|f(z)|>1\)，所以 \(f\) 无零点，\(h=1/f\) 全纯并趋于零。可去奇点判别使 \(h\) 延拓到中心，且 \(h(a)=0\)。它在去心圆盘上处处非零，所以不恒为零。由\cref{b3:prop:01-zero-factorization}，缩小圆盘后可写
\[
 h(z)=(z-a)^m H(z),\quad m\ge1,\quad H(a)\ne0.
\]
进一步缩小圆盘使 \(H\) 无零点，取倒数便得
\(f(z)=(z-a)^{-m}/H(z)\)，这是极点分解。

若只有有限个非零负幂系数，令 \(-m\) 为其中最低指数，则 \((z-a)^m f(z)\) 的展开只有非负幂，并在中心取非零值 \(c_{-m}\)，也给出极点分解。最后，若 \((z-a)^m f(z)\) 有非零有限极限，它由可去奇点判别延拓为中心非零的全纯函数，仍得到同一分解。最低非零指数的唯一性同时保证极点阶数唯一。
\end{proof}

这里 \(|f(z)|\to\infty\) 指对每个 \(M>0\)，都有 \(\delta>0\)，使 \(0<|z-a|<\delta\) 时 \(|f(z)|>M\)。仅仅“在任意去心邻域无界”弱得多，不能代替这一条件。本性奇点也可能在一列点上任意大，却同时在另一列点上保持有限。

\ProseParagraphBreak

商式的分母为零时，应先比较分子与分母的零点阶数。设 \(g,h\) 在 \(a\) 附近全纯且都不局部恒零，令它们在 \(a\) 的阶数分别为 \(p,q\ge0\)，其中中心非零时阶数取零。零点分解给出
\[
 g(z)=(z-a)^pG(z),\quad h(z)=(z-a)^qH(z),
 \quad G(a)H(a)\ne0.
\]
在缩小的邻域内 \(H\) 无零点，故
\[
 \frac{g(z)}{h(z)}=(z-a)^{p-q}\frac{G(z)}{H(z)}.
\]
当 \(p>q\) 时，商延拓后有 \(p-q\) 阶零点；当 \(p=q\) 时，商的可去值为 \(G(a)/H(a)\ne0\)；当 \(p<q\) 时，商有 \(q-p\) 阶极点。分母的形式阶数只有在分子不消失时才直接等于极点阶数。

\ProseParagraphBreak

例如 \(1-\cos z=z^2/2+O(z^4)\)、\(\sin z=z+O(z^3)\)，所以
\[
 \frac{1-\cos z}{z\sin z}
 =\frac{1/2+O(z^2)}{1+O(z^2)}
\]
在零点可去，补入值为 \(1/2\)。而
\[
 \frac{1-\cos z}{(\sin z)^3}
 =\frac1z\frac{1/2+O(z^2)}{1+O(z^2)}
\]
在零点只有一阶极点，尽管分母具有三阶零点。这里各个 \(O(z^2)\) 都来自完整邻域内的收敛 Taylor 级数，分母中的剩余因子在中心等于一，故可以在小圆盘内求全纯倒数。

\ProseParagraphBreak

在开集 \(U\subseteq\mathbb C\) 上，若函数除一组在 \(U\) 内没有聚点的极点外全纯，并在各极点赋值 \(\infty\)，就称它为\term[ya-chun-hanshu]{亚纯函数}[meromorphic function]。极点分解说明它局部是两个全纯函数的商，分母不局部恒零；反过来，上面的公共零点消去也说明，这样的商补入可去值后局部只有极点。例如 \(\cot z=\cos z/\sin z\) 的分母在 \(\pi\mathbb Z\) 为单零点，分子在那里不为零，故这些点都是一阶极点。

\ProseParagraphBreak

对不局部恒零的亚纯函数，可以把阶数统一约定为：零点阶数取正，极点阶数取负，非零全纯点的阶数取零。局部分解立即给出乘积阶数相加、商的阶数相减。求和时最低项可能抵消，例如
\[
 \frac1z+\left(-\frac1z+z\right)=z.
\]
两项各有极点，和却有零点。只有比较首个非零系数以后，才能判断这种抵消是否发生。

\subsection{本性奇点}
\label{b3:subsec:030203}

\Needspace{6\baselineskip}
\begin{definition}\label{b3:def:03-essential}
设 \(a\in\mathbb C\)、\(R>0\)，\(f:D(a,R)\setminus\{a\}\to\mathbb C\) 全纯。若 \(a\) 既不是可去奇点也不是极点，则称它为 \(f\) 的\term[benxing-qidian]{本性奇点}[essential singularity]。
\end{definition}

按前两项判别，本性奇点恰好对应无穷多个非零负幂系数。三类因此互斥且穷尽：主部或者全为零，或者只有有限个非零项，或者有无穷多个非零项。中心是否具有某条路径上的极限，不在这三种代数条件中。

\ProseParagraphBreak

例如
\[
 e^{1/z}=\sum_{n=0}^{\infty}\frac{z^{-n}}{n!}
 \quad(z\ne0)
\]
的负幂无限多，所以零是本性奇点。沿正实轴趋近零时函数趋于无穷，沿负实轴则趋于零。沿虚轴也没有单一极限：取
\[
 z_k=\frac1{2\pi\mathrm{i}k},
 \quad w_k=\frac1{(2k+1)\pi\mathrm{i}}
 \quad(k\ge1),
\]
两列都趋于零，却分别满足 \(e^{1/z_k}=1\)、\(e^{1/w_k}=-1\)。而指数函数没有零点，所以值零始终不被取得。一个本性奇点附近可以同时出现增长、衰减和振荡。

\ProseParagraphBreak

孤立奇点分类的前提是完整去心圆盘上已经有一个单值全纯函数。平方根在零附近的困难不满足这一前提。若存在全纯函数 \(g\) 定义在整个 \(0<|z|<R\) 上并满足 \(g(z)^2=z\)，则 \(|g(z)|=|z|^{1/2}\)，可去奇点判别使它延拓到零，且 \(g(0)=0\)。延拓不恒为零，故其零点阶数是某个正整数 \(m\)，那么 \(g^2\) 的零点阶数为 \(2m\)，不可能等于 \(z\) 的阶数一。这证明全去心圆盘上不存在这样的分支；在割去一条射线的区域选取平方根，不与此矛盾。

\subsection{Casorati–Weierstrass 定理}
\label{b3:subsec:030204}

\term[casorati-weierstrass-dingli]{Casorati--Weierstrass 定理}[Casorati--Weierstrass theorem]说明本性奇点附近的像不能避开任何开圆盘。证明的构造来自一个反问：如果函数始终离某个值有正距离，那么差的倒数就有界，可以用可去奇点判别处理。

\Needspace{6\baselineskip}
\begin{theorem}[Casorati--Weierstrass]\label{b3:thm:03-casorati-weierstrass}
设 \(a\in\mathbb C\)、\(R>0\)，\(f:D(a,R)\setminus\{a\}\to\mathbb C\) 全纯，且 \(a\) 是本性奇点。则对每个 \(0<r<R\)，集合 \(f(D(a,r)\setminus\{a\})\) 在 \(\mathbb C\) 中稠密。等价地，对任意 \(b\in\mathbb C\)、\(\varepsilon>0\) 和 \(0<r<R\)，都存在 \(z\) 满足
\[
 0<|z-a|<r,\quad |f(z)-b|<\varepsilon.
\]
\end{theorem}
\begin{proof}
若对某个 \(0<r<R\) 像集不稠密，其闭包的补集含有一个开圆盘。因此存在 \(b\in\mathbb C\)、\(\varepsilon>0\)，使
\[
 |f(z)-b|\ge\varepsilon\quad(0<|z-a|<r).
\]
令 \(g(z)=1/(f(z)-b)\)。分母无零点使 \(g\) 全纯，并且 \(|g(z)|\le1/\varepsilon\)，所以 \(g\) 可以全纯延拓到 \(a\)。

若延拓满足 \(g(a)\ne0\)，连续性保证在一个小圆盘内 \(g\) 始终非零，于是 \(b+1/g\) 在整个小圆盘全纯，并在去心部分等于 \(f\)。这使 \(a\) 可去，与假设矛盾。

若 \(g(a)=0\)，则 \(g\) 在去心圆盘上非零，因而不局部恒零。零点分解给出
\[
 g(z)=(z-a)^m h(z),\quad m\ge1,\quad h(a)\ne0.
\]
缩小圆盘使 \(h\) 无零点，再写
\[
 f(z)=b+\frac1{(z-a)^m h(z)}
 =(z-a)^{-m}\left(b(z-a)^m+\frac1{h(z)}\right).
\]
括号内是全纯函数，在中心取非零值 \(1/h(a)\)，所以 \(a\) 是 \(m\) 阶极点，也与假设矛盾。两个可能的延拓值都被排除，故像集必须稠密。
\end{proof}

定理同时允许缩小自变量邻域和目标误差。对任意给定 \(b\in\mathbb C\)，取 \(r_k=\min\{R/2,1/k\}\)、\(\varepsilon_k=1/k\)，即可选择 \(z_k\ne a\)，使 \(z_k\to a\) 且 \(f(z_k)\to b\)。它并不保证某一点真正满足 \(f(z)=b\)，例如 \(e^{1/z}\) 就遗漏零值。\chapref{b3:ch:14}的 Picard 大定理将进一步讨论精确取值，而这里的证明只给稠密性。

\ProseParagraphBreak

任意去心邻域的像既然稠密，其中的模就没有有限上界，模的下确界又必为零。这两个结论都针对整个邻域，不能强加到每一条趋近中心的路径上。稠密性给出的离散点列，也不能自动连成一条使函数值收敛的连续路径：连接两点时，函数的像可能离开指定小圆盘。点列聚值、沿路径的极限和精确取值是不同的问题。

\ProseParagraphBreak

无穷远处的孤立奇点通过换坐标归入上述分类。设 \(R_0>0\)，\(f\) 在 \(|z|>R_0\) 全纯，令
\[
 w=\frac1z,\quad F(w)=f(1/w)
 \quad(0<|w|<1/R_0).
\]
约定 \(f\) 在无穷远可去、有 \(m\) 阶极点或本性，分别是指 \(F\) 在零点具有对应类型。这里分类的是函数值；本章下一节在变换积分微分时，还会出现换元导数，不能把两种计算混为一谈。

\Needspace{6\baselineskip}
\begin{proposition}[整函数在无穷远的分类]\label{b3:prop:03-entire-at-infinity}
设 \(f:\mathbb C\to\mathbb C\) 为整函数，通过 \(F(w)=f(1/w)\) 定义它在无穷远处的奇点类型。则无穷远可去当且仅当 \(f\) 为常函数；无穷远为 \(m\) 阶极点当且仅当 \(f\) 是次数恰为 \(m\ge1\) 的多项式；无穷远为本性奇点当且仅当 \(f\) 不是多项式。
\end{proposition}
\begin{proof}
若 \(F\) 在零点可去，它在小圆盘内有界，因此 \(f\) 在某个大圆之外有界。\(f\) 在该大圆所围闭圆盘上连续，也有界，合并两项上界便得到全平面有界。由\cref{b3:thm:02-liouville}，\(f\) 为常函数。反方向直接换坐标即可。

若 \(F\) 在零点有 \(m\) 阶极点，局部分解给出小 \(|w|\) 时 \(|F(w)|\le C|w|^{-m}\)。所以大 \(|z|\) 时有 \(|f(z)|\le C|z|^m\)；再利用紧圆盘上的有界性，增加常数后得到全平面的估计
\[
 |f(z)|\le C'(1+|z|)^m.
\]
\Cref{b3:prop:02-polynomial-growth}说明 \(f\) 是次数不超过 \(m\) 的多项式。它不能是常函数，否则 \(F\) 在零点可去。若其次数为 \(d\ge1\)，最高次系数为 \(a_d\ne0\)，则
\[
 F(w)=w^{-d}\bigl(a_d+a_{d-1}w+\cdots+a_0w^d\bigr).
\]
它在零点的极点阶数恰为 \(d\)，故 \(d=m\)。这个式子也直接证明任意 \(m\) 次多项式在无穷远有 \(m\) 阶极点。最后，由孤立奇点三类的穷尽性，余下的非多项式整函数恰对应本性奇点。
\end{proof}

指数函数和正弦函数的 Taylor 展开都有无穷多个非零系数，因此它们在无穷远都是本性奇点；\(z^m\) 在无穷远则有 \(m\) 阶极点。另一方面，\(1/z\) 换成 \(w\) 后为 \(w\)，在无穷远可去且取值零，却不是常函数，因为它在有限点零处不全纯。命题中的整性负责排除这种有限奇点，不能只由大圆外的衰减推出函数恒零。

% !TeX root = ../../Book3.tex
\section{留数}
\label{b3:sec:0303}

Laurent 展开的所有系数记录局部行为，闭曲线积分却只挑出 \(-1\) 次项。将多个孤立奇点围起来时，各个局部贡献相加，这使许多实积分也可以转化为有限次代数计算。

\subsection{留数}
\label{b3:subsec:030301}

\begin{definition}\label{b3:def:03-residue}
设 \(a\) 是全纯函数 \(f\) 的孤立奇点。其 Laurent 系数 \(c_{-1}\) 称为 \(f\) 在 \(a\) 的\term[liushu]{留数}[residue]，记为 \(\operatorname{Res}(f,a)\)。
\end{definition}
\symidx{\(\operatorname{Res}(f,a)\)：孤立奇点处的留数}
\begin{proposition}\label{b3:prop:03-residue-formulas}
若 \(a\) 是 \(m\) 阶极点，则
\[
 \operatorname{Res}(f,a)=\frac1{(m-1)!}\left.\frac{\mathrm{d}^{m-1}}{\mathrm{d}z^{m-1}}\bigl((z-a)^mf(z)\bigr)\right|_{z=a}.
\]
特别地，若 \(f=g/h\)，\(g,h\) 在 \(a\) 附近全纯，\(h(a)=0\)、\(h'(a)\ne0\)，则 \(\operatorname{Res}(f,a)=g(a)/h'(a)\)。
\end{proposition}
\begin{proof}
令 \(G=(z-a)^mf\)。其 Taylor 展开的 \((m-1)\) 次系数，除以 \((z-a)^m\) 后就是 \(-1\) 次系数。单零分母写成 \(h(z)=(z-a)H(z)\)，\(H(a)=h'(a)\)，故 \((z-a)f\) 在中心值为 \(g(a)/h'(a)\)。即使 \(g(a)=0\)，公式仍给留数0，此时奇点可能已经可去。
\end{proof}
对本性奇点仍可用级数提取留数，例如 \(e^{1/z}\) 在0的留数为1；有限阶求导公式却没有可用的整数 \(m\)。留数为零不意味着奇点可去，\(1/z^2\) 就是一例。

\ProseParagraphBreak

留数在合法的局部换元下属于同一个微分量。
\begin{proposition}[留数的坐标不变性]\label{b3:prop:03-residue-coordinate}
设 \(z=\psi(w)\) 是 \(b\) 附近的双全纯坐标，\(\psi(b)=a\)，而 \(f\) 在 \(a\) 有孤立奇点。则
\[
 \operatorname{Res}(f,a)=\operatorname{Res}\bigl(f(\psi(w))\cdot\psi'(w),b\bigr).
\]
\end{proposition}
\begin{proof}
取足够小的正向 \(w\) 圆，其像是包围 \(a\) 的正向简单闭曲线；全纯非零导数保持方向。曲线换元使两个微分积分相同。像曲线与一个小 \(z\) 圆在去心邻域之间可变形，Cauchy 定理使积分相同。两边各除以 \(2\pi\mathrm{i}\)，得到留数相等。
\end{proof}
因此只把 \(f(\psi(w))\) 拿来算留数通常是错误的，必须乘 \(\psi'(w)\)。例如 \(z=2w\)、\(f(z)=1/z\)，合成函数留数是 \(1/2\)，合成微分乘2后才是1。无穷远留数的 \(-w^{-2}\) 因子就是这一坐标规则的具体实例。

\ProseParagraphBreak

一个全纯函数的导数在任意孤立奇点处留数为零。将 Laurent 展开逐项微分，要产生 \(z^{-1}\) 项只能来自原来的 \(z^0\) 项，但它的导数系数为0。反向对同心环域也成立：留数为零时可逐项积分得到原函数。留数检测的是一个微分是否有非零局部周期，而不只是函数值增长快慢。

\subsection{留数定理}
\label{b3:subsec:030302}

\term[liushu-dingli]{留数定理}[residue theorem]先以简单围道版本建立。
\begin{theorem}[留数]\label{b3:thm:03-residue}
设 \(V\) 是有界 Jordan 区域，边界正向且分段光滑。若 \(f\) 在 \(\overline V\) 的邻域内除有限个内部孤立奇点 \(a_1,\ldots,a_N\) 外全纯，则
\[
 \int_{\partial V}f(z)\,\mathrm{d}z=2\pi\mathrm{i}\sum_{j=1}^N\operatorname{Res}(f,a_j).
\]
\end{theorem}
\begin{proof}
在每个奇点周围选互不相交、闭包包含于 \(V\) 的小圆盘。将挖去这些圆盘的区域沿有限条连接弧切开，再按有限小片细分；各小片上有局部原函数，其边界积分为零。内部边和切口两侧相消，留下外边界及顺时针小圆。故外边界积分等于各个逆时针小圆积分之和。小圆上 Laurent 展开一致收敛，逐项积分只有 \(-1\) 次项留下 \(2\pi\mathrm{i}\)，即得结论。若连接弧或边界不是多边形，可先在局部原函数圆盘中更换为多边形代表，端点公式保持积分。
\end{proof}
定理的有限性来自奇点与边界分离。若奇点向边界积聚，不能未经收敛核对就把有限和改成无穷和。自交闭曲线则需要把每个奇点的绕数作为权重，精确形式在\chapref{b3:ch:10}建立。

\subsection{实积分}
\label{b3:subsec:030303}

计算实积分时，先说明所选围道和附加弧积分为何消失。以
\[
 I=\int_{-\infty}^{\infty}\frac{\mathrm{d}x}{x^2+1}
\]
为例，在上半平面取半径 \(R>2\) 的半圆，内部仅有极点 \(\mathrm{i}\)，留数为 \(1/(2\mathrm{i})\)。半圆上的函数模不超过 \(1/(R^2-1)\)，弧长为 \(\pi R\)，其积分趋于0。实轴截断积分单调趋于反常积分，故 \(I=\pi\)。

\ProseParagraphBreak

若再有振荡因子 \(e^{\mathrm{i}tz}\)，\(t>0\)，在上半平面其模为 \(e^{-t\operatorname{Im}z}\le1\)。同一弧估计仍适用，留数变为 \(e^{-t}/(2\mathrm{i})\)，得到
\[
 \int_{-\infty}^{\infty}\frac{e^{\mathrm{i}tx}}{x^2+1}\,\mathrm{d}x=\pi e^{-t}.
\]
负 \(t\) 时选下半平面，或对称换元，得到 \(\pi e^{-|t|}\)。如果分母只衰减如 \(1/z\)，上述弧长估计并不趋于0，须另外证明振荡提供的衰减，不能直接援引“半圆消失”。

\ProseParagraphBreak

对 \(\int_0^{2\pi}R(\cos\theta,\sin\theta)\,\mathrm{d}\theta\)，置 \(z=e^{\mathrm{i}\theta}\)，则 \(\cos\theta=(z+z^{-1})/2\)、\(\sin\theta=(z-z^{-1})/(2\mathrm{i})\)、\(\mathrm{d}\theta=\mathrm{d}z/(\mathrm{i}z)\)。最后的微分因子不可遗漏，单位圆上有极点时则不能直接使用通常留数定理。

围道选择须同时控制奇点和大圆弧。对 \(\int_{\mathbb R}e^{\mathrm{i}ax}/(x^2+1)\,\mathrm{d}x\)，当 \(a>0\) 时闭向上半平面，因为 \(|e^{\mathrm{i}az}|=e^{-a\operatorname{Im}z}\le1\)；当 \(a<0\) 时应闭向下半平面，并记下顺时针的负号。若不随 \(a\) 改变半平面，大圆弧上的指数会增长，留数计算虽形式正确却不能代表原积分。最终两种选择都给 \(\pi e^{-|a|}\)，而 \(a=0\) 可以由有理函数估计直接处理。

\ProseParagraphBreak
实轴上有极点时还须说明积分的意义。通常的广义积分、Cauchy 主值以及绕开奇点的小半圆贡献并不相同。留数定理只计算闭围道积分，不会替读者选择其中一种意义。先标明每段截断积分和小圆弧的方向，再取极限，才能看清哪些项趋于零、哪些项留下半留数贡献。

\subsection{无穷远点}
\label{b3:subsec:030304}

无穷远点用坐标 \(w=1/z\) 检查函数奇点：研究 \(f(1/w)\) 在 \(w=0\) 的分类。留数属于微分 \(f(z)\,\mathrm{d}z\)，换元又给出额外的 \(-w^{-2}\)：
\begin{definition}\label{b3:def:03-infinity-residue}
若 \(f\) 在某个大圆之外全纯，定义
\[
 \operatorname{Res}(f,\infty)\coloneqq-\operatorname{Res}\left(\frac{f(1/w)}{w^2},0\right).
\]
\end{definition}
例如 \(f(z)=1/z\) 在无穷远作为函数可去，中心值为0，但微分的无穷远留数为 \(-1\)。函数的奇点分类和微分的留数不是同一问题。
\begin{proposition}\label{b3:prop:03-global-residue}
对有理函数 \(f\)，所有有限极点留数与无穷远留数之和为零。
\end{proposition}
\begin{proof}
选大圆包围全部有限极点，有限留数和为大圆积分除以 \(2\pi\mathrm{i}\)。换元 \(w=1/z\) 把逆时针大圆变成顺时针小圆，微分为 \(-w^{-2}\,\mathrm{d}w\)。根据定义，该积分是 \(-2\pi\mathrm{i}\operatorname{Res}(f,\infty)\)，故留数总和为零。
\end{proof}
计算无穷远留数时也可直接取 \(f\) 在大圆外 Laurent 展开的 \(z^{-1}\) 系数再加负号。例如 \(z^2/(z^3-1)\) 的这一系数为1，所以无穷远留数为 \(-1\)。

\ProseParagraphBreak

扩充平面上的亚纯函数实际上全部是有理函数。
\begin{theorem}\label{b3:thm:03-sphere-rational}
若 \(f\) 在扩充复平面上亚纯，则存在多项式 \(P,Q\)，\(Q\) 不恒零，使 \(f=P/Q\)。
\end{theorem}
\begin{proof}
极点在球面上离散，而球面紧，故极点有限；否则无限点列有聚点，与该处亚纯局部模型矛盾。先在每个有限极点减去其有限 Laurent 主部，所得函数 \(g\) 在整个有限平面全纯。在无穷远，原函数与减去的有理主部仍亚纯，故 \(g(1/w)\) 在0可去或有有限阶极点。于是存在整数 \(m\ge0\)，使大 \(|z|\) 时 \(|g(z)|\le C|z|^m\)。由第二章的增长版 Liouville 推论，\(g\) 是多项式。加回有限个有理主部，得到有理函数。
\end{proof}
这一定理利用了无穷远也亚纯的条件。整函数 \(e^z\) 在有限平面无极点，却在无穷远有本性奇点，所以不是反例。它也说明有理函数在球面只需有限多局部数据，整函数的无穷远增长则可能包含无限信息。

% !TeX root = ../../Book3.tex
\section{零点计数}
\label{b3:sec:0304}

函数值绕原点转了多少圈，可以通过 \(f'/f\) 的积分读出。零点和极点分别使辐角增加和减少；小扰动若不能把边界像推过原点，内部零点总数就不会变化。

\subsection{辐角原理}
\label{b3:subsec:030401}

对不经过原点的闭连续路径 \(\eta\)，可沿参数区间选连续辐角 \(\theta\)。具体地，把区间细分，使每段像位于一个不包含0的对数圆盘，在相邻端点给辐角加适当的 \(2\pi\) 整数以接合。因首尾函数值相同，\(\theta(b)-\theta(a)\) 必为 \(2\pi\) 的整数倍，其商称为\term[raoshu]{绕数}[winding number]。若路径分段光滑，局部对数求导后相加，得到
\[
 \operatorname{Ind}(\eta,0)=\frac1{2\pi\mathrm{i}}\int_\eta\frac{\mathrm{d}w}{w}.
\]
\term[fujiao-yuanli]{辐角原理}[argument principle]把该整数转成零点计数。
\begin{theorem}[辐角原理]\label{b3:thm:03-argument-principle}
设 \(f\) 在有界分段光滑 Jordan 区域 \(V\) 的闭包邻域亚纯，边界无零点也无极点。则
\[
 \frac1{2\pi\mathrm{i}}\int_{\partial V}\frac{f'(z)}{f(z)}\,\mathrm{d}z=N-P,
\]
其中 \(N\) 与 \(P\) 分别是内部零点和极点的总重数。
\end{theorem}
\begin{proof}
紧性与边界无零极点使内部只有有限个零极点。在零点 \(a\) 的 \(m\) 阶分解中，\(f=(z-a)^mg\)，故 \(f'/f=m/(z-a)+g'/g\)，留数为 \(m\)。在 \(m\) 阶极点处同样得到留数 \(-m\)。留数定理相加给出公式；链式法则说明左端也等于像曲线 \(f\circ\partial V\) 绕0的绕数。
\end{proof}

\subsection{Rouché 定理}
\label{b3:subsec:030402}

\term[rouche-dingli]{Rouché 定理}[Rouché's theorem]允许把难计数的函数与简单函数比较。
\begin{theorem}[Rouché]\label{b3:thm:03-rouche}
设 \(f,g\) 在有界分段光滑 Jordan 区域 \(V\) 的闭包邻域全纯。若边界处 \(|g|<|f|\)，则 \(f\) 与 \(f+g\) 在 \(V\) 内的零点总重数相同。
\end{theorem}
\begin{proof}
对 \(0\le t\le1\)，令 \(f_t=f+tg\)。边界上 \(|f_t|\ge|f|-|g|>0\)，紧边界使右端有共同正下界。辐角原理给出的数
\(N(t)=(2\pi\mathrm{i})^{-1}\int_{\partial V}f_t'/f_t\,\mathrm{d}z\)
依 \(t\) 连续，又总为整数，故在区间上为常数。取 \(t=0,1\) 即得结论。
\end{proof}
例如在 \(|z|=2\) 上，\(|z^5|=32>2+1\)，故 \(z^5+z+1\) 在盘内有5个零点。严格不等式排除了整个变形中边界出现零点。若只有 \(|g|\le|f|\)，边界可能产生零点，开域中的计数就可能改变。

\ProseParagraphBreak

Rouché 的比较方向由边界决定，不能只看多项式的次数。在 \(|z|=1\) 上，\(z^4+3z+1\) 应比较 \(3z\)，因为余项模不超过2；在 \(|z|=2\) 上，\(z^4\) 的模为16，大于 \(3|z|+1=7\)，应比较最高次项。因此同一个多项式在单位盘有1个零点，在半径2盘有4个零点，其余3个位于两圆之间，均计重数。不同半径的零点计数可以相减定位环域中的零点，但边界必须各自无零点。

\ProseParagraphBreak

参数变化时，若一个紧参数区间内的边界函数始终无零点，零点数就保持不变。共同正下界来自参数区间与边界乘积的紧性。如果参数范围不紧，最低模可能趋于0，即使每个固定参数都无边界零点，也不能得到统一的扰动余量。实际使用时可局部按参数取一个小区间，再连通地传播整数计数；不必先找到一个覆盖所有无界参数的误差常数。

\subsection{重数}
\label{b3:subsec:030403}

零点个数若不计重数，就不具备稳定性。函数 \(z^2-\varepsilon\) 在 \(\varepsilon\ne0\) 时有两个不同零点，极限 \(z^2\) 只有一个零点，然而总重数始终为2。
\begin{proposition}[零点局部稳定]\label{b3:prop:03-local-zero-stability}
设 \(f\) 在 \(a\) 处有 \(m\) 阶零点。在充分小的闭圆盘 \(\overline{D(a,r)}\) 上，若全纯函数 \(h\) 满足边界估计 \(|h-f|<\min_{|z-a|=r}|f(z)|\)，则 \(h\) 在圆盘内恰有 \(m\) 个零点，计重数。
\end{proposition}
\begin{proof}
由孤立零点定理选圆盘，使 \(f\) 在其中只有中心零点且边界无零点。给定边界估计就是 Rouché 定理对 \(f\) 与 \(h-f\) 的条件，故两者零点总重数相同。
\end{proof}
若希望定位单个零点，先在每个已知零点周围选互不相交的小圆盘，再分别估计。只在一个大圆上计总数，不会自动给出各个零点的精确位置。

Rouché 定理中的严格不等式只需在边界成立，不要求内部成立。设 \(p(z)=z^5+10z+1\)，在单位圆上以 \(10z\) 为主项，因为 \(|z^5+1|\le2<10\)，可知内部有一个零点；在半径2的圆上，\(|z^5|=32>21\ge|10z+1|\)，内部有五个零点。因此其余四个零点落在 \(1<|z|<2\)。这比求根公式更直接，并且适用于无法写出根的高次多项式。

\ProseParagraphBreak
若把严格不等式放宽为处处等号允许，结论会失败。单位圆上 \(|z|=|1|\)，但 \(z\) 在内部有一个零点，常数1没有零点。严格性保障整个变形 \(f+t g\) 在边界无零；一旦某个参数的零点经过边界，整数零点数便可能跳变。实际计算通常先在紧边界上寻找一个正的间隙，再允许小参数扰动。这个间隙是解析问题中控制根的稳定性的数值。

\subsection{对数导数}
\label{b3:subsec:030404}

对数导数还可以提取带权的零点信息。
\begin{proposition}\label{b3:prop:03-weighted-zeros}
在辐角原理的条件下，若 \(h\) 在 \(\overline V\) 的邻域全纯，则
\[
 \frac1{2\pi\mathrm{i}}\int_{\partial V}h(z)\cdot\frac{f'(z)}{f(z)}\,\mathrm{d}z
 =\sum_a m_a h(a)-\sum_b n_b h(b),
\]
两和分别遍历内部零点与极点，\(m_a,n_b\) 为阶数。
\end{proposition}
\begin{proof}
对数导数在 \(a\) 的主部只有 \(m_a/(z-a)\)，乘全纯 \(h\) 后留数为 \(m_ah(a)\)；极点处相同计算带负号。直接应用留数定理。
\end{proof}
取 \(h(z)=z\) 得到零点之和与极点之和的差；取 \(h(z)=z^k\) 得到幂和。若 \(f\) 在整个区域无零点，\(f'/f\) 本身全纯；它有没有全局原函数，还须看所有闭曲线周期是否为零。\chapref{b3:ch:05}将用这一条件构造全纯对数。

局部零点稳定性始终按重数计数。不同简单根可以合成重根，重根也可分裂；只有在分别选定无零边界的小圆盘之后，才有可持续追踪的局部零点总数。


\begin{chaptersummary}
环域的内外边界分别产生 Laurent 展开的负幂和正幂。可去奇点、极点与本性奇点由主部的三种形态区分，Casorati--Weierstrass 定理展示本性奇点的局部像。围道积分只保留各奇点的留数；对数导数又把零点和极点的阶数转换为留数。Rouché 定理依赖边界变形始终避开零点，因此稳定的是总重数而不是不同零点的个数。
\end{chaptersummary}

\BookExercises

% 本章题组自拟；以正文已建立的工具设计计算、反例、构造与综合训练。
\begin{exercise}\label{b3:ex:03-1}
分别在 \(0<|z|<1\)、\(1<|z|<2\)、\(|z|>2\) 展开 \(1/((z-1)(z-2))\)。
\end{exercise}
\begin{solution}
先分解为 \(-1/(z-1)+1/(z-2)\)。内环展开为 \(\sum_{n\ge0}(1-2^{-n-1})z^n\)；中环为 \(-\sum_{n\ge0}z^{-n-1}-\sum_{n\ge0}z^n/2^{n+1}\)；外环为 \(\sum_{n\ge0}(2^n-1)z^{-n-1}\)。各几何级数分别按所给半径条件收敛。
\end{solution}

\begin{exercise}\label{b3:ex:03-2}
分类零点处 \(\sin z/z\)、\((e^z-1)/z^3\)、\(e^{1/z}-1/z\) 的奇点并求留数。
\end{exercise}
\begin{solution}
第一项可去，留数0。第二项展开为 \(z^{-2}+\tfrac12z^{-1}+\tfrac16+\cdots\)，是二阶极点，留数 \(1/2\)。第三项负幂仍无限，但 \(z^{-1}\) 项被消掉，是本性奇点，留数0。
\end{solution}

\begin{exercise}\label{b3:ex:03-3}
若 \(f\) 在去心圆盘全纯并满足 \(|f(z)|\le C|z|^{-3/2}\)，证明它在0最多有一阶极点。
\end{exercise}
\begin{solution}
负指标系数满足 \(|c_{-m}|\le C\rho^{m-3/2}\)。对 \(m\ge2\)，令 \(\rho\downarrow0\) 得系数为零；因此主部只可能有 \(c_{-1}/z\)。该系数非零时是一阶极点，为零时可去。
\end{solution}

\begin{exercise}\label{b3:ex:03-4}
计算 \(\int_{-\infty}^{\infty}(x^2+1)^{-2}\,\mathrm{d}x\)。
\end{exercise}
\begin{solution}
在上半圆内只有 \(\mathrm{i}\) 的二阶极点。留数是 \(\left.(\mathrm{d}/\mathrm{d}z)(z+\mathrm{i})^{-2}\right|_{z=\mathrm{i}}=-2/(2\mathrm{i})^3=1/(4\mathrm{i})\)。弧估计不超过 \(\pi R/(R^2-1)^2\)，趋于0。因此积分为 \(2\pi\mathrm{i}/(4\mathrm{i})=\pi/2\)。
\end{solution}

\begin{exercise}\label{b3:ex:03-5}
证明 \(z^4+3z+1\) 在单位圆盘内恰有一个零点。
\end{exercise}
\begin{solution}
在单位圆上 \(|z^4+1|\le2<3=|3z|\)。用 \(3z\) 作为比较函数，Rouché 给出零点总重数为1；因此该零点也必为单零点。
\end{solution}

\begin{exercise}\label{b3:ex:03-6}
设 \(P\) 是次数 \(n\ge1\) 的多项式，所有零点都在 \(|z|<R\)。用围道积分表达这些零点之和与平方和，计重数。
\end{exercise}
\begin{solution}
分别取加权辐角原理中的 \(h(z)=z\) 与 \(z^2\)，得到 \((2\pi\mathrm{i})^{-1}\int_{|z|=R}zP'/P\,\mathrm{d}z\) 和 \((2\pi\mathrm{i})^{-1}\int_{|z|=R}z^2P'/P\,\mathrm{d}z\)。多项式无有限极点，所以无需减极点项。
\end{solution}

\begin{exercise}\label{b3:ex:03-7}
求 \(z^2/(z^3-1)\) 在三个有限极点以及无穷远的留数。
\end{exercise}
\begin{solution}
每个三次单位根 \(a\) 处，单零分母公式给出 \(a^2/(3a^2)=1/3\)。大圆外展开首项为 \(1/z\)，故无穷远留数为 \(-1\)。总和为 \(3\cdot1/3-1=0\)。
\end{solution}

\begin{exercise}\label{b3:ex:03-8}
设 \(f\) 在 \(0<|z|<1\) 全纯且有本性奇点。证明存在趋于0的点列 \(z_n\)，使 \(f(z_n)\rightarrow b\)，其中 \(b\) 是任意给定复数。是否能由此保证 \(f(z_n)=b\)？
\end{exercise}
\begin{solution}
对每个 \(n\)，Casorati--Weierstrass 稠密性在 \(0<|z|<1/n\) 中给出 \(|f(z_n)-b|<1/n\)。所以同时有 \(z_n\rightarrow0\) 与函数值趋于 \(b\)。精确取值不能保证，\(e^{1/z}\) 永远不取0。
\end{solution}

% !TeX root = ../../Book3.tex
\chapter{正规族}
\label{b3:ch:04}
\BookChapterNumber{b3:ch:04}
\begin{chapterabstract}
局部一致收敛保留全纯性和固定阶导数。Cauchy 估计把局部有界族变成等度连续族，经有限中心与对角线抽取得到 Montel 紧性；Rouché 定理又给出 Hurwitz 零点稳定性及单叶极限。球面距离统一处理极点与无穷极限，并区分两种正规性。
\end{chapterabstract}

\begin{learninggoals}
\begin{itemize}
\item\label{b3:item:04-goal-1}能区分逐点、局部一致和全域一致收敛，并核对导数极限。
\item\label{b3:item:04-goal-2}能用紧集共同上界与 Cauchy 估计证明等度连续。
\item\label{b3:item:04-goal-3}能完成两次对角线抽取，说明所得极限为同一个全纯函数。
\item\label{b3:item:04-goal-4}能判断无零性与单叶性的保留范围，并识别常数退化。
\end{itemize}
\end{learninggoals}

列 \(z^n\) 在开单位圆盘的每个内部紧集上趋于0，却在整个圆盘上保持上确界1。列 \(nz\) 的行为更剧烈：原点处固定为0，任何非零点模趋于无穷。函数列的局部尖峰与边界逃逸需要不同处理，正规族理论正是为了从可控函数列中选出仍可使用的极限。

\ProseParagraphBreak

本章使用\cref{b3:thm:02-cauchy-estimate}和\cref{b3:thm:03-rouche}。Joel L. Schiff 的《Normal Families》\cite{Schiff1993Normal}是这一主题的专著；本章先完整证明局部有界版本，共同省略点值的更强版本在\secref{b3:sec:1403}给出证明。

% !TeX root = ../../Book3.tex
\section{局部一致收敛}
\label{b3:sec:0401}

函数列逐点收敛，未必保留连续性；即使极限全纯，也未必可以逐项求导。Cauchy 积分公式给出的控制恰好说明，复分析应采用紧集上的一致收敛，而不是仅固定每个点观察。

\subsection{紧一致收敛}
\label{b3:subsec:040101}

\begin{definition}\label{b3:def:04-compact-convergence}
设 \(U\) 为开集，\(f_n,f:U\rightarrow\mathbb C\)。若对每个紧集 \(K\subset U\) 都有 \(\sup_K|f_n-f|\rightarrow0\)，则称 \(f_n\) 在 \(U\) 上\term[jubu-yizhi-shoulian]{局部一致收敛}[locally uniform convergence]于 \(f\)，也称紧一致收敛。
\end{definition}
开集可由递增紧集穷尽，例如取
\[
 K_j=\{\,z\in U\mid |z|\le j,\ \operatorname{dist}(z,\mathbb C\setminus U)\ge1/j\,\}.
\]
空的早期集合可略去，\(U=\mathbb C\) 时只取闭圆盘。每个 \(K_j\) 闭有界，且任意 \(U\) 内紧集最终包含于其中；因此在这一列紧集上检验一致收敛就足够。

\ProseParagraphBreak

函数列 \(z^n\) 在 \(\mathbb D\) 上局部一致趋于0，因为 \(|z|\le r<1\) 时误差不超过 \(r^n\)；却不在整个开圆盘上一致趋于0，其上确界始终是1。局部一致性不对趋近边界的点给共同误差。后续紧性抽取也只在这一拓扑中进行。

\ProseParagraphBreak

紧集穷尽并不是把开集改成一个闭区域。对于有孔且无界的区域，\(K_j\) 同时截断无穷远，并远离每个孔的边界；一个固定小盘最终包含其中，但随着 \(j\) 增大，仍可能出现新的边界附近部分。若某个 \(K_j\) 为空，不需要为它另设特殊函数空间。穷尽所要求的是每个内部紧集最终被包含，而不是每一步都有非空内部。

\ProseParagraphBreak

局部一致收敛也不会自动控制随指标移动的点。若 \(z_n\) 始终位于同一个紧集且 \(z_n\rightarrow z\)，则
\[
 |f_n(z_n)-f(z)|\le\sup_K|f_n-f|+|f(z_n)-f(z)|\longrightarrow0.
\]
第一项用一致收敛，第二项用极限连续。若 \(z_n\) 逃向边界，这个共同紧集不复存在。取 \(f_n(z)=z^n\)、\(z_n=1-1/n\)，便有 \(f_n(z_n)\rightarrow e^{-1}\)，而固定点处极限为0。两种量词不能交换。

\subsection{全纯极限}
\label{b3:subsec:040102}

\begin{proposition}\label{b3:prop:04-holomorphic-limit}
设 \(f_n\) 在开集 \(U\) 上全纯，且局部一致收敛于 \(f\)。则 \(f\) 在 \(U\) 上全纯。
\end{proposition}
\begin{proof}
先在每个小闭圆盘上用连续函数的一致极限得 \(f\) 连续。固定 \(\overline{D(a,r)}\subset U\)，每个 \(f_n\) 的 Cauchy 公式为
\(f_n(z)=(2\pi\mathrm{i})^{-1}\int_{|\zeta-a|=r}f_n(\zeta)/(\zeta-z)\,\mathrm{d}\zeta\)。圆周上一致收敛使积分在每个较小闭圆盘上一致趋于相同的 \(f\) 核积分。核积分可按 \(z\) 微分，因分母与零保持正距离，所以右端全纯。由 \(z\in D(a,r)\) 时的点值极限，它恰为 \(f(z)\)。小圆盘覆盖 \(U\)，因此处处全纯。
\end{proof}
这里既没有先求 \(f_n'\) 的逐点极限，也没有从实变量可微函数列的性质类推。连续性与全纯性是通过积分表示先后获得的。

\subsection{导数收敛}
\label{b3:subsec:040103}

\begin{proposition}\label{b3:prop:04-derivative-convergence}
在上一命题的条件下，对每个固定整数 \(k\ge1\)，\(f_n^{(k)}\) 在 \(U\) 上局部一致趋于 \(f^{(k)}\)。
\end{proposition}
\begin{proof}
固定紧集 \(K\subset U\)，选 \(\delta>0\)，使所有以 \(z\in K\) 为中心、半径 \(\delta\) 的闭圆盘都在 \(U\) 中。它们的并闭包含于一个紧集 \(K'\subset U\)。对全纯差 \(f_n-f\) 使用 Cauchy 估计，得到
\[
 \sup_K|f_n^{(k)}-f^{(k)}|\le k!\delta^{-k}\sup_{K'}|f_n-f|\longrightarrow0.
\]
选取同一个 \(\delta\) 是紧性的用途，因而结论对 \(K\) 上全部点一致成立。
\end{proof}
导数阶 \(k\) 必须固定；因子 \(k!\delta^{-k}\) 随阶数增长，不能由这一结论让 \(k\) 与 \(n\) 同时趋于无穷。以 \(z^n\) 为例，它在较小圆盘的一阶导数 \(nz^{n-1}\) 仍趋于0，但第 \(n\) 阶导数为 \(n!\)。

\ProseParagraphBreak

反向也有一个带归一化的积分判据。若区域中的 \(f_n'\) 局部一致趋于 \(g\)，且某个基点值 \(f_n(a)\) 收敛，则 \(f_n\) 局部一致收敛。证明时在一个紧集的有限圆盘覆盖中，用从基点到各圆盘中心的有限条固定多边形路径，再在盘内接短线段。积分差由路径长度与导数差控制；有限路径保证同一个误差对整紧集有效。极限函数的导数为 \(g\)，由局部差商或原函数端点式即可核对。

\ProseParagraphBreak

基点值的条件不可省略，常值函数列 \(f_n=n\) 的导数全部为0，却没有有限局部极限。若定义域不连通，应每个分支固定一个值；一个基点不能控制其他分支上的加法常数。这个反向判据把“积分一次多出常数”落实为收敛中的具体条件。

\subsection{Weierstrass 定理}
\label{b3:subsec:040104}

前两项通常合称\term[weierstrass-shoulian-dingli]{Weierstrass 收敛定理}[Weierstrass convergence theorem]。
\begin{theorem}[Weierstrass]\label{b3:thm:04-weierstrass}
全纯函数列的局部一致有限极限全纯，且任意固定阶导数局部一致收敛至极限函数的相应导数。若全纯函数项级数在每个紧集上一致收敛，则其和全纯，并可任意固定阶逐项微分。
\end{theorem}
\begin{proof}
函数列结论是前两命题。函数项级数取有限部分和，它们全纯；导数部分和的局部一致极限就是和函数的导数，因此逐项微分合法。
\end{proof}
若 \(\sum_n\sup_K|g_n|<\infty\) 对每个紧集成立，Weierstrass 判别法保证所需的一致收敛。例子 \(\sum_{n\ge1}e^{-n}z^n\) 在 \(|z|<e\) 局部一致收敛；在更大区域则不能借每项整函数声称和仍整。级数收敛域是另外必须核对的数据。

紧集一致收敛还保持局部积分恒等式。若 \(f_n\rightarrow f\) 在紧集上一致，\(\gamma\) 是区域中一条固定分段光滑曲线，则
\[
 \left|\int_\gamma(f_n-f)\,\mathrm{d}z\right|
 \le L(\gamma)\sup_{\gamma}|f_n-f|\longrightarrow0.
\]
曲线的像是紧集，长度又固定，两个条件都在这里被使用。如果曲线随 \(n\) 增长或靠近边界，上式右端未必趋于0。对单位圆盘中的 \(z^n\)，紧集一致趋于0，但在 \(|z|=1\) 上模恒为1，不能把内点极限自动延伸到闭圆盘。

\ProseParagraphBreak
另一个有用例子是 \(f_n(z)=n\sin(z/n)\)。Taylor 展开在每个固定圆盘上给 \(f_n(z)=z+O(1/n^2)\)，误差常数依赖于圆盘半径。因此 \(f_n\) 紧集一致趋于恒等映射，所有阶导数也有相应极限；但它们在整个平面上的增长与零点分布很不相同。这里的收敛只控制固定的观察窗口，远处的零点可以随 \(n\) 逃向无穷远。

% !TeX root = ../../Book3.tex
\section{正规族}
\label{b3:sec:0402}

研究一个函数时，积分公式恢复它的内部值；研究一族函数时，同一公式控制整个族的变化。正规性要求任意序列都有可用的子列，因而是一种紧性性质，而不是每个成员的光滑程度。

\subsection{定义}
\label{b3:subsec:040201}

\begin{definition}\label{b3:def:04-normal-finite}
设 \(\mathcal F\) 是区域 \(U\) 上的全纯函数族。若任取 \(f_n\in\mathcal F\)，都能抽取在 \(U\) 上局部一致收敛至有限全纯函数的子列，则称 \(\mathcal F\) 在有限值意义下是\term[zhenggui-zu]{正规族}[normal family]。
\end{definition}
极限不必属于原族，所以这里对应相对序列紧性。族 \(\{z^n:n\ge1\}\) 在圆盘上正规，其某些子列极限为0，但0原本不在该族内。族 \(\{nz:n\ge1\}\) 在包含0的圆盘上不正规：若局部一致有有限极限，导数在0也应收敛，而相应导数为无界整数。

\ProseParagraphBreak

有些教材允许子列趋于无穷。为了处理极点，下面改用球面距离统一表述；有限值版本和球面版本不能未经说明互换。

\subsection{序列紧性}
\label{b3:subsec:040202}

局部一致拓扑可以由穷尽紧集定义度量
\[
 \mathrm{d}_{\mathrm{loc}}(f,g)=\sum_{j=1}^{\infty}2^{-j}\min\{1,\sup_{K_j}|f-g|\}.
\]
级数始终收敛；三角不等式从每个截断上确界得到。该度量趋于0当且仅当所有 \(K_j\) 上一致趋于0：前向由第 \(j\) 项的下界得到，反向先使尾和小，再控制有限个前项。因此“抽子列”与这个度量空间的序列紧性相对应。

\ProseParagraphBreak

对角线抽取的顺序不能省略。先在 \(K_1\) 选收敛子列，再在其内选 \(K_2\) 上的子列，依次进行，最后取第 \(j\) 列的第 \(j\) 个成员。对任何固定 \(K_m\)，对角线尾部都属于第 \(m\) 次子列，所以仍一致收敛；各紧集上的极限在重叠处一致，拼成 \(U\) 上同一个函数。若只对每点分别抽子列，就不能保证得到了同一条子列。

\ProseParagraphBreak

局部正规性能够拼成全域正规性。
\begin{proposition}\label{b3:prop:04-normal-locality}
设 \(\mathcal F\) 为开集 \(U\) 上的全纯函数族。若每点都有邻域，使限制族在该邻域有限值正规，则整个族在 \(U\) 上有限值正规。
\end{proposition}
\begin{proof}
平面第二可数性使这些邻域有可数子覆盖 \(V_1,V_2,\ldots\)。对任意函数序列，先在 \(V_1\) 抽取局部一致子列，再依次在 \(V_j\) 抽取，最后用对角线得到每个 \(V_j\) 上局部一致收敛的同一子列。重叠处的点值极限唯一，所以局部极限拼成 \(U\) 上全纯函数。若 \(K\subset U\) 紧，可用有限个位于各 \(V_j\) 内的小闭盘覆盖 \(K\)，这些盘上的一致收敛给出全 \(K\) 一致收敛。
\end{proof}
可数性与紧性承担不同作用：可数子覆盖允许一次对角线抽取，紧集的有限子覆盖把局部误差合成一个全紧集误差。只说“每点可以抽子列”而不经过这两个步骤，还没有得到正规性。

\subsection{球面度量}
\label{b3:subsec:040203}

扩充平面的\term[qiumian-juli]{球面距离}[chordal distance]取为
\[
 \chi(z,w)=\frac{|z-w|}{\sqrt{1+|z|^2}\sqrt{1+|w|^2}},\qquad
 \chi(z,\infty)=\frac1{\sqrt{1+|z|^2}}.
\]
它是单位球面弦长的一半，经立体投影得到，故满足三角不等式。归一化常数不会改变正规性，但应固定一种写法。倒数变换满足 \(\chi(1/z,1/w)=\chi(z,w)\)，包括0和无穷远。

\ProseParagraphBreak

在有限值有界集合上，球面距离与欧氏距离等价：若 \(|z|,|w|\le M\)，则 \(|z-w|/(1+M^2)\le\chi(z,w)\le|z-w|\)。趋向无穷则表示模一致变大。函数 \(n\) 的常值列没有有限局部极限，却在球面意义下一致趋向 \(\infty\)。这一点使两种正规性明确不同。

\subsection{亚纯函数族}
\label{b3:subsec:040204}

\begin{definition}\label{b3:def:04-meromorphic-normal}
区域 \(U\) 上的亚纯函数族若每个序列都有球面局部一致收敛子列，其极限是亚纯函数或常值 \(\infty\)，则称它在球面意义下正规。
\end{definition}
\begin{proposition}\label{b3:prop:04-spherical-limit}
亚纯函数列的球面局部一致极限，若存在，则是亚纯函数或常值 \(\infty\)。
\end{proposition}
\begin{proof}
球面一致极限先是到球面的连续映射 \(f\)。若 \(f(a)\) 有限，缩小圆盘使像避开 \(\infty\) 的一个邻域；充分大的 \(f_n\) 在这个圆盘也有限且一致有界，球面收敛转成欧氏收敛，Weierstrass 定理给出 \(f\) 全纯。若 \(f(a)=\infty\)，在小圆盘中像避开0，充分大的 \(1/f_n\) 全纯有界并欧氏一致收敛到 \(1/f\)。其极限全纯。若极限不恒零，零点孤立，对倒数取回得到 \(f\) 亚纯；若它在一个圆盘恒零，\(f\) 在该盘恒为无穷。全局若同时有有限值点，则连接两种区域的边界附近，倒数极限的恒等定理迫使继续为零，产生矛盾。连通性因此使恒无穷情形传播到整个区域。
\end{proof}
族 \(nz\) 在0附近没有球面一致极限，因为点0处始终为0，任意非零点却趋于无穷，候选极限不连续。允许无穷极限仍不能消除局部尖峰。

\ProseParagraphBreak

球面一致收敛允许极点的位置移动。例子 \(f_n(z)=1/(z-1/n)\) 在每个有限平面紧集上球面一致趋于 \(1/z\)：对有限非极点，
\[
 \chi\left(\frac1{z-1/n},\frac1z\right)
 =\chi(z-1/n,z)\le\frac1n,
\]
等式利用倒数是球面等距变换，并且在移动极点处也按连续性成立。欧氏一致收敛则不可能在含0的紧集上成立，因为函数有极点。球面坐标正是为了让这种自然极限连续可见。

\ProseParagraphBreak

全纯函数列的球面极限若不是恒为无穷，则实际上仍全纯，不会凭空产生一个孤立极点。若极限在 \(a\) 为极点，取小盘使其倒数在盘内全纯、仅中心有零点。充分晚的 \(f_n\) 在小盘内远离0，所以 \(1/f_n\) 全纯无零，局部一致趋于有中心零点但不恒零的 \(1/f\)，与 Hurwitz 无零极限结论矛盾。因此极点只能来自亚纯成员或整域无穷退化。这个细节在不同正规族定义之间换用时尤其重要。

% !TeX root = ../../Book3.tex
\section{Montel 定理}
\label{b3:sec:0403}

局部有界性提供导数上界，导数上界提供等度连续性，而有限覆盖与对角线抽取产生一个共同子列。Montel 定理的有限值版本因此把解析估计转换成紧性。省略值版本需要更深的几何控制，应与这一基本版本区别。

\subsection{局部有界族}
\label{b3:subsec:040301}

每个全纯函数在内部紧集上都连续有界，但这些上界可能随函数任意增大。抽取极限所需的是控制整个族的一个共同常数。

\Needspace{6\baselineskip}
\begin{definition}\label{b3:def:04-locally-bounded-family}
设 \(U\subseteq\mathbb C\) 为开集，\(\mathcal F\) 是从 \(U\) 到 \(\mathbb C\) 的函数族。若对每个紧集 \(K\subset U\)，存在 \(M_K<\infty\)，使
\[
 |f(z)|\le M_K\quad(f\in\mathcal F,\ z\in K),
\]
则称 \(\mathcal F\) 为 \(U\) 上的\term[jubu-youjie-zu]{局部有界族}[locally bounded family]。
\end{definition}

常数可以依赖观察的紧集，不能依赖族的成员。例如 \(\{z^n:n\ge1\}\) 在单位盘上共同有界为 \(1\)，而 \(\{n:n\ge1\}\) 在任意非空集合上都没有共同界。后者每个成员仍是全纯且有界的函数，所以“每个函数有界”与“整个族共同有界”是不同的条件。

\ProseParagraphBreak

定义等价于一个点附近的共同控制：对每个 \(a\in U\)，存在邻域 \(V_a\) 及常数 \(M_a\)，使全部 \(f\in\mathcal F\) 在整个 \(V_a\) 上满足 \(|f|\le M_a\)。从紧集定义推出这一点，只要在 \(a\) 周围选一个内部闭圆盘即可；反过来，用有限个 \(V_a\) 覆盖给定紧集，把对应有限个 \(M_a\) 取最大，就得到紧集版本。这里一个邻域同时控制所有函数及邻域中的所有点，不能简化为只在中心点有共同界。

\ProseParagraphBreak

回忆等度连续的量词：对紧集 \(K\) 上的族，给定 \(\varepsilon>0\)，要选出一个 \(\delta>0\)，使任意成员 \(f\) 和任意 \(z,w\in K\) 在 \(|z-w|<\delta\) 时都满足 \(|f(z)-f(w)|<\varepsilon\)。这个 \(\delta\) 同样不能依赖成员。全纯性通过 Cauchy 导数估计把共同上界转成这样的共同连续模。

\Needspace{6\baselineskip}
\begin{lemma}\label{b3:lem:04-equicontinuity}
设 \(U\subseteq\mathbb C\) 为开集，\(\mathcal F\) 是 \(U\) 上的全纯函数族。若对每个紧集 \(L\subset U\) 有
\[
 \sup_{f\in\mathcal F}\sup_{z\in L}|f(z)|<\infty,
\]
则 \(\mathcal F\) 在每个紧集 \(K\subset U\) 上等度连续。
\end{lemma}
\begin{proof}
空紧集无须讨论。对非空紧集 \(K\subset U\)，取 \(r>0\) 足够小，使
\[
 L=\{\zeta\in\mathbb C:\operatorname{dist}(\zeta,K)\le2r\}
 \subset U.
\]
这是可能的，因为 \(K\) 与闭集 \(\mathbb C\setminus U\) 有正距离；当 \(U=\mathbb C\) 时可任取 \(r>0\)。集合 \(L\) 闭且有界，因而紧。局部有界性给出一个 \(M_L\)，控制所有成员在 \(L\) 上的模。

对任意满足 \(\operatorname{dist}(\zeta,K)\le r\) 的点，闭圆盘 \(\overline{D(\zeta,r)}\) 包含于 \(L\)。\Cref{b3:thm:02-cauchy-estimate}给出
\[
 |f'(\zeta)|\le M_L/r\quad(f\in\mathcal F).
\]
若 \(z,w\in K\) 且 \(|z-w|<r\)，连接线段上的每点与 \(K\) 的距离都小于 \(r\)，所以沿该线段积分可得
\[
 |f(z)-f(w)|
 \le \frac{M_L}{r}\cdot|z-w|.
\]
给定 \(\varepsilon>0\)，取
\(\delta=\min\{r,\varepsilon r/(M_L+1)\}\)，便有统一的等度连续估计。
\end{proof}
证明中先把 \(K\) 向外加厚，再估计导数。这样即使 \(K\) 不凸，两个邻近点的连接线段也有一个共同的可用邻域；不能仅知道 \(K\) 上的导数上界，就默认任意连接线段留在 \(K\) 中。对更高阶导数，同一方法还给出
\[
 \sup_{f\in\mathcal F}\sup_{z\in K}|f^{(k)}(z)|
 \le k!r^{-k}M_L\quad(k\ge1).
\]
以下抽取只需一阶控制，极限得到以后，每个固定阶导数的收敛由 Weierstrass 定理自动提供。

\subsection{Montel 定理}
\label{b3:subsec:040302}

\term[montel-dingli]{Montel 定理}[Montel's theorem]的基本版本把刚才的估计转成一个共同收敛子列。证明不需要先假定函数空间中的紧性定理；我们直接在紧集上完成有限网格论证。
\Needspace{6\baselineskip}
\begin{theorem}[Montel，局部有界形式]\label{b3:thm:04-montel}
设 \(U\subseteq\mathbb C\) 为区域，\(\mathcal F\) 是 \(U\) 上的局部有界全纯函数族。则任取序列 \(f_n\in\mathcal F\)，都能抽取在 \(U\) 上局部一致收敛至全纯函数的子列，即 \(\mathcal F\) 在有限值意义下正规。
\end{theorem}
\begin{proof}
任取 \(f_n\in\mathcal F\)。先固定一个非空紧集 \(K\subset U\)，证明在 \(K\) 上可抽取一致收敛子列。对每个正整数 \(m\)，紧性给出 \(K\) 的有限 \(1/m\) 网：可以选有限个 \(K\) 中的点，使相应半径 \(1/m\) 的圆盘覆盖 \(K\)。把所有网点并在一起得到一个至多可数的稠密集 \(D=\{a_1,a_2,\ldots\}\)；若 \(D\) 有限，下面只需有限次抽取。

由于 \(|f_n(a_1)|\le M_K\)，复平面中的有界数列有收敛子列。再从这个子列中选出在 \(a_2\) 处也收敛的子列，依次进行，并按上一节的办法取对角线。把所得子列记为 \(g_\nu\)。对每个固定的网点 \(a_j\)，数列 \(g_\nu(a_j)\) 都收敛。

接着把稠密点上的收敛提升为整个 \(K\) 上的一致收敛。给定 \(\varepsilon>0\)，由\cref{b3:lem:04-equicontinuity}选 \(\delta>0\)，使任意成员在距离小于 \(\delta\) 的两点之间变化小于 \(\varepsilon/3\)。稠密性和紧性允许选有限个网点 \(a_{j_1},\ldots,a_{j_s}\)，使它们的 \(\delta\) 圆盘覆盖 \(K\)。这些中心上的数列均收敛，所以存在 \(N\)，使 \(\nu,\mu\ge N\) 时
\[
 |g_\nu(a_{j_l})-g_\mu(a_{j_l})|<\varepsilon/3
 \quad(1\le l\le s).
\]
对任意 \(z\in K\)，选一个距离小于 \(\delta\) 的中心 \(a_{j_l}\)，于是
\[
 \begin{aligned}
 |g_\nu(z)-g_\mu(z)|
 &\le |g_\nu(z)-g_\nu(a_{j_l})|\\
 &\quad+|g_\nu(a_{j_l})-g_\mu(a_{j_l})|
       +|g_\mu(a_{j_l})-g_\mu(z)|\\
 &<\varepsilon.
 \end{aligned}
\]
因此 \((g_\nu)\) 在一致范数中是 Cauchy 列。对每个 \(z\)，复数域的完备性给出极限 \(g(z)\)；在最后的不等式中令 \(\mu\to\infty\)，就得到整个 \(K\) 上的一致误差界。极限在 \(K\) 上连续，且有限。

现在取 \(U\) 的递增紧集穷尽 \(K_1,K_2,\ldots\)。对原序列在 \(K_1\) 上作刚才的抽取，再对所取子列在 \(K_2\) 上抽取，依次进行。第二次对角线得到在所有 \(K_j\) 上一致收敛的同一条子列。各极限在重叠处一致，拼成 \(U\) 上的有限函数 \(f\)；每个内部紧集最终包含于某个 \(K_j\)，故收敛局部一致。最后由\cref{b3:thm:04-weierstrass}，\(f\) 全纯。
\end{proof}
两次对角线承担不同的任务：第一次在一个紧集的稠密点上取得共同子列，再用等度连续性控制其余点；第二次把一个紧集扩展到整个开集。稠密点上的逐点收敛本身没有给出函数的一致收敛，有限网点的三角估计补上的正是这一步。

\ProseParagraphBreak

例如所有全纯映射 \(f:U\to\mathbb D\) 都满足共同界 \(1\)，因此是有限值正规族。这里不必事先要求单叶；若极值问题只取其中满足单叶性及某种归一化的子族，仍然正规，因为任意成员序列也是原族中的序列。Montel 只保证极限存在并全纯，极限是否仍在所选子族中，要由下一节的零点稳定性及约束来判断。

\ProseParagraphBreak

共同点值上界也可由平均大小推出。设全纯函数族满足：对每个闭圆盘 \(\overline{D(a,R)}\subset U\)，存在 \(C_{a,R}\)，使
\[
 \int_{D(a,R)}|f(z)|^2\,\mathrm{d}A(z)\le C_{a,R}
 \quad(f\in\mathcal F),
\]
其中 \(\mathrm dA\) 为平面面积。固定这样的圆盘，对 \(z\in D(a,R/2)\) 及 \(0<t<R/2\)，圆周 Cauchy 公式和 Cauchy--Schwarz 不等式给出
\[
 |f(z)|^2
 \le\frac1{2\pi}\int_0^{2\pi}|f(z+te^{\mathrm{i}\theta})|^2\,\mathrm d\theta.
\]
乘以 \(2\pi t\)，再从 \(0\) 积到 \(R/2\)，极坐标积分给出
\[
 \pi(R/2)^2|f(z)|^2
 \le\int_{D(z,R/2)}|f(\zeta)|^2\,\mathrm dA(\zeta)
 \le C_{a,R}.
\]
所以较小圆盘上的共同界为 \(2\sqrt{C_{a,R}}/(\sqrt\pi R)\)。每点附近都有这种界，有限覆盖便给出紧集共同界，Montel 定理于是适用。这个例子说明可积的平均控制也能通过解析性转化为逐点控制。

\ProseParagraphBreak

有限值正规性也反过来迫使局部有界。
\Needspace{6\baselineskip}
\begin{proposition}\label{b3:prop:04-normal-bounded-equivalence}
设 \(U\subseteq\mathbb C\) 为区域，\(\mathcal F\) 是 \(U\) 上的全纯函数族。则 \(\mathcal F\) 在有限值意义下正规，当且仅当它在 \(U\) 上局部有界。
\end{proposition}
\begin{proof}
局部有界推出正规是 Montel 定理。反过来，假设 \(\mathcal F\) 正规却不是局部有界，则存在紧集 \(K\subset U\)，使 \(\sup_{f\in\mathcal F}\sup_K|f|=\infty\)。对每个正整数 \(n\)，可选 \(f_n\in\mathcal F\) 和 \(z_n\in K\)，满足 \(|f_n(z_n)|>n\)。

正规性给出局部一致收敛子列 \(f_{n_k}\to f\)，其中 \(f\) 全纯。因 \(f\) 在 \(K\) 上连续有界，充分大的 \(k\) 满足
\[
 |f_{n_k}(z_{n_k})|
 \le \sup_K|f_{n_k}-f|+\max_K|f|
 \le 1+\max_K|f|.
\]
这与左端大于 \(n_k\to\infty\) 矛盾。此处不必再抽取点列：整个 \(K\) 上的一致误差已经直接控制了移动点 \(z_{n_k}\)。
\end{proof}
该等价是有限值版本的结论。常值族 \(\{n:n\ge1\}\) 在球面意义正规，却不局部有界，所以不能把“有限值”三个字默默删除。教材采用哪一种定义，决定了可使用的判据强弱。

\ProseParagraphBreak

这个等价也解释了局部性。只要每个点附近都能验证族的共同上界，便可在那些邻域抽取极限，再由\cref{b3:prop:04-normal-locality}拼成全域正规性。实际使用时通常先直接估计函数；只有共同上界不易得到而子列结构已知时，才反过来使用必要性。

\subsection{省略值形式}
\label{b3:subsec:040303}

所谓共同省略一个值，指每个成员在整个定义域都不取到该值。若省略的集合还含一个固定圆盘，就能直接化为已证明的有界情形。

\Needspace{6\baselineskip}
\begin{proposition}\label{b3:prop:04-omitted-disk}
设 \(U\subseteq\mathbb C\) 为区域，\(a\in\mathbb C\)、\(\varepsilon>0\)，\(\mathcal F\) 是 \(U\) 上的全纯函数族。若
\[
 |f(z)-a|\ge\varepsilon
 \quad(f\in\mathcal F,\ z\in U),
\]
则 \(\mathcal F\) 在球面意义下正规。
\end{proposition}
\begin{proof}
任取 \(f_n\in\mathcal F\)，令 \(q_n=1/(f_n-a)\)。这些函数全纯且共同有界为 \(1/\varepsilon\)，Montel 定理给出一个局部一致收敛子列，仍记为 \(q_n\to q\)，其中 \(q\) 全纯。

定义扩充平面上的映射
\[
 T(w)=a+\frac1w,\quad T(0)=\infty,\quad T(\infty)=a.
\]
它是球面的连续映射：倒数的连续性已由等距性质得到，平移在有限点连续，在无穷远也因 \(|z+a|\to\infty\) 而连续。紧球面上的连续映射一致连续。因此普通局部一致收敛 \(q_n\to q\) 先给球面局部一致收敛，再在复合 \(T\) 后保持这种收敛，得到
\[
 f_n=T\circ q_n\longrightarrow T\circ q
\]
球面局部一致。由\cref{b3:prop:04-spherical-limit}，极限是亚纯函数或恒为 \(\infty\)，故原族球面正规；由于成员全纯，\cref{b3:cor:04-holomorphic-spherical-limit}进一步排除了非恒常的带极点极限。
\end{proof}

若 \(q\equiv0\)，上述极限就是 \(\infty\)；若 \(q\) 不恒零，它的零点原本可能在倒数中形成极点，但全纯成员的球面极限结论排除了这种可能。因而非零的 \(q\) 在整个区域无零，取回得到有限全纯极限。这一分情况说明，倒数坐标既能保留有限极限，也能记录函数整体逃向无穷远。

\ProseParagraphBreak

更强的经典结论将在\cref{b3:thm:14-montel-omitted}建立：同一区域上的全纯函数族若共同省略两个固定的不同有限复值，则球面正规；亚纯函数族若共同省略三个固定的不同球面值，也球面正规。全纯函数本来就不取 \(\infty\)，因此这两种表述对应相同的三个球面省略值。

\ProseParagraphBreak

省略一个圆盘提供了 \(|f-a|\ge\varepsilon\) 的正距离，省略有限个点只给 \(|f-a|>0\)，其下界可以随成员和自变量一起趋于零，所以刚才的倒数有界证明不够。\secref{b3:sec:1403}将通过球面导数、重缩放及三穿孔球面的覆盖结构建立强版本；这里仅说明其位置，不把它当成本章或下一章的前置。Riemann 映照定理的极值构造使用的是已经完整证明的局部有界 Montel 定理。

\ProseParagraphBreak

省略值的数量不能随意减少。整函数族 \(e^{nz}\) 始终省略 \(0\)，却在含 \(0\) 的圆盘上不球面正规：原点值固定为 \(1\)，正实点沿任意指标趋于无穷的子列都趋于 \(\infty\)，与连续极限矛盾。其在半平面上的行为留作习题。球面正规性容许常数逃向无穷，却仍要求每个内部紧集上有同一个连续极限；共同省略一个有限值本身没有排除这里的失效方式。

% !TeX root = ../../Book3.tex
\section{Hurwitz 定理与全纯极限的应用}
\label{b3:sec:0404}

局部一致收敛已经保证极限全纯，却还没有回答极限能否取得某个值、是否仍为单射。第三章的零点计数提供了进一步的控制：在一个无零边界上，极限函数的模有正下界；函数列与它的误差最终小于这个下界，Rouché 定理便把内部零点的总重数传回函数列。本节先沿这条路线研究零点和映射，再把 Montel 紧性与恒等定理结合，说明何时一组点值极限足以决定整列的局部一致极限。

\subsection{零点稳定性}
\label{b3:subsec:040401}

零点的位置可以变化，几个单零点也可以在极限中合并。稳定的量是一个固定围道内部、按重数计算的零点总数，而保证稳定的条件是极限函数在围道上不为零。\term[hurwitz-dingli]{Hurwitz 定理}[Hurwitz's theorem]把这一观察用于全纯函数列。

\Needspace{6\baselineskip}
\begin{theorem}[Hurwitz]\label{b3:thm:04-hurwitz}
设 \(U\subseteq\mathbb C\) 为区域，\(f_n:U\rightarrow\mathbb C\) 全纯，并在 \(U\) 上局部一致收敛于有限值函数 \(f\)，且 \(f\) 不恒为零。若 \(a\in U\)、\(r>0\)、\(\overline{D(a,r)}\subseteq U\)，且 \(f\) 在 \(|z-a|=r\) 上无零点，则充分大的 \(n\) 使 \(f_n\) 在该圆周无零点，并与 \(f\) 在 \(D(a,r)\) 中具有相同的零点总重数。

\ProseParagraphBreak

此外，对每个满足 \(f(a)=0\) 的点 \(a\in U\)，存在整数 \(N\ge1\) 及点 \(a_n\in U\)、\(n\ge N\)，使 \(f_n(a_n)=0\) 且 \(a_n\rightarrow a\)。
\end{theorem}
\begin{proof}
由\cref{b3:thm:04-weierstrass}，\(f\) 全纯。它不恒零，所以零点孤立，在所给闭圆盘中只有有限个零点。圆周的紧性与无零条件给出
\[
 \delta=\min_{|z-a|=r}|f(z)|>0.
\]
局部一致收敛使充分大的 \(n\) 满足
\[
 \max_{|z-a|=r}|f_n(z)-f(z)|<\delta\le |f(z)|
 \quad(|z-a|=r).
\]
于是 \(f_n\) 在圆周也无零点。对全纯主项 \(f\) 和余项 \(f_n-f\) 应用\cref{b3:thm:03-rouche}，便得圆盘内的零点总重数相同。

再设 \(a\) 是 \(f\) 的 \(m\) 阶零点。取严格递减到零的半径 \(r_j>0\)，使各闭圆盘均包含于 \(U\)，且其中 \(f\) 的唯一零点为 \(a\)。由刚证明的计数结论，对每个 \(j\) 有整数 \(M_j\)，使全部 \(n\ge M_j\) 的 \(f_n\) 在 \(D(a,r_j)\) 中均有 \(m\) 个零点，计重数。将这些最终指标增大，递归取
\[
 N_1\ge M_1,\quad
 N_{j+1}\ge\max\{M_{j+1},N_j+1\}.
\]
对每个 \(n\ge N_1\)，令
\[
 j(n)=\max\{\,j\ge1\mid N_j\le n\,\}.
\]
由于 \(N_j\) 严格增加，这个最大值存在且有限。在 \(D(a,r_{j(n)})\) 中选一个 \(f_n\) 的零点 \(a_n\)。当 \(n\rightarrow\infty\) 时，\(j(n)\rightarrow\infty\)，故 \(|a_n-a|<r_{j(n)}\rightarrow0\)。这一选择对每个足够大的指标 \(n\) 都有定义，给出了所需的整条零点列。
\end{proof}

最后这一步包含两个不同的尺度。先固定半径，才能找到适用于该圆盘的最终指标；随后让半径按层缩小，并相应推迟指标，才能把“每个小盘最终有根”写成 \(a_n\rightarrow a\)。定理不要求给同一重根分裂出的各个根作唯一编号。比如 \(z^2-1/n\) 的两个单零点 \(\pm n^{-1/2}\) 都趋于零，而极限 \(z^2\) 在零点处的重数为二。

\ProseParagraphBreak

作为一个不能直接用多项式求根公式处理的例子，考虑
\[
 f_n(z)=z^2-1+\frac{e^z}{n}.
\]
在每个固定紧集上，末项一致趋于零，所以极限为 \(z^2-1\)。对 \(a=1\) 或 \(a=-1\)，在 \(|z-a|=1/4\) 上都有
\[
 |z^2-1|=|z-a|\cdot|z+a|\ge\frac14\cdot\frac74=\frac7{16},
 \quad \frac{|e^z|}{n}\le\frac{e^{5/4}}n.
\]
因此只要 \(n>16e^{5/4}/7\)，每个圆盘 \(D(1,1/4)\)、\(D(-1,1/4)\) 中都恰有一个计重数零点，因而都是单零点。继续缩小这两个圆盘，Hurwitz 定理说明相应零点分别趋于 \(1\) 和 \(-1\)。这里的边界下界既说明为何能数根，也给出了固定圆盘适用的具体指标。

\ProseParagraphBreak

反过来，若 \(f_n(a_n)=0\)，各点 \(a_n\) 始终落在某个紧集 \(K\subseteq U\) 中，且 \(a_n\rightarrow a\in K\)，则
\[
 |f(a)|\le |f(a)-f(a_n)|+\sup_K|f-f_n|\longrightarrow0,
\]
所以聚点也是极限函数的零点。共同紧集不能省略：\(f_n(z)=1-z/n\) 在全平面局部一致趋于 \(1\)，但它的零点 \(n\) 逃向无穷远。Hurwitz 的结论只约束固定内部区域中的零点。

\ProseParagraphBreak

若极限恒零，计数便失去上述依据。函数列 \(1/n\) 无零，函数列 \(z/n\) 每个成员在零点处有一个单零点，两列却都局部一致趋于零。恒零函数没有离散零点集，也无法在任何圆周上提供严格正的最低模；不能把它看成具有某个有限总重数的函数继续套用计数结论。

\subsection{无零极限}
\label{b3:subsec:040402}

没有零点是一种在极限下几乎封闭的性质。唯一需要留下的退化情形，是整个函数同时缩小到零。

\Needspace{6\baselineskip}
\begin{corollary}[无零极限]\label{b3:cor:04-zero-free-limit}
设 \(U\subseteq\mathbb C\) 为区域，\(f_n:U\rightarrow\mathbb C\) 全纯，且对全部 \(n\ge1\)、\(z\in U\) 有 \(f_n(z)\ne0\)。若 \(f_n\) 在 \(U\) 上局部一致收敛于有限值函数 \(f\)，则 \(f\) 全纯，并且或者在 \(U\) 恒为零，或者在 \(U\) 处处无零。
\end{corollary}
\begin{proof}
全纯性由 Weierstrass 定理给出。若 \(f\) 不恒零却在 \(a\in U\) 为零，取一个闭包包含于 \(U\) 的小圆盘，使 \(a\) 是其中唯一的零点。其重数为某个正整数，边界无零。Hurwitz 定理迫使全部充分晚的 \(f_n\) 在该圆盘中也有零点，与假设矛盾。
\end{proof}

常值函数列 \(1/n\) 说明恒零例外确会发生。若 \(U\) 上的全纯函数列 \(f_n\) 共同省略某个固定值 \(b\in\mathbb C\)，且局部一致趋于有限值 \(f\)，则对 \(f_n-b\) 应用推论：\(f\) 要么恒为 \(b\)，要么仍处处不取值 \(b\)。平移到零点以后，省略值问题也变成了零点稳定性。

\ProseParagraphBreak

当极限仍无零时，倒数也有通常的局部一致极限。具体地，设 \(U\) 上的无零全纯函数 \(f_n\) 局部一致趋于处处非零的 \(f\)。固定非空紧集 \(K\subseteq U\)，有 \(c_K=\min_K|f|>0\)，且充分大的 \(n\) 使 \(|f_n|\ge c_K/2\) 于 \(K\)。于是
\[
 \sup_K\left|\frac1{f_n}-\frac1f\right|
 \le\frac2{c_K^2}\cdot\sup_K|f_n-f|\longrightarrow0.
\]
分母的共同正下界是取倒数与极限交换的依据；若 \(f_n=1/n\)，倒数则趋于无穷，不能得到有限值版本的结论。

\ProseParagraphBreak

区域的连通性同样有作用。令 \(U=D(-2,1)\cup D(2,1)\)，在左边的圆盘取 \(f_n=1/n\)，右边取 \(f_n=1\)。每个成员在 \(U\) 全纯且无零，极限却在一个分支恒零、在另一个分支无零。对一般开集，无零极限推论应逐连通分支使用。若考虑全纯函数列的球面局部一致极限，\cref{b3:cor:04-holomorphic-spherical-limit}已说明它只能是全纯函数或全域恒为无穷，不会在内点新生孤立极点。

\subsection{单叶极限}
\label{b3:subsec:040403}

区域上的单射全纯函数称为\term[danye-hanshu]{单叶函数}[univalent function]。它的不同点不能映到同一值。要研究这一性质在极限中是否保留，可以固定一个点 \(a\)，把“另一个点也映到 \(f(a)\)”改写成差商的零点问题；在中心补上导数，便得到可使用无零极限推论的全纯函数。

\Needspace{6\baselineskip}
\begin{theorem}[单叶极限]\label{b3:thm:04-univalent-limit}
设 \(U\subseteq\mathbb C\) 为区域，\(f_n:U\rightarrow\mathbb C\) 是单射全纯函数，并在 \(U\) 上局部一致收敛于有限值函数 \(f\)。则 \(f\) 是常值函数，或者是单射全纯函数。
\end{theorem}
\begin{proof}
Weierstrass 定理保证 \(f\) 全纯。假设 \(f\) 非常值，任取 \(a\in U\)，在 \(z\ne a\) 处定义
\[
 g_n(z)=\frac{f_n(z)-f_n(a)}{z-a},
 \quad g(z)=\frac{f(z)-f(a)}{z-a},
\]
并令 \(g_n(a)=f_n'(a)\)、\(g(a)=f'(a)\)。Taylor 展开说明补定义后的函数在整个 \(U\) 全纯。对 \(z\ne a\)，\(f_n\) 的单射性给出 \(g_n(z)\ne0\)；在中心处，\cref{b3:cor:01-local-injectivity}说明单射全纯函数的导数不为零，故也有 \(g_n(a)\ne0\)。

须先核对 \(g_n\rightarrow g\) 是局部一致收敛。取 \(\rho>0\)，使 \(\overline{D(a,\rho)}\subseteq U\)。在此闭圆盘上，沿直线段积分给出
\[
 g_n(z)=\int_0^1 f_n'\bigl(a+t(z-a)\bigr)\,\mathrm{d}t,
\]
\(g\) 也有对应表达。所有线段都留在同一个闭圆盘中，故
\[
 \sup_{\overline{D(a,\rho)}}|g_n-g|
 \le\sup_{\overline{D(a,\rho)}}|f_n'-f'|\longrightarrow0,
\]
其中用了\cref{b3:prop:04-derivative-convergence}。对任意紧集 \(K\subseteq U\)，其在 \(D(a,\rho/2)\) 外的部分满足 \(|z-a|\ge\rho/2\)，所以
\[
 |g_n(z)-g(z)|
 \le\frac2\rho\bigl(|f_n(z)-f(z)|+|f_n(a)-f(a)|\bigr).
\]
这部分也一致收敛，与中心附近的估计合起来即得全 \(K\) 上的一致收敛。

由于 \(f\) 非常值，\(g\) 不恒为零。对无零函数列 \(g_n\) 使用\cref{b3:cor:04-zero-free-limit}，得 \(g\) 处处无零。因而对每个 \(z\ne a\) 都有 \(f(z)\ne f(a)\)。点 \(a\) 任意，所以 \(f\) 单射。
\end{proof}

函数列 \(f_n(z)=z/n\) 在任何固定区域上都单叶，却局部一致趋于常值零。若在某个固定点 \(a\in U\) 另有 \(|f_n'(a)|\ge c>0\)，则导数收敛给出 \(|f'(a)|\ge c\)，从而排除常值退化。这里只固定 \(f_n(a)\) 的值还不够，\(z/n\) 在原点已经具有相同的函数值。

\ProseParagraphBreak

映射的收敛还提出了一个定义域问题。即使 \(f_n(z)\rightarrow f(z)\)，给定 \(w\in f(U)\) 后，方程 \(f_n(z)=w\) 是否有解仍需证明；只有先确认这一点，才谈得上 \(f_n^{-1}(w)\) 的极限。单叶全纯函数导数无零，由局部逆函数定理，其逆函数在像区域上全纯。下面证明极限像的每个紧集最终都落在这些逆函数的定义域中。

\Needspace{6\baselineskip}
\begin{proposition}[单叶逆映射的收敛]\label{b3:prop:04-inverse-convergence}
设 \(U\subseteq\mathbb C\) 为区域，\(f_n:U\rightarrow\mathbb C\) 为单叶全纯函数，并在 \(U\) 上局部一致收敛于非常值函数 \(f\)。则 \(f\) 单叶全纯，且对每个紧集 \(K\subseteq f(U)\)，存在整数 \(N\)，使全部 \(n\ge N\) 满足 \(K\subseteq f_n(U)\)，并有
\[
 \sup_{w\in K}|f_n^{-1}(w)-f^{-1}(w)|\longrightarrow0.
\]
\end{proposition}
\begin{proof}
单叶性由前一定理保证。先固定 \(w_0\in f(U)\)，记 \(a=f^{-1}(w_0)\)。取 \(r>0\)，使闭圆盘 \(C=\overline{D(a,r)}\) 包含于 \(U\)。由于 \(f\) 单射，\(f-w_0\) 在该圆盘中只有一个单零点 \(a\)，边界不为零。因此
\[
 m=\min_{|z-a|=r}|f(z)-w_0|>0.
\]
先让目标点在 \(D(w_0,m/3)\) 中移动。对其中的每个 \(w\)，常数扰动 \(|w-w_0|<m/3\) 小于圆周上的 \(|f-w_0|\)。Rouché 定理说明 \(f-w\) 在源圆盘中也恰有一个计重数零点；同时圆周上有共同下界
\[
 |f(z)-w|\ge |f(z)-w_0|-|w-w_0|>2m/3.
\]
局部一致收敛提供一个与 \(w\) 无关的最终指标，使圆周上 \(|f_n-f|<m/3\)。第二次应用 Rouché 定理，比较 \(f-w\) 与 \(f_n-w\)，得 \(f_n-w\) 在该圆盘中同样恰有一个零点。故整个 \(D(w_0,m/3)\) 同时包含于 \(f(U)\) 和这些 \(f_n(U)\)，且相应逆像都落在 \(D(a,r)\) 中。

现在在较小闭盘 \(K_0=\overline{D(w_0,m/4)}\) 上比较逆像。对 \(w\in K_0\)，记 \(z_n=f_n^{-1}(w)\)、\(z=f^{-1}(w)\)。它们都在紧集 \(C\) 内，且
\[
 |f(z_n)-f(z)|=|f(z_n)-f_n(z_n)|
 \le\sup_C|f-f_n|\longrightarrow0.
\]
连续单射 \(f:C\rightarrow f(C)\) 是从紧空间到 Hausdorff 空间的同胚，故其逆映射在紧集 \(f(C)\) 上一致连续。这里 \(f(z_n)\) 和 \(f(z)\) 均属于 \(f(C)\)，所以给定 \(\varepsilon>0\)，可选一个对全部 \(w\in K_0\) 有效的 \(\delta>0\)，使上式左端小于 \(\delta\) 时便有 \(|z_n-z|<\varepsilon\)。因而逆映射在 \(K_0\) 上一致收敛。

对给定紧集 \(K\subseteq f(U)\) 的每个点执行这一构造，所得较小像盘的内部覆盖 \(K\)。取有限子覆盖，并将各处的包含指标及一致误差指标取最大，便同时得到 \(K\subseteq f_n(U)\) 和整个 \(K\) 上的逆映射一致收敛。
\end{proof}

逆映射的定义域 \(f_n(U)\) 随指标变化，所以这里的收敛按紧集理解：先固定极限像中的一个紧集，再从某项起把所有逆映射限制到这个共同集合上。Rouché 比较中的边界正距离同时控制了目标点 \(w\) 的移动和近似函数 \(f_n\) 的误差，才使最终指标对整块像盘统一有效。

\ProseParagraphBreak

例如在单位圆盘 \(\mathbb D\) 上取
\[
 f_n(z)=\frac{z}{1+z/n}.
\]
分母在盘内无零，交叉相乘可核对 \(f_n\) 单射；在任意 \(|z|\le r<1\) 上，\(|f_n(z)-z|\le r^2/(n-r)\)，所以局部一致趋于恒等映射。解方程得到逆函数候选式
\[
 f_n^{-1}(w)=\frac{w}{1-w/n}.
\]
对 \(|w|\le r<1\)，当 \(n>r/(1-r)\) 时，这个式子的模至多为 \(r/(1-r/n)<1\)，确实给出位于原定义域中的逆像，而且
\[
 |f_n^{-1}(w)-w|\le\frac{r^2}{n(1-r/n)}\longrightarrow0.
\]
但对每个固定 \(n\)，选实数 \(n/(n+1)<w<1\)，同一逆函数候选值却大于 \(1\)，所以该 \(w\) 不属于 \(f_n(\mathbb D)\)。因此每个内部紧集最终被覆盖，并不意味着整个极限像 \(\mathbb D\) 最终都包含于 \(f_n(\mathbb D)\)。

\subsection{Vitali 型收敛判据}
\label{b3:subsec:040404}

Montel 定理提供收敛子列，还不能单独确定它们是否具有同一个极限。恒等定理恰好提供所需的唯一性：两个全纯极限只要在具有内部聚点的集合上相等，就在整个区域相等。把这两项结合，得到下面的\term[vitali-shoulian-dingli]{Vitali 收敛定理}[Vitali convergence theorem]的局部有界形式。

\Needspace{6\baselineskip}
\begin{proposition}[Vitali，局部有界形式]\label{b3:prop:04-vitali-bounded}
设 \(U\subseteq\mathbb C\) 为区域，\(f_n:U\rightarrow\mathbb C\) 全纯，并且对每个紧集 \(K\subseteq U\)，存在有限常数 \(M_K\)，使全部 \(n\ge1\) 满足 \(\sup_K|f_n|\le M_K\)。设 \(E\subseteq U\) 在 \(U\) 内有聚点，且对每个 \(z\in E\)，数列 \(f_n(z)\) 收敛于有限值。则存在全纯函数 \(f:U\rightarrow\mathbb C\)，使整列 \(f_n\) 在 \(U\) 上局部一致收敛于 \(f\)。
\end{proposition}
\begin{proof}
局部有界性使\cref{b3:thm:04-montel}适用，先从原序列抽取一条局部一致收敛子列，记其全纯极限为 \(f\)。这一步给出了候选极限的存在。

再任取原序列的一条子列。它仍属于同一个局部有界族，故 Montel 定理使它有局部一致收敛的子子列，记全纯极限为 \(g\)。对每个 \(z\in E\)，两条已收敛的子列都必须趋于原数列 \(f_n(z)\) 的点值极限，所以 \(f(z)=g(z)\)。集合 \(E\) 有位于 \(U\) 内部的聚点，区域 \(U\) 又连通，\cref{b3:thm:01-identity}因此给出 \(g=f\) 于整个 \(U\)。于是原序列的每条子列都能再抽出以 \(f\) 为局部一致极限的子子列。

最后证明整列收敛。若在某个紧集 \(K\subseteq U\) 上不一致趋于 \(f\)，则存在 \(\varepsilon>0\) 及严格增加的指标 \(n_j\)，使
\[
 \sup_K|f_{n_j}-f|\ge\varepsilon\quad(j\ge1).
\]
由前一步，这条子列又有局部一致趋于 \(f\) 的子子列，因而在 \(K\) 上的误差趋于零，与固定正下界矛盾。
\end{proof}

这一判据尤其适合从实变量极限获得复变量结论。令
\[
 F_n(z)=\left(1+\frac zn\right)^n\quad(z\in\mathbb C).
\]
每个 \(F_n\) 都是整函数，且
\[
 |F_n(z)|\le\left(1+\frac{|z|}{n}\right)^n\le e^{|z|}.
\]
任意固定紧集上的 \(|z|\) 有上界，故这给出了全族的局部共同界。在实区间 \((-1,1)\) 上，实变量的指数极限给出 \(F_n(x)\rightarrow e^x\)。Vitali 判据使整列在全平面局部一致趋于某个整函数 \(F\)；\(F\) 与 \(e^z\) 在这个实区间相等，恒等定理于是确定 \(F=e^z\)。

\ProseParagraphBreak

得到局部一致收敛后，才可以应用 Weierstrass 定理取导数。对每个固定整数 \(k\ge1\)，在 \(n\ge k\) 时有
\[
 F_n^{(k)}(z)=\frac{n(n-1)\cdots(n-k+1)}{n^k}
 \cdot\left(1+\frac zn\right)^{n-k}\longrightarrow e^z,
\]
收敛在每个紧集上一致成立。这条推导只在实区间使用了最初的点值极限，区域中的共同上界和全纯唯一性负责把它扩展到整个复平面。

\ProseParagraphBreak

局部有界条件不能删除。在 \(\mathbb D\) 上取 \(f_n(z)=(2z)^n\)，则对实区间 \(E=(-1/2,1/2)\) 中每个点都有 \(f_n(x)\rightarrow0\)，而 \(E\) 具有内部聚点；可是 \(f_n(3/4)=(3/2)^n\) 发散。该序列在包含 \(3/4\) 的紧集上没有共同上界，不能从点值数据抽出覆盖整个区域的收敛子列。

\ProseParagraphBreak

内部聚点的要求也不能改成仅在边界聚集，即使全族在区域上已有统一界。下面在单位圆盘内构造一个有界全纯函数，使它在无内部聚点的无限集合上为零。令
\[
 W(z)=\frac{1+z}{1-z},\quad
 q(z)=e^{-W(z)}-e^{-2}\quad(z\in\mathbb D).
\]
分母在 \(\mathbb D\) 无零，且
\[
 \operatorname{Re}W(z)=\frac{1-|z|^2}{|1-z|^2}>0,
\]
所以 \(q\) 全纯，并有 \(|q(z)|\le1+e^{-2}\)。对每个整数 \(k\ge1\)，取
\[
 z_k=\frac{1+2\pi\mathrm{i}k}{3+2\pi\mathrm{i}k}.
\]
其模平方为 \((1+4\pi^2k^2)/(9+4\pi^2k^2)<1\)，故 \(z_k\in\mathbb D\)；又有 \(z_k\rightarrow1\)，且各点互不相同。直接代入得 \(W(z_k)=2+2\pi\mathrm{i}k\)，从而 \(q(z_k)=0\)。于是 \(E=\{\,z_k\mid k\ge1\,\}\) 完全包含于定义域，却只在边界点 \(1\) 聚集。

\ProseParagraphBreak

现在令 \(f_{2n}=0\)、\(f_{2n-1}=q\)。这是一列在 \(\mathbb D\) 上共同有界的全纯函数；在每个 \(z\in E\) 上，整个数列 \(f_n(z)\) 恒为零，当然收敛。但 \(q(0)=e^{-1}-e^{-2}\ne0\)，所以 \(f_n(0)\) 在两个不同的值之间交替，整列连逐点收敛都不成立。本例中的两条子列极限 \(0\) 与 \(q\) 确实存在，也在 \(E\) 上相等；失效的是恒等定理所需的内部聚点，因而不能把子列的存在提升为整列收敛。


\begin{chaptersummary}
局部一致拓扑由内部紧集控制，不限制边界附近的全域上确界。Cauchy 公式使有限极限保留全纯性，并控制每个固定阶导数；Montel 定理将局部上界转换为可抽取的函数列。Hurwitz 定理保留零点总重数，继而得到无零性和单叶性，常数退化是两者都必须留下的例外。球面正规性允许无穷极限，但仍要求极限连续，不能容许内部尖峰。
\end{chaptersummary}

\BookExercises

% 本章题组自拟；以正文已建立的工具设计计算、反例、构造与综合训练。
\begin{exercise}\label{b3:ex:04-1}
证明 \(nz^n\) 在 \(\mathbb D\) 上局部一致趋于0，却不在 \(\mathbb D\) 上一致收敛。
\end{exercise}
\begin{solution}
在 \(|z|\le r<1\) 上误差不超过 \(nr^n\)，它趋于0；取足够靠近1的实数点可使 \(n|z|^n\) 任意接近 \(n\)，故全域上确界为 \(n\)。
\end{solution}

\begin{exercise}\label{b3:ex:04-2}
若 \(f_n\rightarrow f\) 在开集上局部一致且每个 \(f_n\) 全纯，证明 \(\int_\gamma f_n\,\mathrm{d}z\rightarrow\int_\gamma f\,\mathrm{d}z\)，其中 \(\gamma\) 为固定分段光滑曲线。
\end{exercise}
\begin{solution}
曲线像是内部紧集，长度估计给出两积分之差的模不超过 \(L(\gamma)\sup_\gamma|f_n-f|\)，其趋于0。曲线固定使长度不随指标增长。
\end{solution}

\begin{exercise}\label{b3:ex:04-3}
设 \(f_n\) 在圆盘全纯且 \(|f_n|\le1\)。证明存在子列，使全部阶数的原点 Taylor 系数逐个收敛，并说明为何系数极限确定同一个全纯函数。
\end{exercise}
\begin{solution}
Montel 先抽取函数局部一致收敛子列。任意固定阶导数随之收敛，所以全部阶数的系数都沿这一子列逐个收敛。Cauchy 估计给每个系数 \(|a_k|\le1\)，极限级数在盘内收敛；在较小盘上它等于函数极限的 Taylor 展开。无须对每阶再选不同子列。
\end{solution}

\begin{exercise}\label{b3:ex:04-4}
构造无零整函数列，局部一致极限恰为零；再构造在 \(\mathbb D\) 无零的函数列，其极限仍无零却趋近边界时模可任意小。
\end{exercise}
\begin{solution}
第一列取 \(1/n\)。第二列可取常列 \(1-z\)，它在开盘无零，沿 \(z=r\uparrow1\) 模趋于0。Hurwitz 的无零结论不提供统一的全域正下界。
\end{solution}

\begin{exercise}\label{b3:ex:04-5}
证明 \(\{e^{nz}:n\ge1\}\) 在右半平面球面正规，在任何包含0的圆盘不球面正规。
\end{exercise}
\begin{solution}
右半平面任意紧集上实部有正下界 \(\delta\)，故 \(|e^{nz}|\ge e^{n\delta}\) 一致趋于无穷；任意序列若指标无界可抽此类子列，指标有界则抽常子列。包含0的圆盘内，沿正负实点的候选极限分别为无穷与0，而原点固定为1，任何无界指标子列的候选极限都不连续，故不正规。
\end{solution}

\begin{exercise}\label{b3:ex:04-6}
设 \(f_n:U\rightarrow\mathbb D\) 单叶全纯，固定 \(a\in U\) 且 \(|f_n'(a)|\ge c>0\)。证明存在单叶全纯局部一致子列极限。
\end{exercise}
\begin{solution}
共同有界使 Montel 适用。导数收敛给 \(|f'(a)|\ge c\)，排除常数极限；单叶极限定理因此给出单叶性。最大模原理又使有限极限仍取值于开盘，因为若内部模达到1它为常数。
\end{solution}

\begin{exercise}\label{b3:ex:04-7}
设 \(f_n\) 全纯且局部一致趋于 \(f\)，\(f\) 在 \(a\) 有二阶零点。说明为何不能保证 \(f_n\) 在附近也有二阶零点，并给例子。
\end{exercise}
\begin{solution}
取 \(f_n(z)=z^2-1/n\)。局部一致极限为 \(z^2\)，但每个成员在 \(\pm1/\sqrt n\) 有两个单零点。Hurwitz 保留的是充分小围道内部总重数2，不是某个指定点的阶数。
\end{solution}

\begin{exercise}\label{b3:ex:04-8}
设整函数列 \(f_n\) 在每个闭圆盘共同有界，且 \(f_n(1/j)\) 对每个正整数 \(j\) 都有极限。证明整列局部一致收敛，并说明这为何不要求知道聚点0处的点值极限。
\end{exercise}
\begin{solution}
集合 \(E=\{1/j:j\ge1\}\) 在内部有聚点0。前述局部有界收敛命题直接给整列局部一致收敛。两个候选全纯极限在 \(E\) 上相同，恒等定理只需要这些零点在0聚集，不另要求预先规定0的极限。
\end{solution}

